% ECONOMIC GEOLOGY — Geology Upper Sixth — Cours signé OnBuch+
% Official MINESEC syllabus. Compiler avec : tectonic GEO06-economic-geology.tex
\newcommand{\DOCMATIERE}{Geology}
\newcommand{\DOCNIVEAU}{Upper Sixth}
\newcommand{\DOCMODULE}{Upper Sixth — Geology}
\newcommand{\DOCLECON}{6}
\newcommand{\DOCTITRE}{ECONOMIC GEOLOGY}
% OnBuch+ — préambule « cours signés » ENRICHI (compilé avec tectonic / XeLaTeX)
% Système visuel « Pop » d'OnBuch+ : fond crème (jamais blanc), bordures encre
% épaisses, ombres « dures » décalées, titres Archivo Black en capitales.
% Version MATHÉMATIQUES : graphes (pgfplots/tikz), boîtes pédagogiques
% (définition, propriété, méthode, attention, le savais-tu, exercice résolu,
% exercice, TP, bilan), page de garde et signature OnBuch+ ancrée en bas.
%
% Macros attendues dans main.tex AVANT \input{preamble} :
%   \DOCMATIERE  \DOCNIVEAU  \DOCMODULE  \DOCLECON (numéro)  \DOCTITRE
\documentclass[11pt]{article}

\usepackage{fontspec}
\usepackage[english]{babel}
\usepackage[a4paper,margin=2cm,top=3.4cm,bottom=2.8cm,headheight=1.4cm,footskip=1.3cm]{geometry}
\usepackage{amsmath,amssymb,amsfonts}
\usepackage{mathtools} % \xmapsto, \coloneqq... souvent utilisés par les modèles
\usepackage{cancel} % simplifications barrées (bilans de forces)
% \abs{z} (module, valeur absolue) et \norm{u} (norme), utilisés par les modèles
\providecommand{\abs}{}\let\abs\relax\DeclarePairedDelimiter{\abs}{\lvert}{\rvert}
\providecommand{\norm}{}\let\norm\relax\DeclarePairedDelimiter{\norm}{\lVert}{\rVert}
\usepackage[table]{xcolor} % [table] : fournit aussi \rowcolors (alternance de lignes)
\usepackage{pagecolor}
\usepackage{tikz}
\usetikzlibrary{shapes.symbols,circuits.logic.US,circuits.logic.IEC,arrows.meta,positioning,calc,shapes.geometric,shapes.misc,shapes.arrows,decorations.pathmorphing,decorations.pathreplacing,decorations.markings,patterns,shadows,mindmap,trees,fit,backgrounds,matrix,intersections,angles,quotes}
\usepackage{pgfplots}
\usepackage[version=4]{mhchem} % équations nucléaires \ce{^{235}_{92}U}
\usepackage[european,siunitx]{circuitikz} % schémas électriques
\pgfplotsset{compat=1.18}
\usepgfplotslibrary{fillbetween,groupplots} % aires entre courbes (calcul intégral)
\usepackage{enumitem}
\usepackage{fancyhdr}
% [most] charge toutes les bibliothèques tcolorbox (listings, minted, raster,
% documentation...) et gonfle inutilement la mémoire TeX sur un document long
% (~300 boîtes) : on ne charge que ce qui est réellement utilisé.
\usepackage[skins,breakable]{tcolorbox}
\usepackage{graphicx}
\usepackage{colortbl}
\usepackage{booktabs}
\usepackage{tabularx}
\usepackage{diagbox} % case d'en-tête diagonale des tableaux à double entrée
% Colonnes extensibles centrée / gauche / droite (C, L, R), souvent utilisées par les modèles
\newcolumntype{C}{>{\centering\arraybackslash}X}
\newcolumntype{L}{>{\raggedright\arraybackslash}X}
\newcolumntype{R}{>{\raggedleft\arraybackslash}X}
\usepackage{array}
\usepackage{multicol}
\usepackage{siunitx}
% parse-numbers=false : accepte aussi \SI{2,5 \times 10^{-3}}{...}
\sisetup{locale=UK,output-decimal-marker={.},group-minimum-digits=4,parse-numbers=false}
\DeclareSIUnit\ohm{\ensuremath{\symup{\Omega}}} % U+2126 absent de la police
\DeclareSIUnit\division{div} % réglages d'oscilloscope (ms/div, V/div)
\DeclareSIUnit\year{yr}\DeclareSIUnit\annum{yr}\DeclareSIUnit\Ma{Ma}\DeclareSIUnit\tr{tr} % tours (vitesse de rotation, ex. \SI{1500}{\tr\per\minute})
\usepackage{qrcode}
% \degree : commande fréquemment utilisée par les modèles pour un angle ou une
% température en dehors de \SI{}{} (non fournie par les paquets chargés).
\providecommand{\degree}{^{\circ}}
% \qed (fin de démonstration), utilisé par les modèles sans amsthm
\providecommand{\qed}{\hfill$\square$}
% \barre : double barre des valeurs interdites dans un tableau de variations (tkz-tab)
\providecommand{\barre}{\ensuremath{\|}}
% environnement proof (amsthm non chargé) utilisé par les modèles
\ifdefined\proof\else\newenvironment{proof}[1][Proof]{\par\noindent\textit{#1.}\ }{\qed\par}\fi
\usepackage{lastpage}
\usepackage{hyperref}
\hypersetup{hidelinks}

% ============================================================
% Palette « Pop » OnBuch+ — mêmes hex que la classe P (plus_ui.dart)
% ============================================================
\definecolor{poporange}{HTML}{F59321}
\definecolor{popdark}{HTML}{E07A0C}
\definecolor{popcream}{HTML}{FFF6E8}
\definecolor{popink}{HTML}{1C1714}
\definecolor{poppurple}{HTML}{7A5AE0}
\definecolor{popgreen}{HTML}{2E9E6A}
\definecolor{poppink}{HTML}{E0457B}
\definecolor{popblue}{HTML}{2F7DE1}
\definecolor{popgold}{HTML}{C98A12}
\definecolor{popmuted}{HTML}{8A7F76}
\definecolor{popline}{HTML}{EADFD2}
\definecolor{popgray}{HTML}{8A7F76}
\colorlet{popgrey}{popgray}
% Alias tolérants (le modèle nomme parfois les couleurs différemment)
\colorlet{popwhite}{white}
\colorlet{popblack}{popink}
\colorlet{popred}{poppink}
\colorlet{popyellow}{popgold}
% Teintes claires pour fonds de tableaux / zones
\colorlet{popblueL}{popblue!10!white}
\colorlet{poporangeL}{poporange!14!white}
\colorlet{popgreenL}{popgreen!12!white}
\colorlet{poppurpleL}{poppurple!12!white}
\colorlet{poppinkL}{poppink!12!white}
\colorlet{popgoldL}{popgold!14!white}

\pagecolor{popcream}

% ============================================================
% Typographie
% ============================================================
\defaultfontfeatures{Ligatures=TeX}
\setmainfont{PlusJakartaSans-Medium.ttf}[Path=../../../fonts/,BoldFont=PlusJakartaSans-Bold.ttf]
% Maths en sans-serif (Fira Math) avec chiffres et lettres droites de Plus
% Jakarta, pour que les formules se fondent dans le texte.
\usepackage[math-style=ISO,bold-style=ISO]{unicode-math}
\setmathfont{FiraMath-Regular.otf}[Path=../../../fonts/]
\setmathfont{PlusJakartaSans-Medium.ttf}[Path=../../../fonts/,range={up/num,up/{Latin,latin},\mathup}]
% (icomma retiré : décimales avec point en anglais)
% Caractères mathématiques Unicode absents des polices de texte (Plus Jakarta,
% Archivo Black) : rendus par leur équivalent mathématique, en texte comme en
% formule — sinon ils s'affichent en carré vide (ex. « Arithmétique dans ℤ »).
\usepackage{newunicodechar}
\newunicodechar{ℝ}{\ensuremath{\mathbb{R}}}
\newunicodechar{ℤ}{\ensuremath{\mathbb{Z}}}
\newunicodechar{ℂ}{\ensuremath{\mathbb{C}}}
\newunicodechar{ℕ}{\ensuremath{\mathbb{N}}}
\newunicodechar{ℚ}{\ensuremath{\mathbb{Q}}}
\newunicodechar{𝔻}{\ensuremath{\mathbb{D}}}
\newunicodechar{α}{\ensuremath{\alpha}}
\newunicodechar{β}{\ensuremath{\beta}}
\newunicodechar{γ}{\ensuremath{\gamma}}
\newunicodechar{δ}{\ensuremath{\delta}}
\newunicodechar{ε}{\ensuremath{\varepsilon}}
\newunicodechar{θ}{\ensuremath{\theta}}
\newunicodechar{λ}{\ensuremath{\lambda}}
\newunicodechar{μ}{\ensuremath{\mu}}
\newunicodechar{σ}{\ensuremath{\sigma}}
\newunicodechar{τ}{\ensuremath{\tau}}
\newunicodechar{φ}{\ensuremath{\varphi}}
\newunicodechar{ψ}{\ensuremath{\psi}}
\newunicodechar{ω}{\ensuremath{\omega}}
\newunicodechar{Σ}{\ensuremath{\Sigma}}
% majuscules produites par \MakeUppercase dans les titres (α-aminés -> Α)
\newunicodechar{Α}{\ensuremath{\alpha}}
\newunicodechar{Β}{\ensuremath{\beta}}
\newunicodechar{Γ}{\ensuremath{\Gamma}}
\newunicodechar{Δ}{\ensuremath{\Delta}}
\newunicodechar{Ε}{\ensuremath{\varepsilon}}
\newunicodechar{Θ}{\ensuremath{\Theta}}
\newunicodechar{Λ}{\ensuremath{\Lambda}}
\newunicodechar{Μ}{\ensuremath{\mu}}
\newunicodechar{Τ}{\ensuremath{\tau}}
\newunicodechar{Φ}{\ensuremath{\Phi}}
\newunicodechar{Ψ}{\ensuremath{\Psi}}
\newunicodechar{Ω}{\ensuremath{\Omega}}
\newunicodechar{↦}{\ensuremath{\mapsto}}
\newunicodechar{∈}{\ensuremath{\in}}
\newunicodechar{∉}{\ensuremath{\notin}}
\newunicodechar{∧}{\ensuremath{\wedge}}
\newunicodechar{⇔}{\ensuremath{\Leftrightarrow}}
\newunicodechar{⟺}{\ensuremath{\Longleftrightarrow}}
\newunicodechar{⇒}{\ensuremath{\Rightarrow}}
\newunicodechar{⟹}{\ensuremath{\Longrightarrow}}
\newunicodechar{∘}{\ensuremath{\circ}}
\newunicodechar{∪}{\ensuremath{\cup}}
\newunicodechar{∩}{\ensuremath{\cap}}
\newunicodechar{≡}{\ensuremath{\equiv}}
\newunicodechar{⊂}{\ensuremath{\subset}}
\newunicodechar{⊥}{\ensuremath{\perp}}
\newunicodechar{∃}{\ensuremath{\exists}}
\newunicodechar{∀}{\ensuremath{\forall}}
\newunicodechar{∛}{\ensuremath{\sqrt[3]{\,}}}
\newunicodechar{∜}{\ensuremath{\sqrt[4]{\,}}}
\newunicodechar{⃗}{\ensuremath{{}^{\to}}}
\newunicodechar{ⁿ}{\ifmmode^{n}\else\textsuperscript{\ensuremath{n}}\fi}
\newunicodechar{ˣ}{\ifmmode^{x}\else\textsuperscript{\ensuremath{x}}\fi}
\newunicodechar{ᵇ}{\ifmmode^{b}\else\textsuperscript{\ensuremath{b}}\fi}
\newunicodechar{ʳ}{\ifmmode^{r}\else\textsuperscript{\ensuremath{r}}\fi}
\newunicodechar{ᵘ}{\ifmmode^{u}\else\textsuperscript{\ensuremath{u}}\fi}
\newunicodechar{ʸ}{\ifmmode^{y}\else\textsuperscript{\ensuremath{y}}\fi}
\newunicodechar{ᵃ}{\ifmmode^{a}\else\textsuperscript{\ensuremath{a}}\fi}
\newunicodechar{ᵉ}{\ifmmode^{e}\else\textsuperscript{\ensuremath{e}}\fi}
\newunicodechar{ᵏ}{\ifmmode^{k}\else\textsuperscript{\ensuremath{k}}\fi}
\newunicodechar{ᵖ}{\ifmmode^{p}\else\textsuperscript{\ensuremath{p}}\fi}
\newunicodechar{ᴺ}{\ifmmode^{N}\else\textsuperscript{\ensuremath{N}}\fi}
\newunicodechar{ᶜ}{\ifmmode^{c}\else\textsuperscript{\ensuremath{c}}\fi}
\newunicodechar{ʰ}{\ifmmode^{h}\else\textsuperscript{\ensuremath{h}}\fi}
\newunicodechar{ᵠ}{\ifmmode^{q}\else\textsuperscript{\ensuremath{q}}\fi}
\newunicodechar{ₙ}{\ifmmode_{n}\else\textsubscript{\ensuremath{n}}\fi}
\newunicodechar{ₐ}{\ifmmode_{a}\else\textsubscript{\ensuremath{a}}\fi}
\newunicodechar{ᵢ}{\ifmmode_{i}\else\textsubscript{\ensuremath{i}}\fi}
\newunicodechar{ᵣ}{\ifmmode_{r}\else\textsubscript{\ensuremath{r}}\fi}
\newunicodechar{ₖ}{\ifmmode_{k}\else\textsubscript{\ensuremath{k}}\fi}
\newunicodechar{ᵦ}{\ifmmode_{\beta}\else\textsubscript{\ensuremath{\beta}}\fi}
\newunicodechar{ₓ}{\ifmmode_{x}\else\textsubscript{\ensuremath{x}}\fi}
\newfontfamily\popdisplay{ArchivoBlack-Regular.ttf}[Path=../../../fonts/]
\newfontfamily\poplogo{SpaceGrotesk-Bold.ttf}[Path=../../../fonts/]
\newfontfamily\popsemi{PlusJakartaSans-SemiBold.ttf}[Path=../../../fonts/]
\newfontfamily\popxbold{PlusJakartaSans-ExtraBold.ttf}[Path=../../../fonts/]
\newfontfamily\popxboldls{PlusJakartaSans-ExtraBold.ttf}[Path=../../../fonts/,LetterSpace=3.0]

\setlist[enumerate]{leftmargin=1.7em,itemsep=5pt,topsep=4pt}
\setlist[itemize]{leftmargin=1.7em,itemsep=4pt,topsep=4pt,label={\color{poporange}\textbullet}}
\setlength{\parindent}{0pt}
\setlength{\parskip}{5pt}
\renewcommand{\arraystretch}{1.25}

% pgfplots : style Pop par défaut
\pgfplotsset{
  every axis/.append style={axis line style={popink,line width=1pt},tick style={popink},
    grid style={popline,line width=0.5pt},label style={font=\small},tick label style={font=\footnotesize}},
  popcurve/.style={poporange,line width=1.8pt,no markers,smooth},
}
% Même style disponible dans un \draw TikZ simple (hors pgfplots)
\tikzset{popcurve/.style={poporange,line width=1.8pt}}
\tikzset{popgrid/.style={popline,very thin}}
\colorlet{popcurve}{poporange} % le nom du style est parfois utilisé comme couleur

% ============================================================
% Pill / squiggle
% ============================================================
\newcommand{\poppill}[3]{%
  \tikz[baseline=-0.55ex]\node[draw=popink,line width=1.2pt,rounded corners=2.6pt,%
    fill=#2,inner xsep=6.5pt,inner ysep=3pt]{%
    \popxboldls\fontsize{7}{7}\selectfont\color{#3}\MakeUppercase{#1}};%
}
\newcommand{\popsquiggle}[1]{%
  \tikz[baseline=-0.4ex]{
    \draw[#1,line width=2.6pt,line cap=round]
      plot [smooth] coordinates {(0,0.09) (0.34,0.01) (0.68,0.21) (1.02,0.01) (1.36,0.10)};
  }%
}
% Mot-clé surligné (marqueur orange sous le texte)
\newcommand{\cle}[1]{{\popxbold\color{popdark}#1}}

% ============================================================
% En-tête / pied
% ============================================================
\pagestyle{fancy}
\fancyhf{}
\renewcommand{\headrulewidth}{0pt}
\renewcommand{\footrulewidth}{0pt}
\lhead{\raisebox{-0.15ex}{\poplogo\fontsize{13}{13}\selectfont\color{popink}OnBuch\color{poporange}+}}
\rhead{\poppill{\DOCMATIERE\ \textbullet\ \DOCNIVEAU\ \textbullet\ Chapter \DOCLECON}{white}{popink}}
\lfoot{\popxbold\fontsize{7.6}{7.6}\selectfont\color{popmuted}\MakeUppercase{Signed course OnBuch+}\ \textbullet\ {\popsemi\DOCTITRE}}
\rfoot{\tikz[baseline=-0.6ex]\node[rounded corners=8.5pt,fill=popink,minimum height=17pt,minimum width=17pt,inner xsep=5pt,inner sep=0]{\popxbold\fontsize{8}{8}\selectfont\color{popcream}\thepage\,/\,\pageref*{LastPage}};}

% ============================================================
% Titres
% ============================================================
\newcommand{\chaptitre}[2]{%
  \begin{tcolorbox}[enhanced,colback=poporange,colframe=popink,boxrule=2.6pt,arc=11pt,
      drop shadow=popink,left=15pt,right=15pt,top=12pt,bottom=11pt,before skip=3pt,after skip=18pt]
    \poppill{#2}{popink}{popcream}\par\vspace{9pt}
    {\raggedright\hyphenpenalty=10000\exhyphenpenalty=10000\popdisplay\fontsize{22}{24}\selectfont\color{popcream}\MakeUppercase{#1}\par}
  \end{tcolorbox}
}
% Section numérotée : pastille orange avec le numéro + titre Archivo
\newcounter{popsec}
\newcommand{\coursec}[1]{%
  \stepcounter{popsec}\setcounter{popsub}{0}%
  \par\vspace{16pt}\noindent
  \begin{minipage}{\linewidth}
  \tikz[baseline=-0.7ex]\node[draw=popink,line width=1.6pt,fill=poporange,rounded corners=4pt,minimum size=22pt,inner sep=0]{\popdisplay\fontsize{12}{12}\selectfont\color{popcream}\thepopsec};%
  \hspace{8pt}{\popdisplay\fontsize{15}{17}\selectfont\color{popink}\MakeUppercase{#1}}\par
  \vspace{4pt}\noindent{\color{poporange}\rule{36pt}{3.6pt}}
  \end{minipage}\par\nopagebreak\vspace{8pt}
}
\newcounter{popsub}
\newcommand{\courssub}[1]{%
  \stepcounter{popsub}%
  \par\vspace{9pt}\noindent{\popxbold\fontsize{11.5}{13}\selectfont\color{popink}{\color{poporange}\thepopsec.\thepopsub}\hspace{6pt}#1}\par\nopagebreak\vspace{4pt}
}

% ============================================================
% Boîtes pédagogiques — même recette : fond blanc, bordure épaisse colorée,
% ombre dure encre, badge plein. Titre optionnel [#1].
% ============================================================
\tcbset{popbase/.style={breakable,enhanced,colback=white,boxrule=2.3pt,arc=9pt,
  shadow={3pt}{-3pt}{0pt}{popink},left=14pt,right=14pt,top=11pt,bottom=11pt,before skip=11pt,after skip=15pt,
  fontupper=\popsemi\fontsize{10.6}{14.5}\selectfont\color{popink},
  fontlower=\popsemi\fontsize{10.6}{14.5}\selectfont\color{popink}}}
% Coche dessinée (le glyphe ✓ n'existe pas dans les polices du document)
\newcommand{\popcheck}[1]{\tikz[baseline=-0.5ex]\draw[#1,line width=1.6pt,line cap=round,line join=round](-0.13,0.0)--(-0.04,-0.09)--(0.13,0.1);}
\providecommand{\coche}{\popcheck{popgreen}}
\providecommand{\resultat}[1]{\par\smallskip\noindent\fcolorbox{popgreen}{popgreenL}{\parbox{\dimexpr\linewidth-2\fboxsep-2\fboxrule\relax}{\textbf{Result.} #1}}\par\smallskip}
\newcommand{\popboxhead}[3]{\poppill{#1}{#2}{white}\ifx\relax#3\relax\else\hspace{6pt}{\popxbold\fontsize{10.6}{12}\selectfont\color{popink}#3}\fi\par\vspace{7pt}}

\newtcolorbox{aretenir}[1][]{popbase,colframe=popgreen,before upper={\popboxhead{Key points}{popgreen}{#1}}}
\newtcolorbox{exemplebox}[1][]{popbase,colframe=poporange,before upper={\popboxhead{Exemple}{poporange}{#1}}}
\newtcolorbox{definition}[1][]{popbase,colframe=popblue,before upper={\popboxhead{Definition}{popblue}{#1}}}
\newtcolorbox{propriete}[1][]{popbase,colframe=poppurple,before upper={\popboxhead{Property}{poppurple}{#1}}}
\newtcolorbox{methode}[1][]{popbase,colframe=popgold,before upper={\popboxhead{Method}{popgold}{#1}}}
\newtcolorbox{attention}[1][]{popbase,colframe=poppink,before upper={\popboxhead{Warning}{poppink}{#1}}}
\newtcolorbox{experience}[1][]{popbase,colframe=popblue,colback=popblueL,before upper={\popboxhead{Experiment}{popblue}{#1}}}
% « Le savais-tu ? » : carte violette pleine (contexte, culture, Cameroun)
\newtcolorbox{savaistu}[1][]{popbase,colback=poppurple,colframe=popink,
  fontupper=\popsemi\fontsize{10.4}{14.2}\selectfont\color{white},
  before upper={\poppill{Did you know?}{white}{poppurple}\ifx\relax#1\relax\else\hspace{6pt}{\popxbold\color{white}#1}\fi\par\vspace{7pt}}}
% Exercice résolu : énoncé en haut, \tcblower puis la solution sur fond crème
\newcounter{popexo}\newcounter{popexr}
\newtcolorbox{exoresolu}[1][]{popbase,colframe=poporange,colbacklower=popcream,
  segmentation style={popink,line width=1.2pt,dashed},
  before upper={\stepcounter{popexr}\popboxhead{Worked example \thepopexr}{poporange}{#1}},
  before lower={\poppill{Solution}{popink}{popcream}\par\vspace{6pt}}}
% Exercice (non résolu en place) — la correction est dans la partie Corrigés
\newtcolorbox{exercice}[1][]{popbase,colframe=popink,
  before upper={\stepcounter{popexo}\popboxhead{Exercise \thepopexo}{popink}{#1}}}
% Corrigé
\newtcolorbox{corrige}[1][]{popbase,colframe=popgreen,colback=popgreenL,
  before upper={\popboxhead{Solution}{popgreen}{#1}}}
% Objectifs / prérequis / bilan
\newtcolorbox{objectifs}{popbase,colframe=popink,colback=poporangeL,
  before upper={\popboxhead{Chapter objectives}{popink}{}}}
\newtcolorbox{prerequis}{popbase,colframe=popink,colback=popblueL,
  before upper={\popboxhead{Prerequisites}{popblue}{}}}
\newtcolorbox{bilan}{popbase,colframe=popink,colback=white,boxrule=2.8pt,
  before upper={\popboxhead{Fiche bilan}{poporange}{L'essentiel à connaître pour l'examen}}}
% Figure encadrée avec légende
\newtcolorbox{popfigure}[1][]{enhanced,breakable=false,colback=white,colframe=popline,boxrule=1.4pt,arc=8pt,
  left=10pt,right=10pt,top=10pt,bottom=8pt,before skip=10pt,after skip=12pt,halign=center,
  lower separated=false,halign lower=center,
  fontlower=\popsemi\fontsize{9}{11.5}\selectfont\color{popmuted},
  #1}
\newcommand{\legende}[1]{\par\vspace{6pt}{\popsemi\fontsize{9}{11.5}\selectfont\color{popmuted}#1}}

% ============================================================
% Signature OnBuch+
% ============================================================
\newcommand{\poplogoinline}[1]{{\poplogo\fontsize{#1}{#1}\selectfont\color{popink}OnBuch\color{poporange}+}}
\newcommand{\signatureonbuch}{%
  \begin{tcolorbox}[enhanced,colback=popink,colframe=popink,boxrule=2.2pt,arc=10pt,
      drop shadow=poporange,left=14pt,right=14pt,top=10pt,bottom=10pt,before skip=0pt,after skip=0pt]
    \tikz[baseline=-0.7ex]\node[circle,fill=popgreen,draw=popcream,line width=1.3pt,minimum size=22pt,inner sep=0]{\popcheck{white}};%
    \hspace{9pt}\begin{minipage}[c]{0.62\linewidth}
      {\popxbold\fontsize{12}{13}\selectfont\color{popcream}Signed\ \textbullet\ {\poplogo OnBuch\color{poporange}+}}\par\vspace{2pt}
      {\raggedright\popsemi\fontsize{8}{10.5}\selectfont\color{popline}\MakeUppercase{OnBuch+ teaching team}\\\MakeUppercase{Follows the official MINESEC syllabus}\par}
    \end{minipage}\hfill
    {\poplogo\fontsize{22}{22}\selectfont\color{popcream}OnBuch\color{poporange}+}
  \end{tcolorbox}%
}

% ============================================================
% Page de garde
% #1 titre · #2 matière · #3 niveau · #4 module · #5 n° leçon · #6 durée ·
% #7 description / objectifs (texte ou itemize)
% ============================================================
\newcommand{\pagegarde}[7]{%
  \thispagestyle{empty}
  \newgeometry{a4paper,margin=1.7cm,top=1.6cm,bottom=1.4cm}%
  \begin{tikzpicture}[remember picture,overlay]
    \fill[poporange!18] ([xshift=-3.2cm,yshift=-3.6cm]current page.north east) circle (4.6cm);
    \fill[poppurple!14] ([xshift=2.4cm,yshift=7.6cm]current page.south west) circle (3.8cm);
  \end{tikzpicture}%
  {\poplogo\fontsize{30}{30}\selectfont\color{popink}OnBuch\color{poporange}+}\hfill
  \raisebox{0.6ex}{\poppill{Signed course \textbullet\ Premium}{popink}{popcream}}\par
  \vspace{1pt}\popsquiggle{poppurple}\par
  \vspace{8pt}
  \begin{tcolorbox}[enhanced,colback=poporange,colframe=popink,boxrule=2.8pt,arc=13pt,
      drop shadow=popink,left=18pt,right=18pt,top=12pt,bottom=13pt,after skip=10pt]
    \poppill{#2 \textbullet\ #3}{popink}{popcream}\hspace{5pt}\poppill{#4}{white}{popink}\par\vspace{8pt}
    {\popdisplay\fontsize{14}{14}\selectfont\color{popink}CHAPTER #5}\par\vspace{4pt}
    {\raggedright\hyphenpenalty=10000\exhyphenpenalty=10000\popdisplay\fontsize{25}{27}\selectfont\color{popcream}\MakeUppercase{#1}\par}
    \vspace{8pt}
    {\popxbold\fontsize{9.5}{9.5}\selectfont\color{popink}\MakeUppercase{Suggested time: #6}}
  \end{tcolorbox}
  \begin{tcolorbox}[enhanced,colback=white,colframe=popink,boxrule=2.2pt,arc=10pt,
      drop shadow=popink,left=16pt,right=16pt,top=10pt,bottom=10pt,after skip=8pt]
    \poppill{In this chapter}{poppurple}{white}\par\vspace{5pt}
    \popsemi\fontsize{9.3}{12.6}\selectfont\color{popink}#7
  \end{tcolorbox}
  \noindent\begin{minipage}[t]{0.487\linewidth}
    \begin{tcolorbox}[enhanced,colback=popgreen,colframe=popink,boxrule=2.2pt,arc=10pt,
        drop shadow=popink,left=11pt,right=11pt,top=10pt,bottom=10pt,height=3.1cm,valign=center]
      \tcbox[colback=white,colframe=popink,boxrule=1.3pt,arc=5pt,boxsep=0pt,nobeforeafter,
        left=3pt,right=3pt,top=3pt,bottom=3pt,valign=center]{\qrcode[height=1.9cm]{https://play.google.com/store/apps/details?id=com.luvvix.onbuch&hl=en}}%
      \hspace{8pt}\begin{minipage}[c]{0.5\linewidth}
        {\popdisplay\fontsize{11}{12.5}\selectfont\color{white}\MakeUppercase{Download OnBuch}}\par\vspace{4pt}
        {\raggedright\popsemi\fontsize{8.4}{11}\selectfont\color{white}Free \textbullet\ Google Play\\AI tutor, past papers, results\par}
      \end{minipage}
    \end{tcolorbox}
  \end{minipage}\hfill
  \begin{minipage}[t]{0.487\linewidth}
    \begin{tcolorbox}[enhanced,colback=poppurple,colframe=popink,boxrule=2.2pt,arc=10pt,
        drop shadow=popink,left=11pt,right=11pt,top=10pt,bottom=10pt,height=3.1cm,valign=center]
      \tikz[baseline=-0.7ex]\node[circle,fill=white,draw=popink,line width=1.3pt,minimum size=1.6cm,inner sep=0]{\popdisplay\fontsize{24}{24}\selectfont\color{poppurple}+};%
      \hspace{8pt}\begin{minipage}[c]{0.6\linewidth}
        {\popdisplay\fontsize{11}{12.5}\selectfont\color{white}\MakeUppercase{Go OnBuch+}}\par\vspace{4pt}
        {\raggedright\popsemi\fontsize{8.4}{11}\selectfont\color{white}All signed courses, unlimited AI tutor, PDF sheets — OnBuch+ tab in the app\par}
      \end{minipage}
    \end{tcolorbox}
  \end{minipage}
  \vspace{8pt}
  \popcontenu
  \vfill
  \signatureonbuch
  \restoregeometry
  \newpage
}

% Rangée « Ce cours contient » (page de garde)
\newcommand{\poptile}[3]{% #1 couleur #2 gros texte #3 légende
  \begin{tcolorbox}[enhanced,colback=#1,colframe=popink,boxrule=2pt,arc=9pt,shadow={2.5pt}{-2.5pt}{0pt}{popink},
      left=6pt,right=6pt,top=9pt,bottom=9pt,halign=center,valign=center,height=1.75cm,width=\linewidth,nobeforeafter]
    {\popdisplay\fontsize{12.5}{13}\selectfont\color{white}\MakeUppercase{#2}}\par\vspace{3pt}
    {\centering\popsemi\fontsize{7.8}{9.5}\selectfont\color{white}#3\par}
  \end{tcolorbox}}
\newcommand{\popcontenu}{%
  \par\noindent{\popxboldls\fontsize{7.5}{7.5}\selectfont\color{popmuted}THIS SIGNED COURSE CONTAINS}\par\vspace{7pt}
  \noindent\begin{minipage}[t]{0.235\linewidth}\poptile{popblue}{Course}{complete, illustrated, step by step}\end{minipage}\hfill
  \begin{minipage}[t]{0.235\linewidth}\poptile{poporange}{Methods}{and worked examples}\end{minipage}\hfill
  \begin{minipage}[t]{0.235\linewidth}\poptile{poppink}{Exercises}{exam style + solutions}\end{minipage}\hfill
  \begin{minipage}[t]{0.235\linewidth}\poptile{popgreen}{Summary sheet}{the essentials to remember}\end{minipage}\par}

% Sommaire
\newtcolorbox{sommaire}{enhanced,colback=white,colframe=popink,boxrule=2.2pt,arc=10pt,
  drop shadow=popink,left=18pt,right=18pt,top=13pt,bottom=13pt,before skip=6pt,after skip=20pt,
  before upper={\poppill{Contents}{poppurple}{white}\par\vspace{9pt}\popsemi\fontsize{10.6}{14.5}\selectfont\color{popink}}}

% Signature de fin de cours — poussée tout en bas de la dernière page
\newcommand{\findecours}{%
  \par\vfill\nopagebreak
  \begin{center}\popsquiggle{poporange}\end{center}\vspace{4pt}
  \signatureonbuch
}
\AtBeginDocument{\def\llbracket{\mathopen{[\mkern-3.5mu[}}\def\rrbracket{\mathclose{]\mkern-3.5mu]}}} % intervalles d'entiers (glyphe ⟦ absent de la police)
% Glyphes absents de Fira Math : redéfinis à partir de glyphes disponibles
\AtBeginDocument{%
  \def\setminus{\mathbin{\backslash}}\let\smallsetminus\setminus
  \def\vdots{\mathord{\vbox{\baselineskip3.2pt\lineskiplimit0pt\kern4pt\hbox{.}\hbox{.}\hbox{.}}}}%
  \def\ddots{\mathinner{\mkern1mu\raise6.5pt\hbox{.}\mkern2mu\raise3.6pt\hbox{.}\mkern2mu\raise0.7pt\hbox{.}\mkern1mu}}%
  \def\triangle{\mathord{\scalebox{0.9}{$\symup{\Delta}$}}}%
  \def\nabla{\mathord{\raisebox{\depth}{\scalebox{1}[-1]{$\symup{\Delta}$}}}}%
  \def\checkmark{\popcheck{popdark}}%
  \def\star{\mathbin{\ast}}\def\bigstar{\mathord{\ast}}%
  \def\diamond{\mathbin{\scalebox{0.6}{$\Diamond$}}}%
  \def\bigcup{\mathop{\mathchoice{\raisebox{-0.3em}{\scalebox{1.7}{$\cup$}}}{\scalebox{1.25}{$\cup$}}{\cup}{\cup}}\displaylimits}%
  \def\bigcap{\mathop{\mathchoice{\raisebox{-0.3em}{\scalebox{1.7}{$\cap$}}}{\scalebox{1.25}{$\cap$}}{\cap}{\cap}}\displaylimits}%
  \def\lesseqqgtr{\mathrel{\lessgtr}}\def\bigoplus{\mathop{\oplus}\displaylimits}%
}
\providecommand{\ssi}{if and only if} % abréviation fréquente
\AtBeginDocument{\providecommand{\wideparen}{\overparen}} % arc de cercle

% Fonctions usuelles que les modèles utilisent sans que amsmath les fournisse
\providecommand{\arsinh}{\operatorname{arsinh}}\providecommand{\arcosh}{\operatorname{arcosh}}
\providecommand{\artanh}{\operatorname{artanh}}\providecommand{\arsech}{\operatorname{arsech}}
\providecommand{\arcsch}{\operatorname{arcsch}}\providecommand{\arcoth}{\operatorname{arcoth}}
\providecommand{\arcsinh}{\operatorname{arsinh}}\providecommand{\arccosh}{\operatorname{arcosh}}
\providecommand{\arctanh}{\operatorname{artanh}}\providecommand{\sech}{\operatorname{sech}}
\providecommand{\csch}{\operatorname{csch}}\providecommand{\arcsec}{\operatorname{arcsec}}
\providecommand{\arccsc}{\operatorname{arccsc}}\providecommand{\arccot}{\operatorname{arccot}}
\providecommand{\sgn}{\operatorname{sgn}}\providecommand{\lcm}{\operatorname{lcm}}
\providecommand{\Var}{\operatorname{Var}}\providecommand{\Cov}{\operatorname{Cov}}
\providecommand{\permil}{\ensuremath{{}^{0}\!/_{00}}}\providecommand{\textperthousand}{\ensuremath{{}^{0}\!/_{00}}}

%% FIGFIX-BEGIN v5 — correctifs globaux des figures (docs/onbuch-generation/tools/figfix)
\usepackage{adjustbox}\usepackage{etoolbox}\tcbuselibrary{raster}
% toute figure TikZ/pgfplots plus large que la ligne est réduite à la largeur disponible
\BeforeBeginEnvironment{tikzpicture}{\begin{adjustbox}{max width=\linewidth}}
\AfterEndEnvironment{tikzpicture}{\end{adjustbox}}
% légendes pgfplots lisibles par-dessus les courbes
\pgfplotsset{every axis/.append style={legend style={fill=white,fill opacity=0.92,text opacity=1,draw=black!15,inner sep=3pt}}}
% étiquettes d'arêtes (syntaxe edge["..."]) avec fond blanc
\tikzset{every edge quotes/.append style={fill=white,inner sep=1.2pt}}
%% FIGFIX-END
\begin{document}

\pagegarde{ECONOMIC GEOLOGY}{Geology}{Upper Sixth}{Upper Sixth — Geology}{6}{15 periods}{This chapter explores how geological materials become economic resources: the formation, occurrence and exploitation of coal, petroleum and natural gas, the industrial uses of rocks, and the engineering geology of dams, tunnels and slopes. It closes with the prospection, exploration and mining methods — and their environmental consequences — that underpin Cameroon's extractive economy.
\vspace{4pt}
\begin{multicols}{2}\small
\begin{itemize}
\item By the end of this chapter I can explain the formation of coal, petroleum and natural gas from organic matter.
\item By the end of this chapter I can describe the source rock, reservoir rock, cap rock and trap of a petroleum system.
\item By the end of this chapter I can distinguish the ranks of coal and the types of petroleum traps with labelled diagrams.
\item By the end of this chapter I can state the economic uses of common rocks and industrial minerals in Cameroon.
\item By the end of this chapter I can explain the geological criteria for siting dams, reservoirs and tunnels.
\item By the end of this chapter I can analyse the factors controlling slope stability and propose stabilisation measures.
\end{itemize}\end{multicols}}

\begin{sommaire}
\begin{multicols}{2}
\begin{enumerate}
\item Coal: Origin, Ranks and Occurrence
\item Petroleum and Natural Gas: The Petroleum System
\item Drilling, Production and Integration of Fossil Fuels
\item Applications of Geological Materials and Economic Uses of Rocks
\item Engineering Geology: Dams, Reservoirs, Tunnels and Slope Stability
\item Prospection and Exploration Methods
\item Mining Methods and the Environment
\item Integration activity
\item Methods and worked examples
\item Exercises
\item Solutions to the exercises
\item Summary sheet
\end{enumerate}
\end{multicols}
\end{sommaire}

\begin{objectifs}
\begin{itemize}
\item By the end of this chapter I can explain the formation of coal, petroleum and natural gas from organic matter.
\item By the end of this chapter I can describe the source rock, reservoir rock, cap rock and trap of a petroleum system.
\item By the end of this chapter I can distinguish the ranks of coal and the types of petroleum traps with labelled diagrams.
\item By the end of this chapter I can state the economic uses of common rocks and industrial minerals in Cameroon.
\item By the end of this chapter I can explain the geological criteria for siting dams, reservoirs and tunnels.
\item By the end of this chapter I can analyse the factors controlling slope stability and propose stabilisation measures.
\item By the end of this chapter I can describe geological, geophysical and geochemical prospection and exploration methods.
\item By the end of this chapter I can compare surface and underground mining methods and evaluate their environmental impacts.
\end{itemize}
\end{objectifs}

\begin{prerequis}
\begin{itemize}
\item Rock cycle and classification of igneous, sedimentary and metamorphic rocks (Lower Sixth)
\item Sedimentary structures, stratification and principles of stratigraphy
\item Folds, faults and unconformities (structural geology)
\item Weathering, erosion and mass wasting processes
\item Map reading and interpretation of geological cross-sections
\end{itemize}
\end{prerequis}

\newpage

\coursec{Coal: Origin, Ranks and Occurrence}

In 1977, the discovery of oil at Kole near Rio-del-Rey transformed Cameroon into a petroleum-producing nation, and today the SONARA refinery, the Logbaba gas project in Douala and the Lom Pangar dam all depend on one science: \cle{economic geology}. Yet every rainy season, landslides block the Bamenda--Fundong road and quarry blasts shake houses in Yaound\'e, reminding us that the same rocks that enrich a nation can also threaten it. How do geologists find coal, oil and minerals hidden kilometres underground? How do engineers decide where a dam will hold and where a slope will fail? This chapter takes you from the organic origins of fossil fuels to the geophysical surveys, mining methods and environmental safeguards that turn geology into wealth --- and safety --- for Cameroon.

We begin with the oldest industrial fossil fuel of all: \cle{coal}. Before we can search for it, we must answer a deceptively simple question: how does a forest become a black rock that burns?

\courssub{Fossil Fuels and the Organic Origin of Coal}

\begin{definition}[Fossil fuel]
A \cle{fossil fuel} is any combustible material formed by the transformation of buried organic matter (plants or microscopic organisms) over geological time. The three great fossil fuels are \textbf{coal} (from land plants), \textbf{petroleum} and \textbf{natural gas} (mainly from marine plankton and algae). They are called ``fossil'' because they store solar energy captured by photosynthesis millions of years ago.
\end{definition}

Pick up a lump of coal and you hold a fossilised swamp. Under the microscope, thin sections of coal reveal preserved plant tissues: spores, cuticles, woody fragments, resin bodies. This is direct evidence that coal is a \emph{biogenic sedimentary rock} --- a rock made of accumulated organic debris, just as sandstone is made of accumulated sand grains.

The starting material is \textbf{vegetation growing in wetlands}: swamps, bogs, mires and peatlands. In a tropical coastal swamp --- picture the mangrove-fringed estuaries of the Rio-del-Rey or the Niger Delta --- plants grow fast and die fast. Normally, dead vegetation is attacked by fungi, bacteria and oxygen, and decays completely back to \ce{CO2} and water. For coal to form, this decay must be \emph{interrupted}.

\begin{propriete}[Conditions for preservation of plant debris]
Plant debris is preserved and converted toward coal only when:
\begin{enumerate}
\item \textbf{Waterlogging} excludes oxygen, creating \cle{anaerobic} (oxygen-free) conditions that stop aerobic bacteria and fungi from destroying the organic matter;
\item \textbf{Accumulation exceeds decay}: plant production is rapid and continuous, so debris piles up faster than the slow anaerobic decay can decompose it;
\item \textbf{Rapid burial} by sediment (sand, mud, silt) seals the organic layer away from the atmosphere and further oxidation;
\item \textbf{Subsidence} of the basin keeps the swamp near water level for long periods, allowing thick peat to accumulate vertically.
\end{enumerate}
\end{propriete}

\begin{attention}[Why deserts and mountains make no coal]
Coal is not simply ``old wood''. A forest that burns or rots in the open air leaves \emph{no} coal, because oxidation destroys the carbon. The essential ingredient is the \emph{swamp}: standing water is the lid that keeps oxygen out. This is why coal is associated with \textbf{deltaic and coastal plain environments}, not with uplands or well-drained terrains.
\end{attention}

\courssub{Coal Petrology: Macerals and Microlithotypes}

Before discussing rank, geologists describe coal at the microscopic level. Just as minerals are the building blocks of igneous rocks, \cle{macerals} are the microscopic organic constituents of coal, identifiable by their reflectance, morphology and origin.

\begin{definition}[Maceral groups]
Three major maceral groups are recognised:
\begin{itemize}
\item \textbf{Vitrinite} (from woody tissues, cellulose and lignin): most abundant in most coals; shows medium to high reflectance; the main contributor to coking properties.
\item \textbf{Liptinite (exinite)} (from spores, pollen, cuticles, resins, algae): hydrogen-rich, low reflectance, high volatile yield; important for oil and gas generation.
\item \textbf{Inertinite} (from oxidised/charred material, fungal remains): high reflectance, low volatile matter, inert during coking; derived from forest fires or severe oxidation in the swamp.
\end{itemize}
\end{definition}

\begin{propriete}[Microlithotypes]
Associations of macerals form \cle{microlithotypes}: \textbf{vitrite} ( > 95\% vitrinite), \textbf{liptite}, \textbf{inertite}, \textbf{clarite} (vitrinite + liptinite), \textbf{durite} (inertinite + liptinite), \textbf{virtinite} (vitrinite + inertinite), \textbf{trimacerite} (all three). These associations control the physical behaviour of coal during mining, crushing and carbonisation.
\end{propriete}

\begin{aretenir}[Macerals and rank]
Vitrinite reflectance (\%Ro) is the \textbf{standard rank parameter} in modern coal petrology. It increases systematically with coalification and is used internationally to correlate rank stages.
\end{aretenir}

\courssub{Coalification: From Peat to Anthracite}

Once buried, the peat is squeezed and heated as more sediment piles on top. This progressive transformation is called \cle{coalification}.

\begin{definition}[Coalification]
\cle{Coalification} is the set of physical and chemical changes by which buried plant debris is transformed, under increasing pressure and temperature over geological time, through the series \textbf{peat $\rightarrow$ lignite $\rightarrow$ bituminous coal $\rightarrow$ anthracite}. Each stage is called a \cle{rank} of coal.
\end{definition}

The chemical essence of coalification is simple: the original plant material (essentially cellulose and lignin, rich in C, H and O) progressively \textbf{loses hydrogen and oxygen} --- as water \ce{H2O}, carbon dioxide \ce{CO2} and methane \ce{CH4} --- and becomes \textbf{richer and richer in carbon}. Physically, the material is \textbf{compacted}: a peat layer about 10~m thick may yield a coal seam only about 1~m thick, a compaction ratio of roughly 10:1.

\begin{propriete}[Rank increases with burial and temperature]
The \cle{rank} of a coal increases with the \textbf{depth of burial}, the \textbf{geothermal gradient} (temperature), and the \textbf{duration} of burial. As rank increases: carbon content rises, moisture and volatile matter fall, and the \textbf{calorific value} (heat released per kilogramme on burning) rises. (The direction of these trends is examinable; no numerical proof is required.)
\end{propriete}

\begin{center}
\begin{tabularx}{\linewidth}{|l|X|X|X|X|}
\hline
\rowcolor{popblueL} \textbf{Rank} & \textbf{Fixed carbon (dry, ash-free approx.)} & \textbf{Volatile matter (dry, ash-free)} & \textbf{Aspect} & \textbf{Calorific value (MJ/kg)} \\
\hline
Peat & $<$ 60 \% & $>$ 60 \% & Brown, fibrous, visible plant remains, very wet & $<$ 15 \\
\hline
Lignite & 60--70 \% & 45--55 \% & Brown-black, soft, earthy, high moisture & 15--25 \\
\hline
Sub-bituminous & 70--75 \% & 35--45 \% & Dull black, harder, lower moisture & 25--30 \\
\hline
Bituminous coal & 75--90 \% & 15--35 \% & Black, banded, hard, often coking & 30--35 \\
\hline
Anthracite & $>$ 90 \% & $<$ 15 \% & Black, shiny, hard, conchoidal fracture & $>$ 35 \\
\hline
\end{tabularx}
\end{center}

\begin{popfigure}
\begin{center}
\begin{tikzpicture}[scale=1]
% boxes
\draw[fill=popgreenL, draw=popdark, rounded corners] (0,0) rectangle (2.6,1.4);
\node[align=center] at (1.3,0.7) {\textbf{Peat}\\ \small C $<$ 60 \%};
\draw[fill=popgold!40, draw=popdark, rounded corners] (3.4,0) rectangle (6.0,1.4);
\node[align=center] at (4.7,0.7) {\textbf{Lignite}\\ \small C 60--70 \%};
\draw[fill=poporange!30, draw=popdark, rounded corners] (6.8,0) rectangle (9.4,1.4);
\node[align=center] at (8.1,0.7) {\textbf{Sub-bit.}\\ \small C 70--75 \%};
\draw[fill=popmuted!50, draw=popdark, rounded corners] (10.2,0) rectangle (12.8,1.4);
\node[align=center] at (11.5,0.7) {\textbf{Bituminous}\\ \small C 75--90 \%};
\draw[fill=popink!80, draw=popdark, rounded corners] (13.6,0) rectangle (16.2,1.4);
\node[align=center, text=white] at (14.9,0.7) {\textbf{Anthracite}\\ \small C $>$ 90 \%};
% arrows between
\draw[-{Stealth}, very thick, poporange] (2.6,0.7) -- (3.4,0.7);
\draw[-{Stealth}, very thick, poporange] (6.0,0.7) -- (6.8,0.7);
\draw[-{Stealth}, very thick, poporange] (9.4,0.7) -- (10.2,0.7);
\draw[-{Stealth}, very thick, poporange] (12.8,0.7) -- (13.6,0.7);
% top arrow: carbon and calorific value
\draw[-{Stealth}, very thick, popblue] (1.3,2.3) -- (14.9,2.3);
\node[above, popblue, align=center] at (8.1,2.3) {\small Increasing carbon content and calorific value};
% bottom arrow: depth, temperature
\draw[-{Stealth}, very thick, poppurple] (1.3,-0.6) -- (14.9,-0.6);
\node[below, poppurple, align=center] at (8.1,-0.6) {\small Increasing depth of burial, temperature and time};
% losses
\node[align=center, popdark] at (3.0,1.9) {\small $-$\ce{H2O}};
\node[align=center, popdark] at (6.4,1.9) {\small $-$\ce{H2O}, $-$\ce{CO2}};
\node[align=center, popdark] at (9.8,1.9) {\small $-$\ce{CH4}, volatiles};
\node[align=center, popdark] at (13.2,1.9) {\small $-$\ce{H2}, graphitisation};
\end{tikzpicture}
\end{center}
\legende{The coalification series. Each step expels water and volatile compounds, so the residue becomes richer in carbon and releases more heat per kilogramme. Sub-bituminous coal is an intermediate rank recognised in international classifications (ASTM).}
\end{popfigure}

\begin{methode}[Estimating coal rank from proximate analysis]
In the laboratory, a coal sample is analysed by \textbf{proximate analysis}, giving moisture, volatile matter, ash and fixed carbon. To assign a rank:
\begin{enumerate}
\item Air-dry the sample, then heat to $\SI{105}{\ensuremath{{}^{\circ}\mathrm{C}}}$ to constant mass: mass loss = \textbf{moisture}.
\item Heat the dried sample to about $\SI{950}{\ensuremath{{}^{\circ}\mathrm{C}}}$ in a closed crucible (7 min): mass loss = \textbf{volatile matter} (VM).
\item Burn the residue in air at $\SI{750}{\ensuremath{{}^{\circ}\mathrm{C}}}$: mass loss = \textbf{fixed carbon} (FC); remainder = \textbf{ash}.
\item Calculate on a \textbf{dry, ash-free (daf) basis}:
\[
\%FC_{\mathrm{daf}} = \frac{FC}{FC+VM}\times 100,\qquad
\%VM_{\mathrm{daf}} = \frac{VM}{FC+VM}\times 100
\]
\item Compute the \textbf{fuel ratio} $= \mathrm{FC}/\mathrm{VM}$ (on daf basis).
\item Read the rank: fuel ratio $< 1$ suggests lignite; $1$--$4$ suggests bituminous coal; $> 4$ (very low volatiles) suggests anthracite. High fixed carbon and low volatile matter always mean high rank.
\end{enumerate}
\end{methode}

\begin{exoresolu}[Assigning a rank to a coal sample]
A dried coal sample of mass $50.0\ \mathrm{g}$ loses $12.5\ \mathrm{g}$ of volatile matter on heating, and leaves $5.0\ \mathrm{g}$ of ash after complete combustion. Determine its fixed carbon, its fuel ratio, and its probable rank.
\tcblower
\textbf{Solution.}
\begin{align*}
m_{\mathrm{FC}} &= 50.0 - 12.5 - 5.0 = 32.5\ \mathrm{g}\\
\text{Fuel ratio (dry basis)} &= \dfrac{\mathrm{FC}}{\mathrm{VM}} = \dfrac{32.5}{12.5} = 2.6
\end{align*}
Fixed carbon on a dry basis: $\dfrac{32.5}{50.0}\times 100 = 65\ \%$ of the dried sample. On a dry, ash-free basis:
\[
\%FC_{\mathrm{daf}} = \frac{32.5}{32.5+12.5}\times 100 = 72.2\ \%,\quad
\%VM_{\mathrm{daf}} = 27.8\ \%
\]
A fuel ratio of $2.6$ (between 1 and 4) with substantial volatiles ($27.8\%$ daf) and fixed carbon $72.2\%$ daf indicates a \textbf{bituminous coal} --- a good steam and coking coal, not an anthracite.
\end{exoresolu}

\begin{attention}[Peat is NOT a coal]
A classic GCE trap: candidates write ``the ranks of coal are peat, lignite, bituminous coal and anthracite''. \textbf{Peat is not a true coal.} It is unconsolidated, waterlogged plant debris --- the \emph{parent material} of coal, not a rank of it. Lignite (brown coal) is the \emph{first true coal}. Saying ``peat is the lowest rank of coal'' loses marks; say ``peat is the precursor of coal''.
\end{attention}

\courssub{Coal Seams and Coal-Bearing Sedimentary Environments}

Coal does not occur as random lumps: it forms laterally continuous beds called \cle{coal seams}, interlayered with ordinary sediments. A succession of alternating sandstone, shale and coal seams is called \cle{coal measures}. This interbedding tells a story: the environment \emph{oscillated} between swamp (coal), river channel or flood (sandstone), and quiet floodplain or lake (shale).

The typical home of coal measures is the \textbf{deltaic coastal plain}: a broad, slowly subsiding lowland crossed by distributary rivers, where swamps occupy the abandoned parts of the delta between active channels. Fluvial (river) floodplains and lagoons behind barrier beaches also host coal. Repeated cycles of subsidence and sediment supply produce the rhythmic repetition of seams --- some coal fields contain dozens of seams stacked vertically.

\begin{popfigure}
\begin{center}
\begin{tikzpicture}[scale=1]
% basin fill layers
\fill[popgold!45] (0,0) -- (0,0.8) -- (12,0.5) -- (12,0) -- cycle;
\fill[black!80] (0,0.8) -- (12,0.5) -- (12,0.9) -- (0,1.2) -- cycle;
\fill[popmuted!45] (0,1.2) -- (12,0.9) -- (12,1.9) -- (0,2.2) -- cycle;
\fill[popgold!45] (0,2.2) -- (12,1.9) -- (12,2.7) -- (0,3.0) -- cycle;
\fill[black!80] (0,3.0) -- (12,2.7) -- (12,3.1) -- (0,3.4) -- cycle;
\fill[popmuted!45] (0,3.4) -- (12,3.1) -- (12,4.0) -- (0,4.3) -- cycle;
% labels
\node[right, align=left] at (12.15,3.55) {\small Shale (floodplain/lake)};
\node[right, align=left] at (12.15,2.9) {\small \textbf{Coal seam B}};
\node[right, align=left] at (12.15,2.3) {\small Sandstone (channel fill)};
\node[right, align=left] at (12.15,1.55) {\small Shale (overbank)};
\node[right, align=left] at (12.15,0.85) {\small \textbf{Coal seam A}};
\node[right, align=left] at (12.15,0.3) {\small Sandstone (basal)};
% basement
\fill[popline!60] (0,0) -- (12,0) -- (12,-0.7) -- (0,-0.7) -- cycle;
\node at (6,-0.35) {\small Basement (subsiding)};
% swamp on top
\draw[popgreen, thick] (0,4.3) -- (12,4.0);
\node[popgreen!60!black, align=center] at (3,4.75) {\small Swamp: peat forming today};
\draw[-{Stealth}, popdark] (6,-1.1) -- (6,-1.7);
\node[below, align=center] at (6,-1.7) {\small Subsidence};
\end{tikzpicture}
\end{center}
\legende{Cross-section of a coal-bearing deltaic basin. Coal seams (black) are interbedded with sandstone and shale; continued subsidence buries today's peat to form tomorrow's seam.}
\end{popfigure}

\begin{aretenir}[Key points on coal occurrence]
\begin{itemize}
\item Coal is a \textbf{sedimentary rock of organic origin}, found in \textbf{seams} within coal measures.
\item Coal measures record \textbf{deltaic, fluvial and lagoonal} environments with repeated subsidence.
\item Seam thickness, continuity and rank all depend on the depositional environment and burial history.
\item \textbf{Roof} and \textbf{floor} rocks (shale, sandstone, fireclay/seatearth) control mining conditions and methane emission.
\end{itemize}
\end{aretenir}

\courssub{Uses of Coal}

\begin{itemize}
\item \textbf{Fuel}: burnt in thermal power stations to raise steam and generate electricity; also used for domestic heating and in cement works.
\item \textbf{Coke for metallurgy}: when bituminous coal is heated without air (\emph{destructive distillation} / \cle{carbonisation}), volatiles are expelled and the porous, nearly pure carbon residue called \cle{coke} remains. Coke is the fuel and reducing agent of the blast furnace, converting iron ore (\ce{Fe2O3}) to iron.
\item \textbf{Chemical feedstock}: coal tar and coal gas from distillation yield dyes, drugs, fertilisers, plastics and synthetic fuels (e.g., Fischer--Tropsch process).
\item \textbf{Other}: activated carbon, carbon fibres, silicon carbide, soil conditioner (humic acids from lignite).
\end{itemize}

\courssub{Coal in Cameroon and Africa (GCE-useful extra)}

Cameroon has \textbf{no large exploited coalfield}: its known occurrences are thin, low-rank (lignite to sub-bituminous) seams in Cretaceous to Tertiary sedimentary basins, for example in the \textbf{Mamfe basin} (Mamfe Group), parts of the \textbf{Douala basin} (Mundeck Formation) and \textbf{Garoua basin} (Polí Formation), which remain of limited economic interest compared with hydrocarbons. By contrast, Nigeria's \textbf{Enugu coalfield} (Anambra basin, of Cretaceous age) contains workable bituminous seams that long fed Nigerian railways and the Oji River power station. Elsewhere in Africa, the great Karoo-age (Permian) coalfields of \textbf{South Africa, Zimbabwe, Botswana and Mozambique} rank among the world's largest. The lesson for the GCE candidate: coal distribution is controlled by \emph{geological history} --- the presence of ancient swampy basins and their preservation --- not by present-day climate alone.

\begin{savaistu}[Did you know?]
The same Cretaceous rifting that opened the South Atlantic created a chain of sedimentary basins along the West African margin --- from the Benue Trough through the Rio-del-Rey to Gabon. Some of these basins trapped coal and organic-rich shales; others, more deeply buried and hotter, generated the oil and gas that Cameroon exploits today. Coal and petroleum are thus geological cousins, separated mainly by the type of organic matter (terrestrial vs. marine) and the temperature of burial.
\end{savaistu}

\begin{exercice}[Consolidation]
\begin{enumerate}
\item Define coalification and list its four stages in order of increasing rank (excluding peat).
\item Explain why waterlogging is essential to coal formation.
\item A dried sample of $80.0\ \mathrm{g}$ yields $8.0\ \mathrm{g}$ of volatiles and $4.0\ \mathrm{g}$ of ash. Calculate the fixed carbon, the fuel ratio (dry basis), the daf fixed carbon percentage, and assign a probable rank.
\item Describe the sedimentary environment in which coal measures form, and name the rock types typically interbedded with coal seams.
\item Name the three maceral groups and give one characteristic of each.
\end{enumerate}
\end{exercice}

\begin{corrige}[Exercise 1]
\begin{enumerate}
\item Coalification is the transformation of buried plant debris into coal under rising pressure and temperature. Stages (increasing rank): lignite $\rightarrow$ sub-bituminous $\rightarrow$ bituminous coal $\rightarrow$ anthracite (peat being the precursor, not a true coal).
\item Waterlogging excludes oxygen; anaerobic conditions stop complete decay, so carbon-rich debris accumulates instead of being oxidised to \ce{CO2} and water.
\item $m_{\mathrm{FC}} = 80.0 - 8.0 - 4.0 = 68.0\ \mathrm{g}$; fuel ratio (dry) $= 68.0/8.0 = 8.5$. Dry, ash-free basis: $\%FC_{\mathrm{daf}} = \frac{68}{68+8}\times 100 = 89.5\%$, $\%VM_{\mathrm{daf}} = 10.5\%$. Ratio $> 4$ with very low volatiles ($10.5\%$ daf): \textbf{anthracite}.
\item Subsiding deltaic or fluvial coastal plain with swamps; coal seams interbedded with sandstone (channel fills, crevasse splays) and shale (floodplain, lacustrine, lagoonal). Seatearth (fireclay) often underlies seams.
\item Vitrinite (woody origin, medium-high reflectance, coking); Liptinite/exinite (spores, cuticles, resins, low reflectance, high H, oil-prone); Inertinite (oxidised/charred, high reflectance, inert in coking).
\end{enumerate}
\end{corrige}

\coursec{Petroleum and Natural Gas: The Petroleum System}

Having explored the solid fossil fuel \emph{coal} in the previous section, we now turn to the \emph{fluid} fossil fuels: petroleum (crude oil) and natural gas. While coal forms from land plants in swampy environments, petroleum and natural gas have a predominantly \emph{marine} origin. They are the lifeblood of modern industrial society, powering transport, generating electricity, and providing feedstock for the petrochemical industry (plastics, fertilizers, pharmaceuticals). Understanding their genesis, migration, and entrapment constitutes the science of \emph{petroleum geology}. For Cameroon, this is not abstract theory: the Rio-del-Rey and Douala/Kribi-Campo basins are active petroleum provinces, and the Limbe refinery (SONARA) processes crude oil daily. Mastering the \cle{petroleum system} concept is essential for your GCE examination and for appreciating our national energy resources.

\begin{aretenir}[Learning Objectives for This Section]
By the end of this section, you must be able to:
\begin{enumerate}
    \item Explain the \textbf{organic theory} of petroleum origin and describe the formation of \cle{kerogen} in \cle{source rocks}.
    \item Distinguish the \textbf{four types of kerogen} (I, II, III, IV) and predict their hydrocarbon potential (oil vs. gas).
    \item Define the \textbf{oil window} and \textbf{gas window} in terms of temperature, depth, and geothermal gradient; calculate subsurface temperatures.
    \item Describe \textbf{primary} and \textbf{secondary migration} mechanisms and pathways.
    \item List the \textbf{four essential elements} of a petroleum system (Source, Reservoir, Seal, Trap) and explain the \textbf{Critical Moment}.
    \item Differentiate \textbf{porosity} (storage) from \textbf{permeability} (flow) and evaluate reservoir quality.
    \item Classify traps into \textbf{structural} (anticline, fault, salt dome) and \textbf{stratigraphic} (pinch-out, unconformity, reef) types; sketch and label cross-sections.
    \item Draw \textbf{horizontal fluid contacts} (GOC, OWC) correctly on dipping/folded layers.
    \item Distinguish \textbf{associated}, \textbf{non-associated}, and \textbf{condensate} natural gas.
    \item Apply knowledge to \textbf{Cameroon's basins} (Rio-del-Rey, Douala/Kribi-Campo): source rocks, reservoirs, seals, traps, maturity, infrastructure.
    \item Solve quantitative problems: maturity temperature, buoyancy pressure, column heights.
\end{enumerate}
\end{aretenir}

\courssub{The Organic Theory: From Marine Plankton to Kerogen}

Unlike the outdated \emph{inorganic} (abiogenic) theories, the universally accepted \cle{organic theory} states that petroleum and natural gas originate from the remains of once-living organisms.

\begin{experience}[The Microscopic Origin]
Imagine a warm, shallow epicontinental sea, perhaps similar to the Cretaceous seas that covered parts of the Douala Basin \SI{100}{\mega\year} ago. The surface waters teem with microscopic life: \emph{phytoplankton} (diatoms, coccolithophores) and \emph{zooplankton} (foraminifera, copepods). When these organisms die, their organic matter (lipids, proteins, carbohydrates) rains down through the water column.
\end{experience}

\begin{definition}[Source Rock and Kerogen]
A \cle{source rock} is a fine-grained sedimentary rock (typically \emph{shale} or \emph{marl}) rich in organic matter (\textgreater 0.5\% Total Organic Carbon, TOC) that has generated or can generate hydrocarbons. The insoluble organic residue dispersed in the rock is called \cle{kerogen}. Kerogen is a complex, high-molecular-weight macromolecule; it is the direct precursor of petroleum and natural gas.
\end{definition}

\paragraph{Preservation is key.} In well-oxygenated bottom waters, bacteria decompose the organic rain completely. For a source rock to form, the organic matter must be \emph{preserved}. This requires \cle{anoxic (oxygen-poor) bottom conditions}, often found in restricted basins (like the Black Sea today) or during periods of high sea level and stratification. Rapid burial by clay minerals also protects the organic matter from oxidation. The \cle{Total Organic Carbon (TOC)} content is the primary screening parameter: \textgreater 0.5\% for gas-prone, \textgreater 1\% for oil-prone, \textgreater 2\% for excellent source rocks.

\paragraph{Types of Kerogen.} The chemical composition of kerogen depends on the source organisms and determines the \emph{type} of hydrocarbons generated. The \cle{Van Krevelen diagram} (H/C vs O/C atomic ratios) classifies kerogen:

\begin{center}
\begin{tabularx}{\linewidth}{|l|X|X|X|X|}\hline
\rowcolor{popblueL}
\textbf{Type} & \textbf{Origin / Setting} & \textbf{Key Chemistry} & \textbf{Hydrocarbon Potential} & \textbf{Cameroon Example} \\ \hline
\textbf{I (Algal/Sapropelic)} & Lipid-rich algae (\emph{Botryococcus}, \emph{Tasmanites}); lacustrine / restricted marine anoxic & High H/C (\textgreater 1.5), Low O/C & \textbf{Oil-prone} (high yield); waxy crudes & Barremian lacustrine shales (Douala/Kribi-Campo) \\ \hline
\textbf{II (Mixed Marine)} & Marine plankton (algae + zooplankton); normal marine anoxic & Moderate H/C (1.2--1.5), Low O/C & \textbf{Oil \& Gas}; most common worldwide & Cretaceous marine shales (Aptian-Albian) in both basins \\ \hline
\textbf{III (Humic/Woody)} & Terrestrial higher plants (lignin, cellulose); deltaic, coastal plain & Low H/C (\textless 1.0), High O/C & \textbf{Gas-prone} (dry gas); associated with coal & Tertiary deltaic shales (Rio-del-Rey), Logbaba Fm (Douala) \\ \hline
\textbf{IV (Inert/Residual)} & Highly oxidized, reworked, recycled organic matter & Very low H/C, High O/C & \textbf{None} (spent) & Weathered sediments, basement \\ \hline
\end{tabularx}
\end{center}

\begin{attention}[Exam Precision: Kerogen vs. Coal]
Do not confuse \textbf{Type III kerogen} with \textbf{coal}. Coal is a \emph{rock} (\textgreater 50\% organic matter by volume); Type III kerogen is \emph{dispersed organic matter} in a clastic matrix (shale/sandstone). Both derive from terrestrial plants, but coal is a distinct lithology. In exams, "source rock for gas" = Type III kerogen shale OR coal measures.
\end{attention}

\courssub{Maturation: The Oil Window and Gas Window}

Burial subjects the source rock to increasing temperature and pressure. This thermal stress cracks the large kerogen molecules into smaller hydrocarbon molecules — a process called \cle{catagenesis} (or maturation). The preceding stage, \cle{diagenesis} (\textless \SI{50}{\ensuremath{{}^{\circ}\mathrm{C}}}), involves biogenic methane generation by bacteria.

\begin{propriete}[The Oil and Gas Windows -- Examinable Ranges]
Hydrocarbon generation occurs within specific temperature ranges, primarily controlled by burial depth and the \cle{geothermal gradient} (typically \SI{25--30}{\ensuremath{{}^{\circ}\mathrm{C}}\per\kilo\meter} in sedimentary basins; higher near the Cameroon Volcanic Line).
\begin{itemize}
    \item \textbf{Diagenesis (\textless \SI{50}{\ensuremath{{}^{\circ}\mathrm{C}}}, \textless \SI{2}{\kilo\meter}):} Biogenic methane (shallow gas). Kerogen chemically unchanged. Vitrinite Reflectance $R_o < 0.5\%$.
    \item \textbf{Oil Window (\SIrange{60}{120}{\ensuremath{{}^{\circ}\mathrm{C}}}, approx. \SI{2--4}{\kilo\meter} depth):} Thermal cracking of kerogen (catagenesis) produces liquid petroleum (oil) and associated wet gas. \textbf{Peak oil generation} at \SI{90--100}{\ensuremath{{}^{\circ}\mathrm{C}}} ($R_o \approx 0.85--1.0\%$).
    \item \textbf{Wet Gas / Condensate Window (\SIrange{120}{150}{\ensuremath{{}^{\circ}\mathrm{C}}}, \SI{4--5}{\kilo\meter}):} Oil cracks further to wet gas (C$_2$--C$_5$) and condensate. $R_o \approx 1.0--1.3\%$.
    \item \textbf{Dry Gas Window (\textgreater \SI{150}{\ensuremath{{}^{\circ}\mathrm{C}}}, \textgreater \SI{5}{\kilo\meter}):} All liquid hydrocarbons crack to methane (dry gas). Kerogen becomes \emph{post-mature} (graphitic residue). $R_o > 1.3\%$.
    \item \textbf{Metagenesis (\textgreater \SI{200}{\ensuremath{{}^{\circ}\mathrm{C}}}):} Only methane remains; source rock is spent. $R_o > 2.0\%$.
\end{itemize}
\textbf{Examinable note:} You must be able to state the temperature/depth ranges for the oil and gas windows and explain the role of the geothermal gradient. A high gradient (e.g., near volcanic areas like the Cameroon Volcanic Line) brings the windows shallower; a low gradient pushes them deeper. \textbf{Vitrinite Reflectance ($R_o$)} is the standard maturity indicator; know the approximate $R_o$ values for the oil window onset (0.5\%), peak (0.9\%), and gas window onset (1.3\%).
\end{propriete}

\begin{popfigure}
\begin{tikzpicture}[scale=0.85]
    % Axes
    \draw[->, thick] (0,0) -- (0,7) node[left] {Depth (km)};
    \draw[->, thick] (0,0) -- (10,0) node[below] {Temperature (\textdegree C)};
    % Geothermal gradient line (30 deg/km)
    \draw[popdark, very thick] (0,0) -- (6.5, 6.5) node[above right, pos=0.9] {Geothermal Gradient (30°C/km)};
    % Depth markers
    \foreach \d/\t in {0/0, 1/30, 2/60, 3/90, 4/120, 5/150, 6/180} {
        \draw (\t, -0.1) -- (\t, 0.1) node[below, font=\tiny] {\d};
        \draw (-0.1, \d) -- (0.1, \d) node[left, font=\tiny] {\t};
    }
    % Windows as shaded bands
    \fill[poporange!30] (60,0) rectangle (120,6);
    \fill[popred!20] (120,0) rectangle (150,6);
    \fill[popblue!20] (150,0) rectangle (210,6);
    % Labels
    \node[poporange, font=\bfseries, align=center] at (90, 5.5) {OIL WINDOW\\ (60--120°C)};
    \node[popred, font=\bfseries, align=center] at (135, 5.5) {WET GAS /\\ CONDENSATE};
    \node[popblue, font=\bfseries, align=center] at (180, 5.5) {DRY GAS WINDOW\\ (>150°C)};
    \node[popgreen, font=\bfseries, align=center] at (30, 5.5) {DIAGENESIS\\ (Biogenic Gas)};
    % Maturity indicators
    \draw[dashed, popmuted] (60,0) -- (60,6);
    \draw[dashed, popmuted] (120,0) -- (120,6);
    \draw[dashed, popmuted] (150,0) -- (150,6);
    \node[font=\tiny, align=center] at (30, 0.5) {Immature\\ $R_o<0.5\%$};
    \node[font=\tiny, align=center] at (90, 0.5) {Mature (Oil)\\ $R_o=0.5-1.0\%$};
    \node[font=\tiny, align=center] at (135, 0.5) {Late Mature\\ $R_o=1.0-1.3\%$};
    \node[font=\tiny, align=center] at (180, 0.5) {Post-Mature\\ $R_o>1.3\%$};
\end{tikzpicture}
\legende{Oil and Gas Windows: Depth vs Temperature for a typical geothermal gradient of 30°C/km. The "Oil Window" (60--120°C) is the prime target for liquid petroleum exploration. Vitrinite Reflectance ($R_o$) scale added.}
\end{popfigure}

\courssub{Migration: From Source to Trap}

Oil and gas are less dense than the formation water saturating the sedimentary pores. Once generated, they are squeezed out of the compacting source rock by \cle{primary migration} (expulsion) and move buoyantly upwards through permeable pathways — \cle{secondary migration} — until they are trapped or seep at the surface.

\begin{definition}[Migration]
\begin{itemize}
    \item \textbf{Primary Migration (Expulsion):} Movement of hydrocarbons \emph{out of} the source rock into an adjacent permeable carrier bed. Mechanisms: (1) Micro-fracturing due to overpressure from hydrocarbon generation (volume increase), (2) Oil-phase expulsion as a separate buoyant phase, (3) Diffusion (minor). It is a short-distance (cm to m), inefficient process — only 10--50\% of generated HC may expel; the rest remains as \emph{residual hydrocarbons} or cracks to pyrobitumen.
    \item \textbf{Secondary Migration:} Long-distance movement (metres to kilometres) through permeable \cle{carrier beds} (sandstones, fractured carbonates, fault zones) along the hydraulic gradient (buoyancy + hydrodynamic flow). Oil floats on water; gas floats on oil. Migration is \emph{buoyancy-driven} (vertical) but can be deflected laterally by hydrodynamic flow or permeability barriers.
\end{itemize}
\end{definition}

\begin{attention}[Common Mistake: Underground Lakes]
\textbf{Never} describe or draw oil and gas accumulating in underground "lakes", "pools", or "caves". This is a cardinal error in GCE exams. Hydrocarbons occupy the \emph{pore spaces} (intergranular pores, fractures, vugs) of the reservoir rock, typically saturating 50--80\% of the pore volume (the rest is irreducible water). The "pool" is a \emph{volume of rock}, not a void.
\end{attention}

\courssub{The Petroleum System: Reservoir, Seal, and Trap}

For a commercial accumulation to exist, four essential elements must coincide in time and space. This is the \cle{Petroleum System} concept (Magoon \& Dow, 1994).

\begin{definition}[The Petroleum System (Magoon \& Dow, 1994)]
A \cle{petroleum system} comprises:
\begin{enumerate}
    \item \textbf{Source Rock:} Generates hydrocarbons (mature, adequate TOC, correct kerogen type).
    \item \textbf{Reservoir Rock:} Stores hydrocarbons (porous \& permeable).
    \item \textbf{Seal (Cap Rock):} Prevents escape (impermeable, high capillary entry pressure).
    \item \textbf{Trap:} Geometry that focuses hydrocarbons into an accumulation.
\end{enumerate}
\textbf{Critical Moment:} The time when all elements coincide (trap exists, source is mature, migration pathways open). \textbf{Preservation Time:} Period between critical moment and present (risk of leakage, biodegradation, tectonic breach).
\end{definition}

\paragraph{1. Reservoir Rock.}
Must have high \cle{porosity} (storage capacity) and high \cle{permeability} (flow capacity).

\begin{propriete}[Porosity vs Permeability -- Fundamental Distinction]
\textbf{Porosity ($\phi$)} = Volume of voids / Total volume (\%). \textbf{Permeability ($k$)} = Ability to transmit fluids (Darcies, D; or millidarcies, mD; 1 D = $9.87\times10^{-13}$ m$^2$).
\begin{itemize}
    \item \emph{A rock can be highly porous but impermeable} (e.g., unfractured shale $\phi\sim30\%$, $k\sim0.001$ mD; pumice; some chalks — pores not connected).
    \item \emph{A rock can have low porosity but high permeability} (e.g., tightly cemented but well-fractured granite/quartzite; $\phi\sim2\%$, $k\sim100$ mD via fractures).
\end{itemize}
\textbf{Good reservoirs:} Clean, well-sorted sandstones (intergranular porosity, $\phi>15\%$, $k>100$ mD); fractured/karstified limestones/dolomites (secondary porosity, $\phi>5\%$, $k$ highly variable). \textbf{Effective porosity} excludes isolated pores; only connected pores contribute to flow.
\end{propriete}

\paragraph{2. Seal (Cap Rock).}
An impermeable layer overlying the reservoir. \cle{Shale} (clay alignment, capillary entry pressure $\sim$1--10 MPa), \cle{evaporites} (anhydrite, salt — ultimate seals, entry pressure $\sim$10--50+ MPa), tight carbonates. Capillary entry pressure must exceed buoyancy pressure of the hydrocarbon column. Thickness matters: a \SI{10}{\meter} shale may seal a small column; a \SI{100}{\meter} salt seals a giant field. \textbf{Seal capacity} is quantified by the maximum hydrocarbon column height it can hold: $h_{max} = P_{entry} / (g \Delta\rho)$.

\paragraph{3. Traps.}
A trap is a geometric arrangement of reservoir and seal that allows hydrocarbon accumulation. \cle{Structural traps} (deformation) and \cle{stratigraphic traps} (depositional/erosional) are the two main classes. \textbf{Combination traps} involve both.

\begin{popfigure}
\begin{tikzpicture}[scale=0.9]
    % Anticlinal Trap Cross-Section
    \def\layerA{popblue!30} % Cap rock
    \def\layerB{popyellow!50} % Reservoir
    \def\layerC{popgreen!30} % Source rock
    \def\layerD{popgray!30} % Basement
    
    % Layers (bottom up)
    % Basement
    \fill[\layerD] (-5,-3) -- (5,-3) -- (5,-2.5) -- (-5,-2.5) -- cycle;
    % Source Rock (mature)
    \fill[\layerC] (-5,-2.5) -- (5,-2.5) -- (5,-1.5) -- (-5,-1.5) -- cycle;
    % Reservoir (folded)
    \fill[\layerB] 
        (-5,-1.5) .. controls (-3,-1) and (3,-1) .. (5,-1.5) 
        -- (5,-0.5) .. controls (3,0) and (-3,0) .. (-5,-0.5) -- cycle;
    % Cap Rock (folded)
    \fill[\layerA] 
        (-5,-0.5) .. controls (-3,0) and (3,0) .. (5,-0.5) 
        -- (5,0.5) .. controls (3,1) and (-3,1) .. (-5,0.5) -- cycle;
    % Overburden
    \fill[popgray!10] (-5,0.5) -- (5,0.5) -- (5,2) -- (-5,2) -- cycle;
    
    % Hydrocarbon contacts in reservoir
    \fill[poporange!80] 
        (-3.5, -0.55) .. controls (-2,-0.2) and (2,-0.2) .. (3.5, -0.55) 
        -- (3.5, -0.7) .. controls (2,-0.85) and (-2,-0.85) .. (-3.5, -0.7) -- cycle; % Gas
    \fill[popred!70] 
        (-4.2, -0.7) .. controls (-2.5,-1.1) and (2.5,-1.1) .. (4.2, -0.7) 
        -- (4.2, -1.1) .. controls (2.5,-1.3) and (-2.5,-1.3) .. (-4.2, -1.1) -- cycle; % Oil
    % Water below
    \fill[popblue!50] 
        (-5, -1.5) -- (5, -1.5) 
        -- (4.2, -1.1) .. controls (2.5,-1.3) and (-2.5,-1.3) .. (-4.2, -1.1) -- cycle;
    
    % Contacts lines
    \draw[popdark, dashed, thick] (-4.2, -0.7) .. controls (-2.5,-1.1) and (2.5,-1.1) .. (4.2, -0.7) node[right, pos=0.9, font=\small] {Oil-Water Contact (OWC)};
    \draw[popdark, dashed, thick] (-3.5, -0.55) .. controls (-2,-0.2) and (2,-0.2) .. (3.5, -0.55) node[right, pos=0.9, font=\small] {Gas-Oil Contact (GOC)};
    
    % Labels
    \node[font=\small\bfseries, align=center] at (-6, 1.2) {Overburden};
    \node[font=\small\bfseries, align=center, popblue] at (-6, 0) {Cap Rock (Seal)\\(Shale/Evaporite)};
    \node[font=\small\bfseries, align=center, poporange] at (-6, -1) {Reservoir Rock\\(Sandstone/Limestone)};
    \node[font=\small\bfseries, align=center, popgreen] at (-6, -2) {Source Rock\\(Mature Shale)};
    \node[font=\small\bfseries, align=center, popgray] at (-6, -2.8) {Basement};
    
    % Migration arrows
    \draw[->, thick, popred] (-4.5, -2) -- (-4.5, -1.6) node[midway, left, font=\tiny] {Primary Migration};
    \draw[->, thick, popred] (-4.5, -1.6) .. controls (-4, -1.2) and (-3, -1) .. (-2.5, -0.8) node[midway, above, font=\tiny] {Secondary Migration};
    
    % Well
    \draw[very thick, popdark] (0, 2) -- (0, -0.6);
    \node[font=\small\bfseries] at (0, 2.3) {Exploration Well};
\end{tikzpicture}
\legende{Cross-section of an anticlinal trap. Note the fluid contacts (GOC, OWC) are horizontal (gravity equilibrium), while the rock layers are folded. The well targets the crest (culmination) for maximum hydrocarbon column.}
\end{popfigure}

\courssub{Structural Traps}

Formed by deformation \emph{after} reservoir deposition. Timing is critical: trap must exist \emph{before} or \emph{during} migration (Critical Moment).

\begin{itemize}
    \item \textbf{Anticlines / Domes:} Classic convex-up folds. Hydrocarbons migrate to the \cle{culmination} (highest point). Most giant fields (e.g., Middle East, North Sea). \textbf{Drill the crest!} Plunging anticlines have a \cle{spill point} — the lowest point on the closure where hydrocarbons leak out.
    \item \textbf{Fault Traps:} Reservoir juxtaposed against seal across a fault. \emph{Normal fault:} downthrown reservoir against upthrown seal. \emph{Reverse fault:} upthrown reservoir against downthrown seal. Seal capacity depends on \cle{shale gouge ratio (SGR)} in the fault zone (volume of shale in fault rock / total fault rock). SGR \textgreater 20--30\% generally seals.
    \item \textbf{Salt Domes / Diapirs:} Mobile salt (evaporite) pierces overlying sediments, creating traps on its flanks (subsalt, suprasalt, flank traps). Very important in Gulf of Mexico, Angola, and the \emph{Douala/Kribi-Campo Basin} (Aptian salt). Salt welds (where salt has completely withdrawn) can create leakage paths.
    \item \textbf{Drape / Compaction Anticlines:} Compaction over a basement high or reef creates a structural closure.
\end{itemize}

\courssub{Stratigraphic Traps}

Formed by \emph{depositional} or \emph{erosional} processes, contemporaneous with or prior to reservoir formation. No structural deformation required. Often subtle on seismic; require detailed stratigraphic analysis.

\begin{itemize}
    \item \textbf{Pinch-out (Lateral Facies Change):} Porous reservoir (e.g., sandstone bar, reef, turbidite lobe) pinches out laterally into impermeable shale. \emph{Example:} Cretaceous sandstone pinch-outs in the \emph{Rio-del-Rey Basin}.
    \item \textbf{Unconformity Traps:} Reservoir truncated by an erosion surface, overlain by seal. Hydrocarbons migrate updip to the subcrop edge. \emph{Example:} Pre-Cretaceous basement reservoirs sealed by Cretaceous shales in Douala Basin.
    \item \textbf{Reef / Bioherm Traps:} Porous carbonate buildups (reefs) encased in shale/marl. High porosity but often complex permeability (mouldic, vuggy). \emph{Example:} Albian reefs in the Logbaba area (Douala Basin).
    \item \textbf{Lenticular / Channel Sands:} Isolated sand bodies (fluvial channels, turbidite lobes) encased in shale. Stratigraphic closure.
    \item \textbf{Diagenetic Traps:} Porosity enhancement (dolomitization, leaching) or destruction (cementation) creates lateral permeability barriers.
\end{itemize}

\begin{popfigure}
\begin{tikzpicture}[scale=0.75]
    % Gallery of 4 traps
    \def\res{popyellow!60}
    \def\seal{popblue!30}
    \def\source{popgreen!30}
    \def\base{popgray!30}
    \def\hc{poporange!80}
    
    % 1. Anticline
    \begin{scope}[xshift=0cm, yshift=0cm]
        \fill[\base] (-1.5,-1.5) rectangle (1.5,-1);
        \fill[\source] (-1.5,-1) rectangle (1.5,-0.5);
        \fill[\res] (-1.5,-0.5) .. controls (-1,0) and (1,0) .. (1.5,-0.5) -- (1.5,0) .. controls (1,0.5) and (-1,0.5) .. (-1.5,0) -- cycle;
        \fill[\seal] (-1.5,0) .. controls (-1,0.5) and (1,0.5) .. (1.5,0) -- (1.5,1) .. controls (1,1.5) and (-1,1.5) .. (-1.5,1) -- cycle;
        \fill[\hc] (-0.8,0) .. controls (-0.5,0.3) and (0.5,0.3) .. (0.8,0) -- (0.8,-0.1) .. controls (0.5,-0.2) and (-0.5,-0.2) .. (-0.8,-0.1) -- cycle;
        \draw[->, thick, popred] (0,-0.5) -- (0,-0.1);
        \node[font=\bfseries\small] at (0,1.3) {Anticline};
    \end{scope}
    
    % 2. Fault Trap
    \begin{scope}[xshift=4.5cm, yshift=0cm]
        \fill[\base] (-1.5,-1.5) rectangle (1.5,-1);
        \fill[\source] (-1.5,-1) rectangle (1.5,-0.5);
        % Footwall (left)
        \fill[\res] (-1.5,-0.5) -- (0,-0.5) -- (0.2,0.5) -- (-1.3,0.5) -- cycle;
        \fill[\seal] (-1.5,0.5) -- (0.2,0.5) -- (0.4,1.5) -- (-1.1,1.5) -- cycle;
        % Hanging wall (right) - seal against reservoir
        \fill[\seal] (0.2,0.5) -- (1.5,0.5) -- (1.5,1.5) -- (0.4,1.5) -- cycle;
        \fill[\res] (0,-0.5) -- (1.5,-0.5) -- (1.5,0.5) -- (0.2,0.5) -- cycle;
        \fill[\hc] (-0.2,0.3) -- (0.15,0.3) -- (0.15,-0.2) -- (-0.2,-0.2) -- cycle;
        \draw[thick, popdark] (0,-0.5) -- (0.2,0.5) node[midway, above right, font=\tiny, rotate=45] {Fault Plane};
        \draw[->, thick, popred] (0.5,-0.5) -- (0.5,0.2);
        \node[font=\bfseries\small] at (0,1.8) {Fault Trap};
    \end{scope}
    
    % 3. Salt Dome
    \begin{scope}[xshift=9cm, yshift=0cm]
        \fill[\base] (-1.5,-1.5) rectangle (1.5,-1);
        \fill[\source] (-1.5,-1) rectangle (1.5,-0.5);
        % Salt (pink)
        \fill[poppink!40] (-0.5,-0.5) .. controls (-0.5,0) and (0.5,0) .. (0.5,-0.5) -- (0.5,1.5) .. controls (0.5,1) and (-0.5,1) .. (-0.5,1.5) -- cycle;
        % Reservoir draped
        \fill[\res] (-1.5,-0.5) .. controls (-1,0) and (-0.5,0.5) .. (-0.5,0.5) -- (-0.5,1) .. controls (-0.5,1.2) and (-1,1.2) .. (-1.5,1) -- cycle;
        \fill[\res] (0.5,0.5) .. controls (1,0.5) and (1.5,0) .. (1.5,-0.5) -- (1.5,1) .. controls (1.5,1.2) and (1,1.2) .. (0.5,1) -- cycle;
        \fill[\seal] (-1.5,1) -- (1.5,1) -- (1.5,1.5) -- (-1.5,1.5) -- cycle;
        \fill[\hc] (-1.2,0.8) .. controls (-1,1) and (-0.7,1) .. (-0.55,1) -- (-0.55,0.8) .. controls (-0.7,0.7) and (-1,0.7) .. (-1.2,0.8) -- cycle;
        \fill[\hc] (1.2,0.8) .. controls (1,1) and (0.7,1) .. (0.55,1) -- (0.55,0.8) .. controls (0.7,0.7) and (1,0.7) .. (1.2,0.8) -- cycle;
        \draw[->, thick, popred] (-0.5,-0.5) .. controls (-0.5,0) and (-0.5,0.5) .. (-0.55,0.9);
        \draw[->, thick, popred] (0.5,-0.5) .. controls (0.5,0) and (0.5,0.5) .. (0.55,0.9);
        \node[font=\bfseries\small] at (0,1.8) {Salt Dome};
    \end{scope}
    
    % 4. Pinch-out
    \begin{scope}[xshift=13.5cm, yshift=0cm]
        \fill[\base] (-1.5,-1.5) rectangle (1.5,-1);
        \fill[\source] (-1.5,-1) rectangle (1.5,-0.5);
        \fill[\seal] (-1.5,0.5) rectangle (1.5,1.5);
        % Reservoir pinching right
        \fill[\res] (-1.5,-0.5) -- (1.5,-0.5) -- (1.5,0.5) -- (-0.5,0.5) -- cycle;
        \fill[\hc] (-1.2,0.3) -- (1.2,0.3) -- (1.2,0.1) -- (-1.2,0.1) -- cycle;
        \draw[->, thick, popred] (0,-0.5) -- (0,0.1);
        \node[font=\bfseries\small] at (0,1.8) {Pinch-out (Stratigraphic)};
    \end{scope}
\end{tikzpicture}
\legende{Gallery of four major trap types. (1) Anticline: structural, most common. (2) Fault trap: reservoir sealed by fault juxtaposition. (3) Salt dome: salt diapir creates flank traps. (4) Pinch-out: stratigraphic, reservoir thins into seal.}
\end{popfigure}

\courssub{Natural Gas: Associated and Non-Associated}

Natural gas is predominantly \cle{methane (CH$_4$)}, with varying amounts of ethane (C$_2$), propane (C$_3$), butane (C$_4$), and heavier hydrocarbons (condensate, C$_5$+), plus non-hydrocarbon gases (CO$_2$, N$_2$, H$_2$S). \cle{Hydrogen sulfide (H$_2$S)} makes gas "sour" (toxic, corrosive); "sweet" gas lacks H$_2$S.

\begin{definition}[Gas Classification]
\begin{itemize}
    \item \textbf{Associated Gas:} Gas produced in association with crude oil. Exists as \emph{gas cap} (free gas) above the oil column and \emph{solution gas} (dissolved in oil at reservoir pressure). Released at surface (separator). Also called "wet gas" if rich in C$_2$+ (condensate yield \textgreater 20 bbl/MMscf).
    \item \textbf{Non-Associated Gas:} Gas from reservoirs containing \emph{only} gas (dry gas) and water. Originates from Type III kerogen (coal measures) or over-mature Type II kerogen (dry gas window). No liquid petroleum at reservoir conditions.
    \item \textbf{Gas Condensate:} High-pressure gas rich in light liquids (C$_5$+). At surface conditions (lower P, T), liquids "condense" out (retrograde condensation). High value; requires special surface processing.
    \item \textbf{Biogenic Gas:} Shallow, bacterial methane ($\delta^{13}$C \textless -60\permil). \textbf{Thermogenic Gas:} Deep, thermal cracking ($\delta^{13}$C \textgreater -50\permil). Isotopes distinguish them.
\end{itemize}
\end{definition}

\begin{aretenir}[Key Points: Gas in the Petroleum System]
\begin{itemize}
    \item Gas is generated at higher maturity than oil (or from Type III kerogen at lower maturity).
    \item Gas migrates faster (lower viscosity) and can leak through seals that hold oil (lower capillary entry pressure for gas).
    \item \textbf{Lake Nyos (Cameroon) warning:} CO$_2$ gas accumulation in crater lakes is \emph{volcanic/magmatic}, not petroleum-related. Do not confuse the two in exams.
    \item \textbf{Gas-Oil Ratio (GOR):} Key parameter. Solution GOR (scf/STB) determines if reservoir is oil (GOR \textless 20,000), gas-condensate (20,000--100,000), or dry gas (\textgreater 100,000).
\end{itemize}
\end{aretenir}

\courssub{Cameroon's Petroleum Basins: GCE Context}

Cameroon has two main sedimentary basins with proven petroleum systems. Knowledge of their names, ages, and characteristics is frequently tested.

\begin{center}
\begin{tabularx}{\linewidth}{|l|X|X|}\hline
\rowcolor{popblueL}
\textbf{Feature} & \textbf{Rio-del-Rey Basin} & \textbf{Douala / Kribi-Campo Basin} \\ \hline
\textbf{Location} & Northwest, offshore/onshore border with Nigeria & Southwest, onshore \& offshore (Douala, Kribi, Campo) \\ \hline
\textbf{Tectonic Setting} & Divergent margin (Benue Trough / Gulf of Guinea) & Transform margin (Equatorial Atlantic) \\ \hline
\textbf{Key Source Rocks} & Cretaceous marine shales (Type II/III); Tertiary deltaic shales (Type III) & Cretaceous marine shales (Aptian-Albian, Type II); Lacustrine shales (Barremian, Type I) \\ \hline
\textbf{Reservoirs} & Miocene sands (deltaic), Cretaceous sands & Cretaceous sands (fluvial/deltaic), Fractured Basement, Albian Reefs \\ \hline
\textbf{Seals} & Shales, overpressured shales & Aptian Salt (ultimate seal), Cretaceous/Tertiary shales \\ \hline
\textbf{Traps} & Structural (rollover anticlines, faults), Stratigraphic (pinch-outs) & Structural (salt-related, fault blocks), Stratigraphic (reefs, unconformities) \\ \hline
\textbf{Maturity} & Shallow (high geothermal gradient near CVL) & Variable; deep offshore mature, onshore often immature/early mature \\ \hline
\textbf{Key Fields} & Gas dominant (e.g., Sanaga Sud, Moudi) & Oil \& Gas (Ebome, Kole, Mvia, Logbaba gas) \\ \hline
\textbf{Infrastructure} & Feeds SONARA (via pipeline) & Feeds SONARA (Limbe Refinery) \\ \hline
\end{tabularx}
\end{center}

\begin{savaistu}[SONARA and the Cameroon Volcanic Line]
The \emph{Société Nationale de Raffinage (SONARA)} in Limbe (Victoria) is Cameroon's sole refinery, commissioned in 1981. It processes \emph{light, sweet crude} (API \textgreater 35\ensuremath{{}^{\circ}}, low sulfur) from the Rio-del-Rey and Douala basins (and imported crudes). The \emph{Cameroon Volcanic Line (CVL)} influences the petroleum systems: its high heat flow matures source rocks at shallower depths (good for gas), but excessive heat can over-mature them (destroying oil potential). The CVL also provides CO$_2$ for the killer lakes (Nyos, Monoun) — a distinct geological hazard unrelated to petroleum systems.
\end{savaistu}

\courssub{Worked Example: Interpreting a Petroleum System Cross-Section}

\begin{exoresolu}[Reading an Oil Trap Cross-Section]
\textbf{Statement:} The figure below shows a schematic cross-section of a sedimentary basin. The geothermal gradient is \SI{30}{\ensuremath{{}^{\circ}\mathrm{C}}\per\kilo\meter}. Surface temperature is \SI{20}{\ensuremath{{}^{\circ}\mathrm{C}}}. The source rock (dark grey) is a Type II kerogen shale. The reservoir (yellow) is a sandstone. The seal (blue) is a shale. A fault cuts the section.
\begin{enumerate}
    \item Calculate the temperature at the top of the source rock (depth \SI{3.5}{\kilo\meter}) and state if it is in the oil window.
    \item Identify the trap type at location A (anticlinal crest) and location B (fault juxtaposition).
    \item Where would you drill a well to test the gas cap, oil leg, and water leg? Mark on the diagram.
    \item Explain why the fault at location B might leak.
\end{enumerate}

\begin{tikzpicture}[scale=0.7]
    % Layers
    \fill[popgray!20] (-5,0) rectangle (5,-1); % Overburden
    \fill[popblue!30] (-5,-1) rectangle (5,-1.5); % Seal
    \fill[popyellow!50] (-5,-1.5) .. controls (-3,-1) and (3,-1) .. (5,-1.5) -- (5,-2.5) .. controls (3,-2) and (-3,-2) .. (-5,-2.5) -- cycle; % Reservoir (anticline)
    \fill[popgreen!30] (-5,-2.5) rectangle (5,-3.5); % Source
    \fill[popgray!40] (-5,-3.5) rectangle (5,-4); % Basement
    % Fault
    \draw[thick, popred] (2,-1) -- (3,-4);
    % Hydrocarbons
    \fill[poporange!80] (-2.5,-1.5) .. controls (-1.5,-1.2) and (1.5,-1.2) .. (2.5,-1.5) -- (2.5,-1.7) .. controls (1.5,-1.4) and (-1.5,-1.4) .. (-2.5,-1.7) -- cycle; % Gas
    \fill[popred!70] (-3.5,-1.7) .. controls (-2,-2) and (2,-2) .. (3.5,-1.7) -- (3.5,-2.2) .. controls (2,-2.4) and (-2,-2.4) .. (-3.5,-2.2) -- cycle; % Oil
    \fill[popblue!50] (-5,-2.5) -- (5,-2.5) -- (3.5,-2.2) .. controls (2,-2.4) and (-2,-2.4) .. (-3.5,-2.2) -- cycle; % Water
    % Contacts
    \draw[dashed, thick] (-3.5,-1.7) .. controls (-2,-2) and (2,-2) .. (3.5,-1.7);
    \draw[dashed, thick] (-2.5,-1.5) .. controls (-1.5,-1.2) and (1.5,-1.2) .. (2.5,-1.5);
    % Labels
    \node at (-5.5, -0.5) {Seal};
    \node at (-5.5, -2) {Reservoir};
    \node at (-5.5, -3) {Source};
    \node at (3.5, -2.5) [font=\bfseries] {A};
    \node at (3.5, -1) [font=\bfseries] {B};
    \draw[<->] (5.5,0) -- (5.5,-3.5) node[midway, right] {3.5 km};
\end{tikzpicture}

\tcblower
\textbf{Detailed Solution:}
\begin{enumerate}
    \item \textbf{Temperature at \SI{3.5}{\kilo\meter}:} $T = T_0 + (\text{gradient} \times \text{depth}) = 20 + (30 \times 3.5) = 20 + 105 = \SI{125}{\ensuremath{{}^{\circ}\mathrm{C}}}$.
    This is \emph{above} the oil window (\SI{120}{\ensuremath{{}^{\circ}\mathrm{C}}}) and inside the \emph{wet gas / condensate window}. The source rock is \textbf{late mature for oil, peak for wet gas}. It would generate gas/condensate, not much oil. $R_o \approx 1.1\%$.
    \item \textbf{Trap A:} Anticlinal crest (structural trap, four-way dip closure). \textbf{Trap B:} Fault trap (reservoir on left juxtaposed against seal on right across the fault).
    \item \textbf{Drilling locations:}
        \begin{itemize}
            \item \textbf{Gas Cap:} Drill at the \textbf{culmination (crest)} of the anticline (Location A, high point of yellow layer).
            \item \textbf{Oil Leg:} Drill on the \textbf{flanks} of the anticline, below the GOC but above the OWC.
            \item \textbf{Water Leg:} Drill \textbf{downdip} (far left or right on the structure), below the OWC.
        \end{itemize}
    \item \textbf{Fault Seal Risk at B:} The fault juxtaposes reservoir (permeable) against seal (shale). For a seal to hold, the \textbf{Shale Gouge Ratio (SGR)} in the fault zone must be high enough (typically \textgreater 20--30\%). If the fault displacement is large relative to the shale thickness, the fault zone may contain too much sand/sandstone cataclasite, creating a permeable pathway (fault leak). Also, clay smearing must be continuous. If the fault cuts the seal entirely, it leaks.
\end{enumerate}
\end{exoresolu}

\begin{methode}[Reading an Oil-Trap Cross-Section (Exam Technique)]
\begin{enumerate}
    \item \textbf{Identify the layers:} Colour/pattern code: Source (dark), Reservoir (porous pattern), Seal (tight pattern), Overburden.
    \item \textbf{Find the Trap Geometry:} Anticline? Fault? Pinch-out? Salt? Label it.
    \item \textbf{Locate Fluid Contacts:} Draw \textbf{horizontal} lines for GOC (Gas-Oil Contact) and OWC (Oil-Water Contact). They are \emph{always horizontal} (equipotential surfaces), cutting across dipping beds.
    \item \textbf{Determine Hydrocarbon Column:} Vertical thickness between GOC and OWC (oil column) and crest to GOC (gas column).
    \item \textbf{Select Drill Site:} Target the \textbf{culmination} (highest structural point) for gas/oil contact; target \textbf{flank} for oil only; avoid spill point.
    \item \textbf{Check Timing:} Did the trap exist \emph{before} migration? (Critical Moment). If the fault moved \emph{after} charge, it might leak.
\end{enumerate}
\end{methode}

\courssub{Integration Exercise: The Petroleum System in a Basin Context}

\begin{exercice}[Basin Analysis Synthesis]
\textbf{Statement:} A petroleum exploration company acquires seismic data over a block in the Douala Basin. The interpretation reveals a large, elongated anticline cored by a salt pillow (Aptian salt). The reservoir target is the \emph{Melange Formation} (Cretaceous fractured limestone). The seal is the \emph{Cap Shale} (Cenomanian-Turonian). The source rock is the \emph{Logbaba Shale} (Albian, Type II kerogen). Present-day geothermal gradient is \SI{28}{\ensuremath{{}^{\circ}\mathrm{C}}\per\kilo\meter}. Surface temp \SI{25}{\ensuremath{{}^{\circ}\mathrm{C}}}. Top of Logbaba Shale at \SI{3.2}{\kilo\meter}. Top of Melange Reservoir at \SI{2.8}{\kilo\meter}.
\begin{enumerate}
    \item Calculate the temperature at the top of the source rock. Is it mature for oil? For gas?
    \item Describe the complete petroleum system elements (Source, Reservoir, Seal, Trap, Timing).
    \item The seismic shows the salt pillow grew actively during the Cretaceous. Discuss the impact on (a) maturation history and (b) trap integrity.
    \item A well drilled at the crest of the anticline encounters \SI{150}{m} of gas, \SI{300}{m} of oil, and water below. Calculate the pressure difference (buoyancy) between the crest and the OWC if oil density is \SI{0.85}{g/cm^3} and water density is \SI{1.05}{g/cm^3}. ($g = \SI{9.81}{m.s^{-2}}$).
\end{enumerate}
\end{exercice}

\begin{corrige}[Exercise: Basin Analysis Synthesis]
\begin{enumerate}
    \item \textbf{Temp at Source Rock Top:} $T = 25 + (28 \times 3.2) = 25 + 89.6 = \SI{114.6}{\ensuremath{{}^{\circ}\mathrm{C}}}$.
    \textbf{Maturity:} Inside the \textbf{Oil Window} (\SIrange{60}{120}{\ensuremath{{}^{\circ}\mathrm{C}}}). Peak oil generation (\SI{~90--100}{\ensuremath{{}^{\circ}\mathrm{C}}}) has passed; it is in the \textbf{late oil / early wet gas window}. Good for oil + associated gas. $R_o \approx 0.95\%$.
    \item \textbf{Petroleum System Elements:}
    \begin{itemize}
        \item \textbf{Source:} Logbaba Shale (Albian, Type II, TOC \textgreater 2\%), mature ($T \approx \SI{115}{\ensuremath{{}^{\circ}\mathrm{C}}}$).
        \item \textbf{Reservoir:} Melange Formation (Cretaceous fractured limestone). Porosity: fracture/vuggy. Permeability: fracture-dependent.
        \item \textbf{Seal:} Cap Shale (Cenomanian-Turonian marine shale). Thick, regional. Also Aptian Salt (ultimate seal at depth).
        \item \textbf{Trap:} Salt-cored anticline (structural/halokinetic). Four-way dip closure at crest.
        \item \textbf{Timing (Critical Moment):} Salt movement (pillow growth) created the anticline \emph{during} the Late Cretaceous. Logbaba Shale entered oil window in Late Cretaceous/Early Tertiary. \textbf{Trap formed before/during charge $\rightarrow$ FAVOURABLE.}
    \end{itemize}
    \item \textbf{Salt Pillow Impact:}
    \begin{itemize}
        \item (a) \textbf{Maturation:} Salt has high thermal conductivity ($\sim \SI{5--7}{\watt\per\meter\per\kelvin}$ vs $\sim \SI{1.5}{\watt\per\meter\per\kelvin}$ for shale). It acts as a \textbf{heat pipe}, drawing heat upwards. This \textbf{accelerates maturation} of the source rock adjacent to the salt (shallower depth for same maturity). The "maturity high" over the salt pillow is a classic exploration play.
        \item (b) \textbf{Trap Integrity:} Active salt growth creates \textbf{radial faults} and \textbf{extensional fractures} on the crest (crestal collapse graben). These can breach the seal! However, the salt itself is the ultimate seal. Risk: fault leakage at crest; opportunity: fault-assisted migration up the flanks.
    \end{itemize}
    \item \textbf{Buoyancy Pressure ($\Delta P$):}
    Height of oil column $h = \SI{300}{\meter} = \SI{300}{\meter}$.
    $\Delta \rho = \rho_w - \rho_o = 1050 - 850 = \SI{200}{\kilo\gram\per\cubic\meter}$.
    $\Delta P = \Delta \rho \cdot g \cdot h = 200 \times 9.81 \times 300 = \SI{588600}{\pascal} \approx \SI{5.89}{\mega\pascal} (\approx \SI{59}{\bar})$.
    This is the pressure the seal must withstand at the OWC level. The capillary entry pressure of the Cap Shale must exceed \SI{5.9}{\mega\pascal}.
\end{enumerate}
\end{corrige}

\begin{attention}[GCE Exam Traps: Petroleum System]
\begin{itemize}
    \item \textbf{Confusing Source and Reservoir:} "The sandstone generated the oil." \textbf{WRONG.} Shale generates; sandstone stores.
    \item \textbf{Ignoring Timing:} Describing a perfect trap but forgetting to check if it existed \emph{before} hydrocarbon expulsion. "The trap formed 10 Ma ago; charge occurred 50 Ma ago." $\rightarrow$ \textbf{No accumulation.}
    \item \textbf{Fluid Contacts not Horizontal:} Drawing GOC/OWC parallel to bedding. \textbf{WRONG.} They are horizontal (gravity).
    \item \textbf{Porosity = Permeability:} Stating "high porosity means good flow." \textbf{WRONG.} Need connectivity (permeability). Shale has high porosity, near-zero permeability.
    \item \textbf{Cameroon Specifics:} Confusing Rio-del-Rey (NW, Niger Delta affinity) with Douala Basin (SW, Transform margin). Know the \emph{Aptian Salt} is the key seal in Douala/Kribi-Campo; \emph{overpressured shales} key in Rio-del-Rey.
    \item \textbf{Maturity Calculation:} Forgetting to add surface temperature ($T_0$). Formula: $T = T_0 + (\text{gradient} \times \text{depth})$.
    \item \textbf{Gas vs Oil Window:} Stating gas forms at lower temperature than oil. \textbf{WRONG.} Gas window is hotter/deeper.
\end{itemize}
\end{attention}

\courssub{Summary: The Complete Petroleum System}

\begin{aretenir}[Chapter 2 Essentials for GCE]
\begin{enumerate}
    \item \textbf{Origin:} Organic theory. Marine plankton $\rightarrow$ Kerogen (Type I, II, III) in Source Rock (Shale). TOC \textgreater 0.5\%.
    \item \textbf{Maturation:} Temperature-driven. Oil Window \SI{60--120}{\ensuremath{{}^{\circ}\mathrm{C}}} ($R_o$ 0.5--1.0\%); Gas Window \textgreater \SI{120}{\ensuremath{{}^{\circ}\mathrm{C}}} ($R_o$ \textgreater 1.3\%). Geothermal gradient controls depth.
    \item \textbf{Migration:} Primary (expulsion from source, inefficient) $\rightarrow$ Secondary (buoyant flow up carrier beds).
    \item \textbf{Reservoir:} Porosity (storage) + Permeability (flow). Sandstone (primary), Limestone (secondary/fracture). $\phi>15\%$, $k>100$ mD good.
    \item \textbf{Seal:} Shale, Evaporites (Salt). Capillary entry pressure $>$ Buoyancy pressure ($\Delta\rho g h$).
    \item \textbf{Traps:} Structural (Anticline, Fault, Salt Dome) + Stratigraphic (Pinch-out, Unconformity, Reef). Combination traps common.
    \item \textbf{Fluid Contacts:} GOC and OWC are HORIZONTAL. Gas-Oil-Water from top to bottom.
    \item \textbf{Cameroon:} Rio-del-Rey (Gas/Oil, Tertiary/Cretaceous) \& Douala/Kribi-Campo (Oil/Gas, Cretaceous Salt). SONARA refinery (Limbe).
    \item \textbf{Exam Skills:} Read cross-sections, calculate maturity temperature, identify trap types, explain timing, calculate buoyancy pressure.
\end{enumerate}
\end{aretenir}

This section has equipped you with the conceptual framework of the Petroleum System. In the next section, we will move from the static geological model to the dynamic industrial processes: \emph{Drilling, Production, and the Integration of Fossil Fuels}.

\coursec{Drilling, Production and Integration of Fossil Fuels}

In the previous section we saw how a complete \cle{petroleum system} — source rock, maturation, migration, reservoir, seal and trap — concentrates hydrocarbons in the subsurface. But an oil accumulation locked two kilometres beneath the Rio del Rey basin is of no economic value until someone finds it, reaches it, brings it to the surface and transforms it into usable products. This section follows that journey: drilling the well, completing and logging it, recovering the oil in successive stages, refining the crude, processing natural gas, and finally integrating everything we have learned about fossil fuels in the Cameroonian context.

\courssub{Oil Well Drilling: The Rotary Method}

\textbf{From observation to principle.} If you watch water-well drillers at work in Douala or Yaoundé, you see the same basic idea used for oil wells: a rotating bit grinds the rock, and a fluid is pumped down to flush out the fragments. Oil-well drilling is simply this idea pushed to great depths and pressures.

\begin{definition}[Rotary drilling]
\cle{Rotary drilling} is the standard method of drilling oil and gas wells. A \cle{drill bit} (usually a tricone roller bit for softer formations or a diamond-studded PDC bit for hard rock) is screwed to the bottom of a \cle{drill string} — a column of hollow steel pipes (\emph{drill pipes}) with heavier \emph{drill collars} near the bit to provide weight on bit (WOB) without buckling the drill pipes. The whole string is rotated from the surface either by a \emph{rotary table} (which turns the kelly and the string) or, on modern rigs, by a \emph{top drive} (a motor that travels up and down the derrick and rotates the string directly). The rig tower, called the \cle{derrick} (or mast), supports the hoisting system (crown block, travelling block, hook and drawworks) that raises and lowers the drill string when pipes must be added (\emph{making a connection}) or the bit changed (\emph{tripping}).
\end{definition}

\begin{popfigure}
\begin{center}
\begin{tikzpicture}[scale=0.85, every node/.style={font=\small}]
% ground
\draw[thick, popdark] (-4,0) -- (4,0);
\fill[popblueL] (-4,-0.15) rectangle (4,0);
% derrick
\draw[very thick, poporange] (-1.4,0) -- (0,6.2) -- (1.4,0);
\draw[poporange] (-1.05,1.5) -- (1.05,1.5);
\draw[poporange] (-0.7,3.0) -- (0.7,3.0);
\draw[poporange] (-0.35,4.6) -- (0.35,4.6);
\draw[poporange] (-1.4,0) -- (1.05,1.5) (-1.05,1.5) -- (0.7,3.0) (-0.7,3.0) -- (0.35,4.6);
\node[popdark, align=center] at (2.9,5.6) {Derrick (mast)};
\draw[->, popmuted] (2.2,5.5) -- (0.5,5.4);
% crown block and hook
\fill[popdark] (-0.25,5.9) rectangle (0.25,6.3);
\node[popdark] at (-2.6,6.1) {Crown block};
\draw[->, popmuted] (-1.7,6.05) -- (-0.3,6.05);
% travelling block / hook
\fill[popdark] (-0.3,4.2) rectangle (0.3,4.5);
\node[popdark] at (-2.6,4.35) {Travelling block};
\draw[->, popmuted] (-1.7,4.35) -- (-0.35,4.35);
% drill string
\draw[very thick, popblue] (0,4.2) -- (0,-4.6);
\node[popblue, align=center] at (2.9,2.6) {Drill string\\(drill pipes + collars)};
\draw[->, popmuted] (2.2,2.5) -- (0.1,2.3);
% rotary table
\fill[popgold] (-0.9,-0.05) rectangle (0.9,0.35);
\node[popdark] at (-2.9,0.35) {Rotary table};
\draw[->, popmuted] (-2.0,0.25) -- (-0.95,0.15);
% kelly
\draw[very thick, popblue] (0,4.2) -- (0,4.8);
\node[popblue, font=\footnotesize] at (1.2,4.5) {Kelly};
% bit
\draw[very thick, popdark] (0,-4.6) -- (-0.28,-5.15) (0,-4.6) -- (0.28,-5.15) (-0.28,-5.15) -- (0.28,-5.15);
\node[popdark] at (2.5,-4.9) {Drill bit};
\draw[->, popmuted] (1.8,-4.85) -- (0.3,-4.95);
% borehole walls
\draw[popmuted] (-0.55,0) -- (-0.55,-4.6) (0.55,0) -- (0.55,-4.6);
% mud circulation arrows
\draw[->, very thick, popgreen] (0.18,0.6) -- (0.18,-3.6);
\draw[->, very thick, popgreen] (-0.42,-3.6) -- (-0.42,-0.4);
\node[popgreen, align=center] at (-3.0,-2.6) {Mud returns\\with cuttings};
\draw[->, popmuted] (-1.9,-2.5) -- (-0.5,-2.2);
\node[popgreen, align=center] at (2.9,-1.4) {Mud pumped\\down the string};
\draw[->, popmuted] (2.1,-1.3) -- (0.25,-1.6);
% mud pump and pits
\draw[thick, popgreen] (3.4,0) rectangle (4.4,0.8);
\node[popdark] at (3.9,0.4) {Mud pump};
\draw[thick, popgreen] (3.4,0.4) -- (0.6,0.4);
\draw[thick, popgreen] (-0.6,0.4) -- (-2.5,0.4);
\draw[thick, popgreen] (-2.5,0.4) -- (-2.5,-0.5);
\node[popgreen, font=\footnotesize] at (-2.5,-0.8) {Shale shaker};
% strata
\draw[popmuted, dashed] (-4,-2.2) -- (-0.55,-2.2) (0.55,-2.2) -- (4,-2.2);
\draw[popmuted, dashed] (-4,-3.6) -- (-0.55,-3.6) (0.55,-3.6) -- (4,-3.6);
\node[popmuted, font=\footnotesize] at (-3.1,-1.1) {sedimentary strata};
\end{tikzpicture}
\end{center}
\legende{Schematic of a rotary drilling rig. Mud is pumped down the hollow drill string, exits through the bit nozzles, and rises back up the annulus carrying rock cuttings to the shale shaker.}
\end{popfigure}

\begin{definition}[Drilling mud and its four functions]
\cle{Drilling mud} (drilling fluid) is a specially formulated fluid (water-based or oil-based, usually weighted with baryte, \ce{BaSO4}) circulated continuously during drilling. Its four essential functions are:
\begin{enumerate}
\item \textbf{Cooling and lubricating} the drill bit, which would otherwise overheat and wear out rapidly;
\item \textbf{Carrying cuttings} (rock fragments) up the annulus to the surface, where they are separated on a \emph{shale shaker} and examined by the well-site geologist;
\item \textbf{Controlling subsurface pressure}: the hydrostatic pressure of the mud column balances formation pressures and prevents a dangerous \emph{blowout} (uncontrolled eruption of formation fluids);
\item \textbf{Plastering the borehole wall} with an impermeable \emph{mud cake} (filter cake), which stabilises the hole and limits fluid loss into permeable formations.
\end{enumerate}
\end{definition}

\begin{savaistu}[Did you know?]
The cuttings brought up by the mud are the geologist's first direct samples of the subsurface. By describing them under the microscope as drilling progresses, the well-site geologist builds the \emph{mud log} (lithology, gas shows, rate of penetration) and can announce in real time when the bit enters a promising sandstone — long before any sophisticated measurement is made.
\end{savaistu}

\courssub{Well Completion, Casing and Well Logging}

\textbf{Observation.} A freshly drilled hole through soft sediments would collapse or let fluids from different layers mix. The hole must therefore be stabilised and then "interrogated" to find out exactly what it has crossed.

\begin{definition}[Casing and completion]
\cle{Casing} consists of steel tubes lowered into the hole and cemented in place; successive strings of decreasing diameter are set as the well deepens:
\begin{itemize}
\item \emph{Conductor pipe} (largest, sets the cellar and supports the wellhead);
\item \emph{Surface casing} (protects freshwater aquifers and provides a foundation for the blowout preventer);
\item \emph{Intermediate casing} (isolates troublesome zones: lost circulation, high pressure, salt);
\item \emph{Production casing} (runs through the reservoir to total depth).
\end{itemize}
\cle{Well completion} is the set of operations that turn a drilled hole into a producing well: the production casing is \emph{perforated} opposite the reservoir using shaped charges, production tubing is installed inside the casing, a \emph{packer} seals the annulus, and a \cle{Christmas tree} (assembly of valves, chokes and gauges) is placed at the wellhead to control flow.
\end{definition}

\begin{definition}[Well logging]
\cle{Well logging} (wireline logging) is the recording of physical properties of the rocks traversed, using instruments (sondes) lowered on a cable after drilling (or in real time via \emph{logging while drilling}, LWD). Two logs are fundamental at this level:
\begin{itemize}
\item The \cle{gamma-ray log} (GR) measures natural radioactivity. \textbf{Shales} are rich in radioactive potassium, uranium and thorium, so they give \emph{high} readings; clean \textbf{sandstones} and limestones give \emph{low} readings.
\item \cle{Electric logs} measure electrical behaviour. The \emph{spontaneous potential} (SP) log detects permeability contrasts; the \emph{resistivity log} (induction or laterolog) measures formation resistivity. Resistivity is high in hydrocarbon-bearing or tight rocks, low in salty-water-bearing or shaly rocks.
\end{itemize}
\end{definition}

\begin{methode}[Interpreting a gamma-ray log to spot a reservoir sandstone]
\begin{enumerate}
\item Draw a vertical \emph{shale baseline} through the highest (rightmost) gamma-ray values — these correspond to shales.
\item Look for intervals where the curve deflects strongly to the \textbf{left} (low radioactivity): these are clean, potentially porous sandstones or limestones.
\item Confirm with the resistivity log: if the clean interval also shows \emph{high resistivity}, it probably contains hydrocarbons rather than salt water.
\item Identify the overlying high-gamma-ray shale: it is a potential \textbf{cap rock} (seal).
\item Conclude: a low-GR sandstone with high resistivity, sealed by shale above, is a candidate reservoir to be tested.
\end{enumerate}
\end{methode}

\begin{attention}[Common mistake]
Do not claim that a low gamma-ray reading \emph{proves} the presence of oil. A low reading only indicates a clean (non-shaly) rock; the rock may be water-filled or even tight (low porosity). Hydrocarbons are suspected only when low gamma ray is combined with high resistivity and favourable porosity (from density/neutron/sonic logs).
\end{attention}

\courssub{Recovery of Oil: Primary, Secondary and Enhanced}

\textbf{Observation.} When the first wells were drilled in some oil provinces, oil gushed out by itself; later the flow weakened and pumps had to be installed. This everyday observation hides the three successive stages of oil recovery.

\begin{definition}[The three stages of recovery]
\begin{itemize}
\item \cle{Primary recovery}: oil flows naturally, driven by the reservoir's own energy — gas under pressure dissolved in the oil (\emph{solution gas drive}) or overlying it (\emph{gas cap drive}), expansion of underlying water (\emph{water drive}), and gravity drainage. It typically recovers only about 10--20\,\% of the oil originally in place (OOIP).
\item \cle{Secondary recovery}: when natural pressure declines, energy is supplied from outside by \textbf{injecting water} (waterflooding) or \textbf{gas} through injection wells to maintain reservoir pressure and sweep the oil towards producing wells. Recovery rises to roughly 30--40\,\% of OOIP.
\item \cle{Enhanced oil recovery} (EOR, tertiary recovery): more advanced techniques — injecting steam to reduce the viscosity of heavy oil (\emph{thermal EOR}), \ce{CO2} or hydrocarbon gases to achieve miscibility (\emph{miscible gas injection}), or chemicals such as polymers (to increase water viscosity) and surfactants (to lower interfacial tension) — can push recovery beyond 50\,\% of OOIP.
\end{itemize}
\end{definition}

\begin{attention}[Common mistake — primary vs secondary recovery]
A very frequent examination error is to say that secondary recovery means "pumping oil out with a pump". Using a surface pump (the familiar "nodding donkey" or beam pump) is still \textbf{primary} recovery, because no energy is added to the \emph{reservoir}. Recovery becomes \textbf{secondary} only when a fluid (water or gas) is \emph{injected into the reservoir} to maintain pressure and sweep the oil towards the producers.
\end{attention}

\courssub{Refining: Fractional Distillation of Crude Oil}

\textbf{Observation.} Crude oil straight from the well is a thick, dark mixture that cannot be used directly in an engine. Yet from it we obtain cooking gas, petrol, kerosene, diesel and road tar. The separation exploits a simple physical fact: the different hydrocarbons in crude oil have different \cle{boiling ranges}.

\begin{definition}[Fractional distillation]
\cle{Fractional distillation} is the primary refining process in which crude oil is heated (to about $350$--$400\ ^{\circ}\mathrm{C}$) and fed into a tall \cle{fractionating column}. Vapours rise and cool; each fraction condenses at the level where the column temperature matches its boiling range. \textbf{Light, small molecules} (low boiling point) are collected near the \textbf{top}; \textbf{heavy, large molecules} near the \textbf{bottom}. The residue (bitumen) leaves the base. Further processing (cracking, reforming, treating) converts heavy fractions into lighter, more valuable products and removes sulphur.
\end{definition}

\begin{popfigure}
\begin{center}
\begin{tikzpicture}[scale=0.95]
% Column
\draw[thick, fill=popblueL] (0,0) rectangle (2.2,8);
% Temperature gradient arrow
\draw[-latex, thick, poporange] (-1.1,0.3) -- (-1.1,7.7) node[midway, left, align=center, text width=2.2cm] {decreasing temperature};
% Trays
\foreach \y in {1.2,2.5,3.8,5.1,6.4}{
  \draw[popdark] (0,\y) -- (2.2,\y);
}
% Fraction outlets
\foreach \y/\name/\range in {
  7.0/{Refinery gas (LPG)}/{below $25\ ^{\circ}\mathrm{C}$},
  5.6/{Petrol (gasoline)}/{$25$--$150\ ^{\circ}\mathrm{C}$},
  4.3/{Kerosene (jet fuel)}/{$150$--$250\ ^{\circ}\mathrm{C}$},
  3.0/{Diesel (gas oil)}/{$250$--$350\ ^{\circ}\mathrm{C}$},
  1.7/{Lubricating oil, waxes}/{$300$--$370\ ^{\circ}\mathrm{C}$}
}{
  \draw[-latex, thick, popblue] (2.2,\y) -- (3.4,\y);
  \node[right, align=left, text width=5.5cm] at (3.4,\y) {\textbf{\name}\\ \range};
}
% Crude inlet
\draw[-latex, very thick, popdark] (-2.2,0.6) -- (0,0.6) node[midway, above, align=center, text width=2.4cm] {heated crude oil ($\sim 400\ ^{\circ}\mathrm{C}$)};
% Residue outlet
\draw[-latex, very thick, popdark] (1.1,0) -- (1.1,-0.9) node[below] {bitumen (residue)};
\end{tikzpicture}
\end{center}
\legende{Fractionating column: crude oil vapours separate according to boiling range; light products leave at the cool top, heavy products and bitumen at the hot base.}
\end{popfigure}

\begin{propriete}[Crude oil quality: API gravity and sulphur content]
The commercial quality of a crude oil is defined mainly by two parameters:
\begin{itemize}
\item \textbf{API gravity} (American Petroleum Institute): $\mathrm{API} = \dfrac{141.5}{\text{specific gravity at }60^{\circ}\mathrm{F}} - 131.5$. The \emph{higher} the API, the \emph{lighter} the crude and the greater its yield of valuable light products (petrol). Light crude: API $> 31$; medium: $22$--$31$; heavy: API $< 22$.
\item \textbf{Sulphur content}: \emph{sweet} crudes (low sulphur, $< 0.5\,\%$) are more valuable than \emph{sour} crudes (high sulphur), because sulphur must be removed during refining (hydrodesulphurisation) to meet environmental specifications.
\end{itemize}
A light, sweet crude therefore commands the highest price. (Knowledge of the formula is useful; its derivation is not examinable.)
\end{propriete}

\courssub{Natural Gas: Processing and Uses}

Natural gas, mostly \cle{methane} (\ce{CH4}), may occur with oil (associated gas) or alone (non-associated gas). Before use it is processed: removal of water (dehydration), \ce{CO2} and hydrogen sulphide (\ce{H2S}) (acid gas removal), and heavier hydrocarbons (which are recovered as \emph{natural gas liquids} — NGLs — and LPG). Its main uses are:
\begin{itemize}
\item \textbf{Electricity generation} in gas-fired thermal power plants (combined cycle for high efficiency);
\item \textbf{Fertiliser manufacture}: methane is the feedstock for ammonia (\ce{NH3}) via steam reforming, and hence urea;
\item \textbf{Domestic and industrial fuel} (cooking gas, industrial heating);
\item \textbf{Petrochemical feedstock} for methanol, ethylene, etc.
\end{itemize}

\begin{savaistu}[Did you know? — the Logbaba gas project]
The \textbf{Logbaba} gas field, on the outskirts of Douala, was Cameroon's first natural gas development (operated by Victoria Oil \& Gas / GEC). Gas from Logbaba is treated and distributed by pipeline to industrial customers in the Douala area, where it fuels boilers and electricity generators, reducing dependence on imported fuel oil. It is a concrete illustration that a modest gas accumulation, close to an industrial market, can be highly economic even without export. The national oil company, \textbf{SNH} (Société Nationale des Hydrocarbures), plays a key role in managing the state's hydrocarbon interests.
\end{savaistu}

\courssub{Worked Example and Integration 4}

\begin{exoresolu}[From crude to products]
Statement: A crude oil has a specific gravity of $0.85$. (a) Calculate its API gravity and classify the crude. (b) A refinery receives this crude; state, with justification, whether you expect a high or low yield of petrol. (c) Name the product collected at the base of the fractionating column and give one use.
\tcblower
Solution:
\begin{enumerate}
\item $\mathrm{API} = \dfrac{141.5}{0.85} - 131.5 \approx 166.47 - 131.5 = 34.97 \approx 35.0$. Since API $> 31$, this is a \textbf{light crude}.
\item A light crude contains a large proportion of small molecules with low boiling points, so it yields a \textbf{high proportion of petrol} (and other light fractions) — which is why it is commercially valuable.
\item The base product is \textbf{bitumen} (asphalt), used for \textbf{surfacing roads} (and roofing).
\end{enumerate}
\end{exoresolu}

\begin{exercice}[Integration 4 — Cameroon's fossil-fuel system]
The Rio del Rey basin (offshore South-West Region) contains Tertiary deltaic sediments.
\begin{enumerate}
\item Explain why thick accumulations of buried plant matter in a delta can produce both coal seams and, at greater depth and temperature, a petroleum source rock. Refer to the stages of coalification and to the "oil window".
\item Describe the complete petroleum system you would expect in this deltaic basin (source, reservoir, seal, trap).
\item A well is drilled into a sandstone at $2500\ \mathrm{m}$. The gamma-ray log shows low values in the sandstone and high values above it; resistivity is high in the sandstone. Interpret these logs.
\item The field's production declines after several years. Propose a secondary recovery method and explain its principle.
\item Give two reasons why developing the Logbaba-type gas resource is attractive for the Cameroonian economy.
\end{enumerate}
\end{exercice}

\begin{corrige}[Integration 4]
\begin{enumerate}
\item In a delta, plant debris accumulates in swamps and is rapidly buried by sediment. Burial first drives \textbf{coalification}: peat $\rightarrow$ lignite $\rightarrow$ bituminous coal $\rightarrow$ anthracite as temperature and pressure rise. Where organic matter (especially planktonic/algal matter in marine shales) is buried deeper, into the \textbf{oil window} (roughly $60$--$120\ ^{\circ}\mathrm{C}$), kerogen cracks into liquid petroleum; at still higher temperatures ($>150\ ^{\circ}\mathrm{C}$) only gas forms (gas window). Thus the same deltaic burial history yields coal at shallow levels and oil/gas source rocks deeper down.
\item \textbf{Source rock}: organic-rich marine/deltaic shale matured in the oil window. \textbf{Reservoir}: porous deltaic sandstones (channel sands, distributary mouth bars). \textbf{Seal}: impermeable shales interbedded with or overlying the sands. \textbf{Trap}: rollover anticlines and fault traps associated with deltaic growth faults, plus stratigraphic pinch-outs of sand bodies.
\item Low gamma ray $=$ clean sandstone (potential reservoir); high gamma ray above $=$ shale (cap rock/seal); high resistivity in the sandstone suggests hydrocarbon saturation rather than salt water. Conclusion: a sealed, hydrocarbon-bearing reservoir — worth testing.
\item Use \textbf{waterflooding}: inject water through injection wells to maintain reservoir pressure and sweep the remaining oil towards the production wells (secondary recovery).
\item Any two: it supplies cheaper, cleaner fuel for Douala's industries and power generation; it reduces imports of refined products, saving foreign currency; it provides feedstock for fertiliser production (ammonia/urea); it creates jobs and infrastructure; it monetises a resource that might otherwise be flared.
\end{enumerate}
\end{corrige}

\begin{aretenir}[Key points]
\begin{itemize}
\item Rotary drilling: a rotating bit on a hollow drill string, supported by the derrick; drilling mud cools the bit, lifts cuttings, controls pressure and cakes the borehole wall.
\item Wells are cased, cemented, perforated and completed; gamma-ray and electric logs identify clean reservoir sandstones (low GR, high resistivity) and sealing shales (high GR).
\item Recovery: primary (natural reservoir energy), secondary (water/gas injection), enhanced (steam, \ce{CO2}, chemicals). Pumping alone is still primary recovery.
\item Fractional distillation separates crude by boiling range: LPG, petrol, kerosene, diesel, lubricants, bitumen. Quality is set by API gravity (light vs heavy) and sulphur (sweet vs sour).
\item Processed natural gas serves for electricity, fertiliser and domestic/industrial fuel — as at Logbaba near Douala.
\end{itemize}
\end{aretenir}

\coursec{Applications of Geological Materials and Economic Uses of Rocks}

Having examined the origin, exploration and production of fossil fuels — coal, petroleum and natural gas — we now turn to the solid geological materials that underpin almost every aspect of modern infrastructure and industry. From the aggregate in the concrete of a classroom block in Yaoundé to the limestone that feeds the Figuil cement plant, from the volcanic scoria that lightens horticultural substrates to the iron ore of Mbalam destined for export, the economic value of a rock is dictated by its \cle{mineralogy}, \cle{texture}, \cle{physical properties} and \cle{geological setting}. This section surveys the principal categories of geological materials used in Cameroon and worldwide, linking each use to the rock properties that make it fit for purpose.

\courssub{Building Stones: Properties and Selection}

A \cle{building stone} is any massive rock quarried and cut into blocks or slabs for structural or decorative use. The choice of rock is governed by a combination of \cle{strength}, \cle{durability}, \cle{workability} and \cle{appearance}.

\begin{propriete}[Qualities of a Good Building Stone]
A rock suitable for dimension stone must satisfy the following criteria:
\begin{itemize}
    \item \textbf{High compressive strength} — typically $> 100\ \mathrm{MPa}$ for load-bearing masonry.
    \item \textbf{Low water absorption} — $< 1\%$ by weight to resist freeze-thaw and salt crystallisation.
    \item \textbf{Durability} — resistance to chemical weathering (acid rain, pollution) and physical disintegration.
    \item \textbf{Workability} — can be cut, dressed and finished with reasonable effort; uniform texture without deleterious joints or veins.
    \item \textbf{Aesthetic appeal} — colour, pattern and ability to take a polish for decorative use.
    \item \textbf{Availability in large, unfractured blocks} — controlled by joint spacing and quarry geometry.
\end{itemize}
\end{propriete}

\begin{center}
\begin{tabularx}{\linewidth}{|l|X|X|X|}\hline
\rowcolor{popblueL}
\textbf{Rock type} & \textbf{Typical strength / durability} & \textbf{Workability} & \textbf{Common uses in Cameroon} \\ \hline
Granite & Very high strength, excellent durability & Hard to cut; requires diamond tools & Kerbstones, monuments, polished cladding (e.g. Yaoundé city hall) \\ \hline
Gneiss & High strength, good durability; foliation may weaken & Moderate; splits along foliation & Foundations, walling stone in highland areas (Dschang, Bafoussam) \\ \hline
Basalt & Very high strength, excellent durability & Very hard; fine grain takes polish well & Road setts, aggregate, occasional dimension stone \\ \hline
Sandstone & Variable (50--150 MPa); durability depends on cement & Easy to cut and dress & Building blocks in sedimentary basins (Garoua, Maroua) \\ \hline
Limestone & Moderate strength; susceptible to acid attack & Very easy to cut and carve & Traditional construction in North (Rey Bouba), cement raw material \\ \hline
Marble & High strength, takes high polish & Moderate; cuts well & Decorative cladding, flooring, sculpture \\ \hline
\end{tabularx}
\end{center}

\emph{Worked example:} A contractor must choose between a fine-grained sandstone (compressive strength $85\ \mathrm{MPa}$, water absorption $3.5\%$) and a medium-grained granite (compressive strength $180\ \mathrm{MPa}$, water absorption $0.4\%$) for the exterior cladding of a five-storey building in Douala. The granite is preferred despite higher quarrying cost because its low absorption and high strength withstand the humid, salt-laden coastal climate far better than the sandstone, which would suffer rapid salt weathering and strength loss.

\courssub{Aggregates, Sand and Gravel for Concrete and Road Construction}

\cle{Aggregates} are granular materials (sand, gravel, crushed stone) that form $60$--$75\%$ of concrete volume and the bulk of road base and asphalt. Their quality directly controls the strength, durability and workability of the final product.

\begin{methode}[Choosing an Aggregate for Concrete]
\begin{enumerate}
    \item \textbf{Hardness and abrasion resistance} — Test with Los Angeles abrasion value (LAV); aim for LAV $< 30\%$ for high-strength concrete. Basalt, quartzite and dense granite are excellent; soft sandstone or weathered rock are poor.
    \item \textbf{Particle shape} — Cubical or rounded particles give better workability and higher strength than flaky or elongated particles (flakiness index $< 15\%$, elongation index $< 20\%$). Crushed rock tends to be angular; alluvial gravel is rounded.
    \item \textbf{Cleanliness} — Free from clay lumps, organic matter, mica, sulfides and soluble salts. Clay coats particles and weakens the cement paste bond; organics retard setting. Wash if necessary.
    \item \textbf{Grading} — Well-graded (continuous particle size distribution) minimises voids and cement paste demand. Follow standard grading envelopes (e.g. BS 882, EN 12620).
    \item \textbf{Durability (soundness)} — Sodium or magnesium sulfate soundness test loss $< 12\%$ after 5 cycles. Critical in climates with wet-dry or freeze-thaw cycles.
    \item \textbf{Alkali-silica reactivity (ASR)} — Avoid reactive silica (opal, chalcedony, volcanic glass) in high-alkali cement; use low-alkali cement or pozzolanic addition if reactive aggregate is unavoidable.
\end{enumerate}
\end{methode}

In Cameroon, the alluvial sands of the Wouri, Sanaga and Benue rivers supply much of the construction sand, while crushed basalt and granite from quarries along the Cameroon Volcanic Line (e.g. Buea, Dschang, Ngaoundéré) provide coarse aggregate. The volcanic scoria (pozzolana) from recent cones is also used as lightweight aggregate and supplementary cementitious material (see below).

\courssub{Cement Raw Materials: The Figuil Cement Works}

Portland cement is manufactured by heating a precise mixture of \cle{limestone} (calcium carbonate, $\ce{CaCO3}$) and \cle{clay} (aluminosilicates) to $\sim 1450^\circ\mathrm{C}$ in a rotary kiln. The chemistry requires:
\begin{itemize}
    \item $\ce{CaO}$ (from limestone) $\approx 60$--$67\%$
    \item $\ce{SiO2}$ (from clay/sand) $\approx 19$--$24\%$
    \item $\ce{Al2O3}$ (from clay) $\approx 4$--$7\%$
    \item $\ce{Fe2O3}$ (from clay/iron ore) $\approx 2$--$4\%$
\end{itemize}

\begin{experience}[Figuil Cement Works — North Region, Cameroon]
The Figuil plant, commissioned in the 1970s and expanded since, exploits a Cretaceous limestone sequence in the Benue Trough (Figuil Formation) and overlying clay/shale. The limestone is relatively pure ($> 90\%\ \ce{CaCO3}$), while the clay provides the necessary silica, alumina and iron. The quarry operates by bench blasting; raw materials are crushed, blended in a pre-homogenisation pile, then fed to a dry-process kiln. Annual capacity exceeds $500\,000$ tonnes. The plant illustrates the classic \cle{limestone-clay couple} in a sedimentary basin setting — a model replicable wherever similar stratigraphy exists (e.g. Logone Birni, Koutaba).
\end{experience}

\courssub{Industrial Minerals of the Cameroon Volcanic Line}

The Cameroon Volcanic Line (CVL) is a treasure trove of industrial minerals thanks to its varied volcanic and plutonic rocks.

\begin{itemize}
    \item \textbf{Marble} — Metamorphosed limestone in the Pan-African belt (e.g. near Maroua, Bidzar). Used for dimension stone, lime production, and as a flux in metallurgy.
    \item \textbf{Pozzolana (volcanic scoria)} — Vesicular, glassy basaltic ejecta from recent strombolian cones (e.g. Mount Cameroon flanks, Manengouba, Bambouto). When finely ground, it reacts with portlandite ($\ce{Ca(OH)2}$) to form additional C-S-H gel, improving concrete durability and reducing cement demand. Widely used in Cameroon as a \cle{supplementary cementitious material (SCM)} and lightweight aggregate.
    \item \textbf{Clay for bricks and pottery} — Residual kaolinitic clays from weathering of granites/gneisses (e.g. Foumban, Baham) and alluvial clays in floodplains. Fired bricks are a major building material in rural and peri-urban areas.
    \item \textbf{Salt} — Evaporitic deposits in the Lake Chad basin (e.g. Boudouri) and coastal lagoons. Produced by solar evaporation of brine.
\end{itemize}

\courssub{Metallic Ore Minerals: Definitions and Cameroonian Examples}

While the GCE syllabus focuses on non-metallic materials, understanding the terminology of metallic deposits is examinable and essential for geological literacy.

\begin{definition}[Ore, Gangue, Grade and Cut-off Grade]
\begin{itemize}
    \item \textbf{Ore} — A naturally occurring solid material from which a valuable mineral or metal can be extracted \emph{profitably}. Profitability depends on market price, extraction cost and technology.
    \item \textbf{Gangue} — The unwanted mineral(s) associated with the ore mineral in a deposit (e.g. quartz, calcite, pyrite). Must be separated during processing.
    \item \textbf{Grade} — The concentration of the valuable element in the ore, usually expressed as weight percent (\% or g/t for precious metals). Example: $40\%\ \ce{Fe}$ in hematite ore.
    \item \textbf{Cut-off grade} — The minimum grade below which mining is uneconomic. It varies with metal price, mining method (open pit vs underground), processing recovery and costs. Material above cut-off is \emph{ore}; below is \emph{waste}.
\end{itemize}
\end{definition}

\begin{attention}[Common Mistake: Calling Every Extracted Rock an 'Ore']
\textbf{Do not} call limestone quarried for cement, sand for concrete, or granite for building stone an \emph{ore}. These are \cle{industrial minerals} or \cle{construction materials}. The term \emph{ore} is reserved for metallic (or metalliferous) deposits where a metal is extracted. A rock is an ore \emph{only} if the target commodity is a metal and the grade exceeds the cut-off grade. In exams, misusing 'ore' for limestone or aggregate loses marks.
\end{attention}

\begin{savaistu}[Cameroon's Major Metallic Deposits (GCE-useful Extra)]
\begin{itemize}
    \item \textbf{Mbalam (South-East)} — World-class \cle{iron ore} deposit (hematite/goethite) in Archean Banded Iron Formation (BIF) of the Congo Craton. Estimated resources $> 2\ \mathrm{Gt}$ at $\sim 60\%\ \ce{Fe}$. Planned for rail export via Kribi deep-sea port.
    \item \textbf{Minim-Martap (Adamaoua)} — Large \cle{bauxite} deposit (gibbsite/boehmite) formed by lateritic weathering of Cretaceous syenite/nepheline syenite. Resources $> 1\ \mathrm{Gt}$; low silica, suitable for alumina refining.
    \item \textbf{Nkamouna (East)} — \cle{Cobalt-Nickel-Manganese} laterite deposit over ultramafic rocks; strategic for battery metals.
    \item \textbf{Akonolinga (Centre)} — \cle{Tin-tantalum-niobium} mineralisation in rare-metal granites (pegmatites/greisen).
\end{itemize}
\end{savaistu}

\courssub{Dimension Stone vs Crushed Stone: Quarrying Basics}

\begin{itemize}
    \item \textbf{Dimension stone} — Quarried by \cle{wire saw}, \cle{diamond wire} or \cle{controlled blasting} to produce large, intact blocks ($> 1\ \mathrm{m^3}$). Requires massive, unfractured rock masses. High value per tonne; low volume.
    \item \textbf{Crushed stone (aggregate)} — Produced by \cle{drilling and blasting} followed by \cle{crushing} (jaw, cone, impact crushers) and \cle{screening} to specified size fractions. Can exploit more fractured rock. High volume; lower value per tonne.
\end{itemize}

Quarry design considers \cle{bench height}, \cle{blast pattern}, \cle{haul road gradient} and \cle{environmental mitigation} (dust suppression, noise bunds, progressive rehabilitation).

\courssub{Rocks as Sources of Soil Fertility and Groundwater}

The economic value of rocks extends beyond direct extraction; they control two vital resources: \cle{soil fertility} and \cle{groundwater}.

\begin{propriete}[Rock Control on Soil Fertility]
\begin{itemize}
    \item \textbf{Basic/mafic rocks (basalt, gabbro, dolerite)} — Weather to clay-rich, cation-rich soils (smectite, vermiculite) with high CEC (cation exchange capacity), good water retention and natural fertility. The volcanic highlands of the CVL (Bamenda, Dschang, Foumban) support intensive agriculture on such soils.
    \item \textbf{Acidic/felsic rocks (granite, gneiss, sandstone)} — Yield sandy, kaolinitic soils with low CEC, low base saturation, prone to leaching and acidity. Require liming and fertiliser for sustained cropping.
    \item \textbf{Limestone/marble} — Produce calcareous, high-pH soils; may induce micronutrient deficiencies (Fe, Zn) but excellent for legumes.
    \item \textbf{Volcanic ash/scoria} — Rapidly weather to release nutrients; andosols have high phosphate retention but high organic matter.
\end{itemize}
\end{propriete}

\begin{propriete}[Aquifers in Volcanic and Sedimentary Terrains]
\begin{itemize}
    \item \textbf{Volcanic aquifers (CVL)} — \cle{Vesicular basalt}, \cle{scoria} and \cle{fractured lava flows} form high-yield, unconfined aquifers. Springs issue at lithological contacts (e.g. basalt over clay). Recharge is rapid; vulnerability to pollution is high. Examples: Mount Cameroon flanks, Bamenda highlands.
    \item \textbf{Sedimentary aquifers (Benue Trough, Logone Basin, Douala Basin)} — \cle{Sandstone} (e.g. Figuil, Poli, Koutaba formations) and \cle{limestone} (karstic) aquifers. Often confined (artesian) in deeper parts. Transmissivity controlled by primary porosity (sandstone) or secondary karst conduits (limestone).
    \item \textbf{Basement aquifers} — \cle{Weathered zone (saprolite)} and \cle{fractured bedrock} in granites/gneisses. Low to moderate yield ($1$--$10\ \mathrm{m^3/h}$); success depends on fracture density and regolith thickness. Critical for rural water supply in East, South and Centre regions.
\end{itemize}
\end{propriete}

\begin{exemplebox}[Integrated Resource Assessment]
A geological survey team maps a $50\ \mathrm{km^2}$ area in the Adamaoua region underlain by basaltic lava flows (CVL) overlying Cretaceous sandstone. They identify:
\begin{enumerate}
    \item \textbf{Construction aggregate} — Basalt quarries for road metal and concrete (high LAV resistance).
    \item \textbf{Pozzolana} — Scoria cones for cement additive.
    \item \textbf{Groundwater} — Basalt aquifer (spring lines at basalt-sandstone contact) + sandstone aquifer (deeper boreholes).
    \item \textbf{Soil fertility} — Basalt-derived soils for coffee/food crops; sandstone-derived soils for grazing.
    \item \textbf{Dimension stone} — Not viable (basalt too fractured; sandstone too soft).
\end{enumerate}
This multi-resource view is the hallmark of modern economic geology.
\end{exemplebox}

\courssub{Illustrations}

\begin{popfigure}
\centering
\begin{tikzpicture}[node distance=1.2cm and 1.5cm, every node/.style={font=\small}]
% Rock types
\node[draw, rounded corners, fill=popblueL, text=white, minimum width=2.5cm, minimum height=1cm] (granite) {Granite / Gneiss};
\node[draw, rounded corners, fill=popblueL, text=white, minimum width=2.5cm, minimum height=1cm, below=of granite] (basalt) {Basalt / Dolerite};
\node[draw, rounded corners, fill=popblueL, text=white, minimum width=2.5cm, minimum height=1cm, below=of basalt] (sandstone) {Sandstone};
\node[draw, rounded corners, fill=popblueL, text=white, minimum width=2.5cm, minimum height=1cm, below=of sandstone] (limestone) {Limestone / Marble};
\node[draw, rounded corners, fill=popblueL, text=white, minimum width=2.5cm, minimum height=1cm, below=of limestone] (scoria) {Volcanic Scoria (Pozzolana)};
\node[draw, rounded corners, fill=popblueL, text=white, minimum width=2.5cm, minimum height=1cm, below=of scoria] (clay) {Clay / Shale};

% Uses
\node[draw, rounded corners, fill=popgreenL, minimum width=3.5cm, minimum height=0.9cm, right=of granite, align=center] (dim){Dimension Stone\\ (blocks, cladding, monuments)};
\node[draw, rounded corners, fill=popgreenL, minimum width=3.5cm, minimum height=0.9cm, right=of basalt, align=center] (agg){Crushed Aggregate\\ (concrete, road base, rail ballast)};
\node[draw, rounded corners, fill=popgreenL, minimum width=3.5cm, minimum height=0.9cm, right=of sandstone, align=center] (build){Building Blocks\\ (walling, lintels)};
\node[draw, rounded corners, fill=popgreenL, minimum width=3.5cm, minimum height=0.9cm, right=of limestone, align=center] (cement){Cement Raw Material\\ (Figuil, Koutaba)};
\node[draw, rounded corners, fill=popgreenL, minimum width=3.5cm, minimum height=0.9cm, right=of scoria, align=center] (pozz){Pozzolana / SCM\\ (blended cement, lightweight agg.)};
\node[draw, rounded corners, fill=popgreenL, minimum width=3.5cm, minimum height=0.9cm, right=of clay, align=center] (brick){Bricks, Pottery,\\ Ceramics, Tile};

% Arrows
\draw[->, thick, popdark] (granite.east) -- (dim.west);
\draw[->, thick, popdark] (granite.east) -- (agg.west);
\draw[->, thick, popdark] (basalt.east) -- (agg.west);
\draw[->, thick, popdark] (basalt.east) -- (dim.west);
\draw[->, thick, popdark] (sandstone.east) -- (build.west);
\draw[->, thick, popdark] (sandstone.east) -- (agg.west);
\draw[->, thick, popdark] (limestone.east) -- (cement.west);
\draw[->, thick, popdark] (limestone.east) -- (dim.west);
\draw[->, thick, popdark] (scoria.east) -- (pozz.west);
\draw[->, thick, popdark] (scoria.east) -- (agg.west);
\draw[->, thick, popdark] (clay.east) -- (brick.west);
\draw[->, thick, popdark] (clay.east) -- (cement.west);

% Title
\node[above=0.5cm of granite, font=\bfseries\large, text=popdark] {Matching Rock Types to Economic Uses};
\end{tikzpicture}
\legende{Flowchart linking principal rock types in Cameroon to their main economic applications. Arrows show multiple uses per rock.}
\end{popfigure}

\begin{popfigure}
\centering
\begin{tikzpicture}[scale=0.85]
% Cameroon outline (simplified)
\draw[thick, popdark] (0,0) .. controls (1,2) and (3,4) .. (2,6) .. controls (1,7) and (-1,7) .. (-2,6) .. controls (-3,4) and (-1,2) .. (0,0);
% Major features
\draw[popblue, thick] (-1.5,1) .. controls (-0.5,2) .. (0,3); % Sanaga
\draw[popblue, thick] (-1,0.5) .. controls (-0.5,1.5) .. (0,2); % Wouri
\draw[popblue, thick] (0.5,0.5) .. controls (1,1.5) .. (1.5,2.5); % Benue (N)
% Volcanic Line
\draw[poporange, thick, dashed] (-1.8,0.5) -- (-1,2) -- (-0.5,3.5) -- (0,5) -- (0.5,6);
% Locations
\node[circle, fill=popred, minimum size=6pt, inner sep=0] (figuil) at (0.8,1.2) {};
\node[circle, fill=popred, minimum size=6pt, inner sep=0] (mbalam) at (-1.5,4.5) {};
\node[circle, fill=popred, minimum size=6pt, inner sep=0] (minim) at (0.2,3.8) {};
\node[circle, fill=popred, minimum size=6pt, inner sep=0] (pozz1) at (-0.8,2.5) {};
\node[circle, fill=popred, minimum size=6pt, inner sep=0] (pozz2) at (0,4.5) {};
\node[circle, fill=popred, minimum size=6pt, inner sep=0] (pozz3) at (-0.5,5.2) {};
% Labels
\node[right=2pt of figuil, font=\tiny, anchor=west] {Figuil (cement)};
\node[left=2pt of mbalam, font=\tiny, anchor=east] {Mbalam (Fe)};
\node[right=2pt of minim, font=\tiny, anchor=west] {Minim-Martap (Bauxite)};
\node[above=2pt of pozz1, font=\tiny, anchor=south] {Pozzolana quarries};
\node[above=2pt of pozz2, font=\tiny, anchor=south] {Mt Cameroon};
\node[above=2pt of pozz3, font=\tiny, anchor=south] {Bambouto};
% Legend
\begin{scope}[shift={(3.5,4)}]
\draw[poporange, thick, dashed] (0,0.6) -- (0.8,0.6); \node[right, font=\tiny] at (0.8,0.6) {Cameroon Volcanic Line};
\draw[popblue, thick] (0,0.3) -- (0.8,0.3); \node[right, font=\tiny] at (0.8,0.3) {Major rivers};
\fill[popred] (0,0) circle (3pt); \node[right, font=\tiny] at (0.8,0) {Key economic sites};
\end{scope}
\node[below=0.2cm, font=\bfseries\small, text=popdark] at (0,-0.5) {Simplified Sketch Map of Cameroon — Major Economic Geology Sites};
\end{tikzpicture}
\legende{Sketch map locating Figuil (cement limestone/clay), Mbalam (iron ore), Minim-Martap (bauxite) and pozzolana quarries along the Cameroon Volcanic Line.}
\end{popfigure}

\courssub{Worked Exercise: Aggregate Quality Assessment}

\begin{exoresolu}[Aggregate Selection for a Bridge Deck]
\textbf{Statement:} A bridge deck in the Bamenda highlands requires $500\ \mathrm{m^3}$ of C30/37 concrete. Two aggregate sources are available:
\begin{itemize}
    \item Source A: Crushed basalt from a CVL quarry. LAV $= 18\%$, flakiness index $= 12\%$, water absorption $= 0.6\%$, sulfate soundness loss $= 6\%$.
    \item Source B: Alluvial gravel from the Mbam River. LAV $= 35\%$, flakiness index $= 8\%$, water absorption $= 1.8\%$, sulfate soundness loss $= 14\%$.
\end{itemize}
Which source meets the specification for high-durability concrete in a freeze-thaw environment? Justify.

\tcblower
\textbf{Detailed solution:}
\begin{enumerate}
    \item \textbf{Hardness (LAV):} Spec for high-strength/durable concrete typically requires LAV $\le 25\%$ (BS 882). Source A ($18\%$) passes; Source B ($35\%$) fails — particles will degrade during mixing and under traffic.
    \item \textbf{Shape:} Both have flakiness index $< 15\%$ (acceptable). Source B is slightly better (rounded), but shape alone cannot compensate for hardness failure.
    \item \textbf{Water absorption:} Source A ($0.6\%$) is very low — minimal effect on w/c ratio. Source B ($1.8\%$) is acceptable but requires moisture correction.
    \item \textbf{Soundness (freeze-thaw proxy):} Sulfate soundness loss $< 12\%$ is a common threshold for freeze-thaw durability. Source A ($6\%$) passes comfortably; Source B ($14\%$) fails — indicates presence of weak, porous or clay-coated particles that will disintegrate in wet-freeze cycles.
\end{enumerate}
\textbf{Conclusion:} Source A (crushed basalt) is the only suitable aggregate. Source B fails on hardness and soundness, critical for a bridge deck exposed to de-icing salts and freeze-thaw. The higher cost of crushing basalt is justified by the 100-year design life requirement.
\end{exoresolu}

\courssub{Summary of Key Points}

\begin{aretenir}[Economic Uses of Rocks — Essentials for GCE]
\begin{itemize}
    \item Building stone selection balances \textbf{strength, durability, workability, appearance} and \textbf{block size}. Granite/gneiss/basalt for high durability; sandstone/limestone for easy working.
    \item Aggregates must be \textbf{hard (low LAV), well-shaped (low flakiness), clean, well-graded, sound} and \textbf{non-reactive}. Basalt and granite from CVL are premium sources in Cameroon.
    \item Cement requires \textbf{limestone + clay} in correct proportions. Figuil exploits a Cretaceous limestone-clay couple in the Benue Trough.
    \item CVL provides \textbf{pozzolana (scoria)} for blended cement, \textbf{marble} for dimension stone/lime, \textbf{clay} for bricks.
    \item \textbf{Ore} = profitable metal-bearing material; \textbf{gangue} = waste mineral; \textbf{cut-off grade} = economic threshold. \emph{Never call limestone, sand or building stone 'ore'.}
    \item Cameroon's flagship metallic deposits: \textbf{Mbalam (Fe)}, \textbf{Minim-Martap (bauxite)}, \textbf{Nkamouna (Co-Ni-Mn)}, \textbf{Akonolinga (Sn-Ta-Nb)}.
    \item \textbf{Dimension stone} needs massive rock and careful extraction; \textbf{crushed stone} uses blasting and crushing.
    \item \textbf{Soil fertility} follows rock chemistry: mafic $\to$ fertile; felsic $\to$ infertile. \textbf{Groundwater} occurs in volcanic (vesicular/fractured), sedimentary (sandstone/karst) and basement (weathered/fractured) aquifers.
\end{itemize}
\end{aretenir}

\begin{exercice}[Self-Assessment Questions]
\begin{enumerate}
    \item List five properties a rock must possess to be a good dimension stone. Explain why a highly foliated gneiss may be unsuitable despite high compressive strength.
    \item A concrete specification requires coarse aggregate with LAV $\le 25\%$, flakiness index $\le 15\%$ and sulfate soundness loss $\le 12\%$. A quarry produces crushed granite with LAV $= 22\%$, flakiness $= 18\%$, soundness loss $= 8\%$. Does it meet the spec? If not, what processing step could improve the flakiness index?
    \item Define \emph{ore}, \emph{gangue}, \emph{grade} and \emph{cut-off grade}. Why is the cut-off grade not a fixed value for a given deposit?
    \item Describe the geological setting and raw materials of the Figuil cement works. What would be the consequence of using a limestone with $15\%\ \ce{MgO}$ (dolomitic) without adjustment?
    \item Compare the aquifer characteristics of (a) vesicular basalt on Mount Cameroon and (b) fractured granite in the East Region. Which is more vulnerable to surface pollution and why?
    \item A lateritic soil developed on basalt in the Bamenda highlands has pH $5.8$, CEC $25\ \mathrm{cmol_+/kg}$ and high exchangeable $\ce{Ca^{2+}}$ and $\ce{Mg^{2+}}$. A soil on granite in the same climate has pH $4.5$, CEC $4\ \mathrm{cmol_+/kg}$ and low bases. Explain the differences in terms of parent rock mineralogy and weathering products.
\end{enumerate}
\end{exercice}

\begin{corrige}[Exercise 1]
\textbf{1.} Strength, durability, low water absorption, workability, aesthetic appeal, large unfractured blocks. Foliation creates planes of weakness; blocks split along foliation, reducing recoverable block size and creating anisotropic strength.
\textbf{2.} Fails on flakiness index ($18\% > 15\%$). Use a \emph{cubical crusher} (vertical shaft impactor) or adjust jaw/cone crusher settings (closed side setting, choke feeding) to produce more equant particles.
\textbf{3.} Ore: profitable metal-bearing material. Gangue: unwanted associated minerals. Grade: concentration of valuable element. Cut-off grade: minimum economic grade; varies with metal price, costs, recovery, mining method.
\textbf{4.} Benue Trough, Cretaceous Figuil Formation (limestone) + overlying clay/shale. High $\ce{MgO}$ causes periclase ($\ce{MgO}$) in clinker, leading to delayed expansion and cracking (unsound cement). Must blend with low-$\ce{MgO}$ limestone or add $\ce{SiO2}$ to form forsterite.
\textbf{5.} (a) Basalt: high primary porosity (vesicles) + fractures; high yield, rapid recharge, unconfined, \textbf{highly vulnerable} to pollution (short travel time, little filtration). (b) Granite: low primary porosity; flow in weathered zone and fractures; lower yield, slower recharge, often semi-confined, \textbf{less vulnerable} (thicker unsaturated zone, more attenuation).
\textbf{6.} Basalt weathers to smectite/vermiculite (high CEC, base-rich) + releases $\ce{Ca, Mg, K}$; granite weathers to kaolinite (low CEC) + quartz; bases leached rapidly in humid climate $\to$ acid, infertile soil.
\end{corrige}

\coursec{Engineering Geology: Dams, Reservoirs, Tunnels and Slope Stability}

In the previous section we saw how rocks and minerals serve as \emph{materials}: aggregates, cement, building stone, road metal. But geology does not only supply the materials of construction --- it also controls the \emph{ground} on which every structure stands. A dam, a tunnel or a highway cut is only as safe as the rock mass beneath and around it. When the geology is ignored, the consequences are measured in collapsed slopes, leaking reservoirs and lost lives. This is the domain of \cle{engineering geology}: the application of geological knowledge to the design, construction and maintenance of civil engineering works.

\courssub{The Role of Geology in Civil Engineering}

Every civil engineering structure interacts with the ground in two directions: the structure loads the ground (a dam pushes on its foundation), and the ground acts on the structure (water pressure, swelling clays, unstable slopes). The engineer therefore needs to know, \emph{before} construction:

\begin{itemize}
\item the \textbf{rock types} present and their strength (massive granite behaves very differently from fractured shale);
\item the \textbf{structure} of the rock mass: bedding, joints, faults, and their orientations;
\item the \textbf{hydrogeology}: water table depth, permeability, springs;
\item the \textbf{geological hazards}: landslides, seismicity, subsidence.
\end{itemize}

\begin{methode}[Site investigation before construction]
A proper site investigation proceeds in stages:
\begin{enumerate}
\item \textbf{Desk study}: examination of geological maps, aerial photographs and satellite images; study of previous work in the area.
\item \textbf{Field reconnaissance}: geological mapping of outcrops, measurement of dip and strike of beds and joints, location of springs and old landslides.
\item \textbf{Subsurface investigation}: boreholes and test pits to determine rock type, depth of weathering, and the position of the water table; recovery of core samples.
\item \textbf{Laboratory testing}: measurement of uniaxial compressive strength, shear strength, permeability and Atterberg limits of soils.
\item \textbf{Geophysical surveys}: seismic refraction and electrical resistivity to map bedrock depth and fracture zones between boreholes.
\item \textbf{Synthesis}: production of an engineering geological map and report recommending (or rejecting) the site.
\end{enumerate}
\end{methode}

\begin{attention}[The cost of skipping site investigation]
A borehole campaign costs a small fraction of the project budget, yet it routinely saves the whole project. Dams built on unexplored foundations have leaked, cracked or failed; roads cut without studying joint orientations have been buried by landslides. \textbf{Never site a major structure on geological guesswork.}
\end{attention}

\courssub{Dams: Site Criteria and Dam Types}

A dam stores water by blocking a valley. The ideal site is dictated by geometry and by geology.

\textbf{Favourable geological criteria}
\begin{itemize}
\item A \cle{narrow valley} (a gorge), which minimises the length --- and therefore the cost --- of the dam wall while maximising the stored volume upstream.
\item A \textbf{strong, impermeable foundation}: massive, unweathered rock (granite, gneiss, basalt) able to bear the load and resist seepage.
\item \textbf{Low jointing and no active faults}: joints and faults are pathways for leakage and planes of weakness along which the foundation can slide.
\item Abundant local \textbf{construction materials} (sand, gravel, rock fill, clay core material).
\item A reservoir basin floored by \textbf{impermeable rocks}, so that stored water does not escape sideways or downwards.
\end{itemize}

\begin{methode}[Geological checklist for siting a dam]
\begin{enumerate}
\item Is the valley narrow at the dam axis and wide upstream (good storage-to-cost ratio)?
\item Is the foundation rock massive, fresh and strong? Check borehole cores for weathering and fracture frequency.
\item Are joints frequent, open, or dipping downstream (sliding risk)? Is any fault crossing the site?
\item Is the rock permeable? Will a grout curtain be needed?
\item Do permeable beds (sandstone, limestone with karst) or faults cross the reservoir rim, allowing leakage out of the basin?
\item Is the catchment heavily eroded (silting risk)? Is the region seismic?
\item Are sand, gravel and clay available nearby for construction?
\end{enumerate}
If several answers are unfavourable, change the site --- not the geology.
\end{methode}

\textbf{Main dam types}

\begin{center}
\begin{tabularx}{\linewidth}{|l|X|X|}
\hline
\rowcolor{popblueL} \textbf{Type} & \textbf{Principle} & \textbf{Geological requirement} \\
\hline
Gravity dam & Resists water pressure by its own weight; usually concrete, triangular cross-section & Very strong, rigid foundation (massive igneous or metamorphic rock) \\
\hline
Arch dam & Curved wall transfers the water load sideways into the valley walls & Narrow gorge with strong, unjointed abutments on both flanks \\
\hline
Earth-fill / rock-fill dam & Embankment of compacted earth or rock with an impermeable clay core & Tolerates weaker, compressible foundations; needs wide valley and local fill material \\
\hline
\end{tabularx}
\end{center}

\begin{popfigure}
\begin{tikzpicture}[scale=0.85, font=\small]
% ---- Favourable: massive bedrock ----
\begin{scope}
\fill[popblueL] (0,0) rectangle (6.4,-3);
\draw[popdark, thick] (0,0) -- (6.4,0);
\node[popdark] at (3.2,-2.6) {\textbf{Massive granite}};
% water
\fill[popblue!35] (0,0) -- (0,1.6) -- (2.6,1.6) -- (2.6,0) -- cycle;
\draw[popblue] (0,1.6) -- (2.6,1.6);
% dam
\fill[gray!55] (2.6,0) -- (3.6,0) -- (3.1,2.0) -- (2.85,2.0) -- cycle;
\draw[black] (2.6,0) -- (3.6,0) -- (3.1,2.0) -- (2.85,2.0) -- cycle;
\node at (4.6,1.7) {Gravity dam};
\node[popgreen!60!black] at (3.2,-1.2) {no leakage path};
\node at (3.2,2.5) {\textbf{Favourable site}};
\end{scope}
% ---- Unfavourable: jointed bedrock ----
\begin{scope}[xshift=8.2cm]
\fill[popblueL] (0,0) rectangle (6.4,-3);
\draw[popdark, thick] (0,0) -- (6.4,0);
% joints dipping downstream
\foreach \x in {0.4,1.4,2.4,3.4,4.4,5.4}{\draw[poporange, thick] (\x,0) -- (\x+1.0,-3);}
\node[popdark] at (5.4,-2.6) {\textbf{Jointed rock}};
\fill[popblue!35] (0,0) -- (0,1.6) -- (2.6,1.6) -- (2.6,0) -- cycle;
\draw[popblue] (0,1.6) -- (2.6,1.6);
\fill[gray!55] (2.6,0) -- (3.6,0) -- (3.1,2.0) -- (2.85,2.0) -- cycle;
\draw[black] (2.6,0) -- (3.6,0) -- (3.1,2.0) -- (2.85,2.0) -- cycle;
% leakage arrows
\draw[->, popblue, thick] (2.5,-0.3) -- (3.6,-1.4);
\draw[->, popblue, thick] (2.3,-0.6) -- (3.4,-1.9);
\node[poporange!80!black] at (4.7,-0.8) {leakage \& sliding};
\node at (3.2,2.5) {\textbf{Unfavourable site}};
\end{scope}
\end{tikzpicture}
\legende{Cross-sections of a gravity dam on (left) massive, unjointed bedrock --- favourable; (right) bedrock with joints dipping downstream --- unfavourable: joints allow leakage and provide sliding planes under the dam.}
\end{popfigure}

\courssub{Reservoir Problems}

Impounding water creates new geological problems:

\begin{itemize}
\item \textbf{Leakage}: water escapes through permeable beds (sandstone, fractured lava, karstic limestone) or along faults that cross the reservoir rim. A reservoir can literally empty itself underground.
\item \textbf{Silting (sedimentation)}: rivers carry sand and silt into the reservoir, where velocity drops and sediment settles. The storage capacity shrinks year by year, especially where the catchment is deforested and deeply weathered, as on the savanna plateaux of northern Cameroon.
\item \textbf{Reservoir-induced seismicity}: the weight of impounded water and its infiltration along faults can trigger small to moderate earthquakes.
\item \textbf{Bank instability}: rising water saturates slope materials around the reservoir rim, provoking landslides.
\end{itemize}

\begin{savaistu}[Cameroonian examples: Lagdo and Lom Pangar]
The \textbf{Lagdo dam} (earth-fill dam on the Benue River, North Region) regulates floods, supplies irrigation and generates hydroelectricity; its management requires constant attention to silting from the eroded Benue catchment and to downstream flood releases. The \textbf{Lom Pangar dam} (East Region), built on the Sanaga basin, regulates the flow of the Sanaga River to improve the output of downstream hydroelectric plants at Song Loulou and Ed\'ea. Both projects required detailed geological studies of foundations, borrow materials and reservoir tightness before construction.
\end{savaistu}

\courssub{Tunnels}

Tunnels are excavated for roads, railways, water transfer and hydroelectric pressure shafts. Their behaviour is governed by the rock mass:

\begin{itemize}
\item \textbf{Rock type}: massive, strong rock (granite, gneiss) may stand unsupported; weak rocks (shale, highly weathered material) deform and collapse.
\item \textbf{Joints}: closely spaced joints divide the rock into blocks that can fall from the tunnel roof; joint orientation relative to the tunnel axis controls block stability.
\item \textbf{Faults}: fault zones are crushed, weak and often water-bearing --- the worst ground a tunneller can meet.
\item \textbf{Water inflow}: permeable rocks and fractures below the water table drain into the tunnel, flooding the works and reducing face stability.
\end{itemize}

\textbf{Support and lining}: depending on ground quality, tunnels are supported by steel arches, \cle{rock bolts} (steel rods anchored into the rock to stitch blocks together), shotcrete (sprayed concrete), and finally a permanent concrete lining. Water is controlled by drainage channels and by grouting (injecting cement into fractures).

\courssub{Slope Stability and Mass Movement}

A slope is a permanent battle between \textbf{driving forces} (gravity pulling material downslope) and \textbf{resisting forces} (the shear strength of the rock or soil). Stability is expressed quantitatively:

\begin{definition}[Factor of safety of a slope]
The \cle{factor of safety} $F$ of a slope is the ratio of resisting forces to driving forces along the potential slip surface:
\[
F \;=\; \dfrac{\text{sum of resisting (shear strength) forces}}{\text{sum of driving (gravitational) forces}}.
\]
If $F > 1$ the slope is stable; if $F = 1$ it is at the point of failure; if $F < 1$ it fails. Engineers typically design for $F \geq 1.5$.
\end{definition}

\textbf{Factors controlling slope stability}
\begin{itemize}
\item \textbf{Slope angle}: the steeper the slope, the greater the driving force. Over-steepening by road cuts or quarrying is a classic cause of failure.
\item \textbf{Rock strength}: weak or weathered materials (shale, lateritic soil, saprolite) fail more easily than fresh massive rock.
\item \textbf{Joint and bedding orientation}: discontinuities dipping \emph{out of} the slope face (daylighting) are planes on which blocks slide; discontinuities dipping \emph{into} the slope are generally stable.
\item \textbf{Water content}: see the property below --- the single most important trigger.
\item \textbf{Vegetation}: roots bind the soil and plants pump out water; deforestation destabilises slopes.
\item \textbf{Vibrations}: earthquakes, blasting and heavy traffic shake slopes and can trigger failure, especially when pores are saturated.
\end{itemize}

\begin{propriete}[Water reduces shear strength]
Water reduces the shear strength of slope materials in two ways: (1) it adds \textbf{weight}, increasing the driving force; (2) more importantly, it increases \textbf{pore water pressure}, which pushes grains apart and reduces the effective normal stress, and hence the frictional resistance, along the slip surface. This is why the great majority of landslides in the West and North-West Regions of Cameroon occur during or just after heavy rains. (Statement examinable; no derivation required.)
\end{propriete}

\textbf{Mass movement types relevant to engineering}
\begin{itemize}
\item \textbf{Rockfall}: free fall of detached blocks from steep cliffs --- common along the Dschang cliffs and the Bamenda escarpment, where jointed volcanic and crystalline rocks overhang roads.
\item \textbf{Slide (translational)}: movement of a mass along a planar surface, usually a bedding plane or joint dipping out of the slope.
\item \textbf{Slump (rotational slide)}: movement along a curved, spoon-shaped surface; the head sinks and tilts backward while the toe bulges forward. Typical of thick soils and weathered material on moderate slopes.
\end{itemize}

\begin{popfigure}
\begin{tikzpicture}[scale=0.9, font=\small]
% ground
\fill[popgreenL] (0,0) -- (0,2.2) -- (2.2,2.2) .. controls (3.4,2.0) and (4.0,1.2) .. (5.2,0.4) -- (8,0.4) -- (8,0) -- cycle;
\draw[popdark] (0,2.2) -- (2.2,2.2) .. controls (3.4,2.0) and (4.0,1.2) .. (5.2,0.4) -- (8,0.4);
% slip surface
\draw[poporange, very thick] (2.4,2.2) .. controls (3.2,0.6) and (4.6,-0.1) .. (5.6,0.4);
% slumped block
\draw[popink, thick, dashed] (2.4,2.2) .. controls (3.0,1.6) and (3.4,1.0) .. (4.4,0.9);
% water table
\draw[popblue, dashed, thick] (0,1.3) -- (8,1.3);
\node[popblue] at (7.2,1.55) {water table};
% labels
\node at (1.1,2.45) {crown / scarp};
\draw[->, popink] (2.4,2.5) -- (2.4,2.25);
\node[poporange!80!black] at (4.6,-0.35) {curved slip surface};
\node at (3.3,1.85) {head (back-tilted)};
\node at (6.2,0.75) {toe (bulge)};
\draw[->, popink] (5.6,0.85) -- (5.35,0.5);
% rotation arrow
\draw[->, popdark, thick] (3.0,1.5) arc (140:60:0.7);
\node at (2.9,1.15) {rotation};
\end{tikzpicture}
\legende{Rotational slump: the mass moves along a spoon-shaped slip surface; the head sinks and tilts backward, the toe bulges. A high water table (dashed) raises pore pressure and triggers failure.}
\end{popfigure}

\begin{popfigure}
\begin{tikzpicture}[scale=0.9, font=\small]
% LEFT: joints dip into slope - stable
\begin{scope}
\fill[popblueL] (0,0) -- (0,2.6) -- (3.2,2.6) -- (4.6,0) -- cycle;
\draw[popdark, thick] (0,2.6) -- (3.2,2.6) -- (4.6,0);
\foreach \t in {0.25,0.55,0.85,1.15,1.45,1.75,2.05,2.35}{\draw[poporange] (0.4,\t) -- (1.6,\t+0.9);}
\foreach \t in {0.2,0.7,1.2}{\draw[poporange] (2.0,\t+0.9) -- (2.9,\t+1.6);}
\node at (2.3,-0.5) {\textbf{Stable}: joints dip};
\node at (2.3,-1.0) {\textbf{into} the road cut};
\draw[->, popgreen!60!black, thick] (1.0,1.2) -- (0.4,1.7);
\node[popgreen!60!black] at (0.6,2.1) {dip};
\end{scope}
% RIGHT: joints dip out of slope - unstable
\begin{scope}[xshift=6.4cm]
\fill[popblueL] (0,0) -- (0,2.6) -- (3.2,2.6) -- (4.6,0) -- cycle;
\draw[popdark, thick] (0,2.6) -- (3.2,2.6) -- (4.6,0);
\foreach \t in {0.3,0.9,1.5,2.1}{\draw[poporange, thick] (1.2,\t+0.5) -- (3.6,\t-0.5);}
% sliding block
\fill[poppink!45] (2.5,1.9) -- (3.9,1.2) -- (4.3,0.7) -- (3.2,1.0) -- cycle;
\draw[popink] (2.5,1.9) -- (3.9,1.2) -- (4.3,0.7);
\draw[->, poppink!70!black, very thick] (3.6,1.5) -- (4.6,0.9);
\node at (2.3,-0.5) {\textbf{Unstable}: joints dip};
\node at (2.3,-1.0) {\textbf{out of} the road cut};
\node[poppink!70!black] at (4.9,1.6) {sliding};
\end{scope}
\end{tikzpicture}
\legende{Joint orientation relative to a road-cut face. Left: joints dip into the slope, blocks cannot slide out --- stable. Right: joints ``daylight'' (dip out of the face), providing sliding planes --- unstable.}
\end{popfigure}

\begin{attention}[Common mistake: ignoring joint orientation relative to the slope face]
Students often write ``jointed rock = unstable slope''. This is incomplete. A \emph{closely jointed} slope can be perfectly stable if the joints dip \textbf{into} the face or are horizontal; conversely a single set of joints dipping \textbf{out of} the face at an angle less than the slope angle can cause catastrophic planar slides. In any GCE question on slope stability, always compare the \textbf{dip direction and dip angle of discontinuities} with the \textbf{orientation of the slope face}.
\end{attention}

\begin{exoresolu}[Assessing a proposed road cut]
A road cut near Dschang is excavated in jointed basalt. The cut face dips at $70^{\circ}$ towards the east. One joint set dips at $40^{\circ}$ towards the east; a second set dips at $55^{\circ}$ towards the west. The slope stands in a region of heavy seasonal rainfall. Assess the stability and propose two stabilisation measures.
\tcblower
\textbf{Step 1 --- compare orientations.} The east-dipping joints ($40^{\circ}$E) dip \emph{out of} the cut face and are less steep than the face ($40^{\circ} < 70^{\circ}$): they \emph{daylight} in the slope and form potential planar sliding surfaces. The west-dipping set dips \emph{into} the slope and is stable.

\textbf{Step 2 --- consider water.} Heavy rains raise pore water pressure along the joints, reducing frictional resistance; the factor of safety $F$ falls during the wet season.

\textbf{Conclusion:} the cut is \textbf{potentially unstable} by planar sliding along the east-dipping joints, especially during rains.

\textbf{Stabilisation measures (any two):}
\begin{itemize}
\item \textbf{Drainage}: cut-off drains above the face and horizontal drain holes to lower pore pressure.
\item \textbf{Rock bolting}: bolts anchored across the joints to clamp blocks to the stable mass behind.
\item Reduce the slope angle by \textbf{benching} (cutting the face into steps).
\end{itemize}
\end{exoresolu}

\courssub{Stabilisation Measures}

Engineering geology does not stop at diagnosis; it prescribes remedies:

\begin{itemize}
\item \textbf{Drainage} (the most effective and cheapest): surface ditches intercept runoff; subsurface drains and horizontal boreholes lower the water table and hence pore pressure.
\item \textbf{Retaining walls}: concrete or gabion (wire cages filled with stones) walls support the toe of the slope, adding resisting force.
\item \textbf{Benching / regrading}: cutting the slope into steps or flattening it reduces the driving force and catches falling debris.
\item \textbf{Rock bolting and shotcrete}: bolts stitch jointed rock together; sprayed concrete protects the face from weathering and ravelling.
\item \textbf{Reforestation}: deep-rooted trees bind the soil, and evapotranspiration removes water --- a long-term, low-cost protection well suited to the humid western highlands of Cameroon.
\end{itemize}

\begin{aretenir}[Key points]
\begin{itemize}
\item Engineering geology studies the ground as a construction material and as a hazard; every major work begins with a staged site investigation.
\item A good dam site combines a narrow valley, a strong impermeable foundation, low jointing and a tight reservoir basin; dam type (gravity, arch, earth-fill) is chosen to match the geology.
\item Reservoirs suffer leakage, silting and induced seismicity; tunnels are controlled by rock strength, joints, faults and water inflow.
\item Slope stability is expressed by the factor of safety $F = \dfrac{\text{resisting forces}}{\text{driving forces}}$; failure occurs when $F < 1$.
\item Water is the main landslide trigger because pore pressure reduces shear strength; joint orientation relative to the slope face decides whether joints are dangerous.
\item Stabilisation combines drainage, retaining structures, benching, bolting and reforestation.
\end{itemize}
\end{aretenir}

\begin{exercice}[Slump along a hill road]
After exceptional rains, a section of road on the flank of the Bamenda escarpment drops by two metres along a curved scarp, while the ground below the road bulges upward. (a) Name the type of mass movement and justify your answer. (b) Explain the role of the rainfall. (c) Propose three stabilisation measures, explaining how each one raises the factor of safety.
\end{exercice}

\begin{corrige}[Exercise 1]
(a) A \textbf{rotational slump}: the curved scarp, the sinking and back-tilting of the road (head) and the bulging toe are diagnostic of movement along a spoon-shaped slip surface.

(b) Rainwater infiltrated the slope, adding weight and, above all, raising \textbf{pore water pressure} along the potential slip surface. This reduced the effective normal stress and therefore the frictional shear strength, so that $F$ fell below 1 and failure occurred.

(c) \textbf{Drainage} (surface ditches plus horizontal drains): lowers pore pressure, restoring shear strength. \textbf{Retaining wall or gabions at the toe}: adds resisting force against the bulging toe. \textbf{Regrading/benching of the slope}: reduces the slope angle and hence the driving force. (Reforestation is an acceptable long-term alternative: roots bind the soil and evapotranspiration removes water.) Each measure either increases resisting forces or decreases driving forces, raising $F$ above 1.
\end{corrige}

\coursec{Prospection and Exploration Methods}

In the previous section, we saw how geological knowledge is applied to civil engineering projects like dams, tunnels, and slope stability. We now turn to the upstream side of the mining and petroleum industry: \textbf{how do we find the resources in the first place?} Discovering an economic mineral deposit or a hydrocarbon trap is a scientific detective story. It progresses from regional reconnaissance over thousands of square kilometres to the precise definition of ore reserves measured in metres. This section details the systematic stages, the geological, geochemical, and geophysical tools, and the critical transition from \emph{prospection} (searching) to \emph{exploration} (evaluating).

\courssub{Stages of Mineral Exploration: From Reconnaissance to Feasibility}

Mineral exploration is a funnel-shaped process: the area of interest shrinks while the density of data and the cost per unit area increase dramatically. Each stage acts as a filter, discarding unpromising ground and focusing resources on targets with higher probability of success.

\begin{popfigure}
\begin{tikzpicture}[node distance=1.5cm and 2cm, >=Stealth,
    stage/.style={rectangle, draw=popdark, thick, fill=popblueL, text width=3.2cm, align=center, minimum height=1.8cm, rounded corners},
    arrow/.style={->, thick, popdark},
    decision/.style={diamond, draw=poporange, thick, fill=poporange!20, text width=2.5cm, align=center, aspect=2}
]
\node[stage, align=center] (recon){\textbf{1. Reconnaissance}\\(Regional Scale\\1:250,000 -- 1:1,000,000)\\\textit{Desk study, Remote sensing,\\ Regional geochem/geophysics}};
\node[stage, below of=recon, align=center] (prospect){\textbf{2. Prospection}\\(Target Scale\\1:50,000 -- 1:25,000)\\\textit{Geological mapping,\\ Stream sediment/soil grids,\\ Ground geophysics}};
\node[stage, below of=prospect, align=center] (detailed){\textbf{3. Detailed Exploration}\\(Prospect Scale\\1:10,000 -- 1:2,000)\\\textit{Trenching, Pitting,\\ Diamond Drilling (Core),\\ Resource Estimation}};
\node[stage, below of=detailed, align=center] (feas){\textbf{4. Feasibility Study}\\(Mine Scale\\Detailed Engineering)\\\textit{Metallurgy, Hydrogeology,\\ Environmental Impact (ESIA),\\ Economic Modelling}};
\node[decision, right of=prospect, xshift=4cm] (dec1) {Anomaly?};
\node[decision, right of=detailed, xshift=4cm] (dec2) {Viable Resource?};
\node[decision, right of=feas, xshift=4cm] (dec3) {Profitable?};

\draw[arrow] (recon) -- (prospect);
\draw[arrow] (prospect) -- (detailed);
\draw[arrow] (detailed) -- (feas);
\draw[arrow] (feas) -- (dec3);
\draw[arrow] (prospect.east) -- node[above] {No} (dec1.west);
\draw[arrow] (dec1.south) -- ++(0,-1) -| node[near start, above] {Yes} (detailed);
\draw[arrow] (detailed.east) -- node[above] {No} (dec2.west);
\draw[arrow] (dec2.south) -- ++(0,-1) -| node[near start, above] {Yes} (feas);
\draw[arrow] (dec3.east) -- node[above] {Yes} node[below, xshift=1.5cm] {\textbf{Mine Development}} (dec3.east);
\draw[arrow] (dec3.south) -- ++(0,-1) node[below] {Abandon / Sell};
\end{tikzpicture}
\legende{The exploration funnel: successive stages reduce area and increase geological confidence. A "No" at any decision gate typically returns the project to an earlier stage or terminates it.}
\end{popfigure}

\begin{aretenir}[The Four Stages]
\begin{itemize}
\item \textbf{Reconnaissance (Desk Study \& Regional Survey):} Uses existing geological maps, satellite imagery (Landsat, Sentinel), and regional aeromagnetic/radiometric data (e.g., data from the \emph{Ministry of Mines, Industry and Technological Development -- MINMIDT} in Yaoundé). Goal: Identify \textbf{Target Areas} (e.g., the Pan-African belt in North Cameroon for gold, or the Douala/Kribi-Campo basins for hydrocarbons).
\item \textbf{Prospection (Follow-up):} "Boots on the ground". Geological mapping at 1:50,000, systematic stream sediment and soil sampling grids (e.g., 200m x 50m), ground geophysics (magnetics, IP). Goal: Define \textbf{Drill Targets}.
\item \textbf{Detailed Exploration (Drilling):} Diamond core drilling (HQ/NQ size) on a grid (e.g., 50m x 50m for indicated resources). Core logging, assaying, bulk density measurement. Goal: Convert \emph{Mineral Resources} (Inferred $\to$ Indicated $\to$ Measured) into \emph{Mineral Reserves} (Probable $\to$ Proven).
\item \textbf{Feasibility Study:} Not purely geological. Includes metallurgical testwork (recovery \%), geotechnical studies (pit slope angles), hydrogeology, Environmental and Social Impact Assessment (ESIA -- mandatory in Cameroon under Law No. 96/12 of 5 August 1996), and financial modelling (NPV, IRR). Goal: Bankable Feasibility Study (BFS) for project financing.
\end{itemize}
\end{aretenir}

\begin{attention}[Prospection vs. Exploitation -- A Classic GCE Trap]
\textbf{Never confuse \emph{prospection} (or exploration) with \emph{exploitation} (mining).}
\begin{itemize}
\item \textbf{Prospection/Exploration:} Searching for and evaluating a deposit. Low volume, high information density. Methods: mapping, drilling, geophysics. \textit{Permit: Research Permit (Permis de Recherche).}
\item \textbf{Exploitation (Mining):} Extracting the mineral for sale. High volume, low information density (you already know the geology). Methods: open pit, underground, processing plant. \textit{Permit: Mining Concession (Concession Minière).}
\end{itemize}
In exam questions asking for "methods of prospection", do \emph{not} describe open-pit mining or blasting.
\end{attention}

\courssub{Geological Methods: The Foundation}

Geological methods provide the \emph{context} without which geochemical and geophysical anomalies are uninterpretable.

\begin{methode}[Systematic Geological Mapping for Exploration]
\begin{enumerate}
\item \textbf{Lithological Mapping:} Identify host rocks (e.g., Banded Iron Formation for Fe, Birimian greenstones for Au, Cretaceous sandstones for U). Note alteration halos (sericitization, silicification, propylitic).
\item \textbf{Structural Mapping:} Measure strike/dip of bedding ($S_0$), foliation ($S_1$), faults, joints, veins. Mineralization is often structurally controlled (e.g., gold in shear zones, veins in tension gashes). Plot on stereonets to find intersection lineations.
\item \textbf{Stream Sediment Sampling:} The classic first-pass tool in Cameroon's rainforest (e.g., Batouri, Bétaré-Oya gold areas). Sample \textbf{B-horizon soil} or \textbf{active stream sediment} (fine fraction < 80 mesh) at confluences. Analyze for pathfinder elements (As, Sb, Hg for Au; Cu, Zn, Pb for VMS).
\item \textbf{Soil Sampling (Grid):} On a regular grid (e.g., 100m x 50m) over a target. Define \textbf{geochemical anomalies} (see box below). In lateritic terrain (South Cameroon), sample the \textbf{ferruginous duricrust} or \textbf{mottled zone} rather than transported topsoil.
\item \textbf{Trenching \& Pitting:} Mechanical excavation (backhoe) or hand-dug pits to expose bedrock below lateritic cover. Channel sampling across strike (e.g., 1m channels, 2kg samples). Essential for lateritic gold (e.g., \emph{Kamina}) or lateritic Ni-Co-Mn (e.g., \emph{Nkamouna}, \emph{Mada}) where drilling is premature.
\end{enumerate}
\end{methode}

\begin{definition}[Geochemical Anomaly]
A \textbf{geochemical anomaly} is a statistically significant concentration of one or more elements in a sample medium (rock, soil, sediment, water, vegetation) that is distinctly higher (or lower) than the \textbf{background value} (normal abundance for that rock type/environment) and which \emph{may} indicate the presence of mineralization.
\begin{itemize}
\item \textbf{Background:} Natural variation in barren rock. Established by statistical analysis (mean $\pm$ 2 standard deviations) of samples from unmineralized areas.
\item \textbf{Threshold:} The upper limit of background (e.g., Mean + 2$\sigma$). Values above are "anomalous".
\item \textbf{Contrast:} Ratio of anomaly peak to background. High contrast = easier detection.
\item \textbf{Dispersion Halo:} Primary (syngenetic/epigenetic in rock) or Secondary (mechanical/hydromorphic in soil/stream). In Cameroon's tropical climate, \textbf{hydromorphic dispersion} (downslope movement in groundwater) creates broad, diffuse anomalies downslope from the source. \textbf{Mechanical dispersion} (soil creep, erosion) spreads anomalies downslope physically.
\end{itemize}
\end{definition}

\courssub{Geophysical Methods: Seeing Through the Cover}

Geophysics measures physical property contrasts between the target ore body (or trap) and the host rock. The choice of method depends entirely on the \textbf{physical property contrast}.

\begin{propriete}[Matching Geophysical Method to Physical Property Contrast]
\begin{center}
\begin{tabularx}{\linewidth}{|l|X|X|X|}\hline
\rowcolor{popblueL} \textbf{Method} & \textbf{Physical Property Measured} & \textbf{Target Contrast} & \textbf{Typical Targets (Cameroon Context)} \\ \hline
\textbf{Magnetics} & Magnetic Susceptibility ($\kappa$, SI) & High $\kappa$ (Magnetite, Pyrrhotite) vs Low $\kappa$ (Sediments, Granite) & Banded Iron Formation (BIF), Mafic/Ultramafic complexes (Ni-Cu-PGE), Basement structure mapping (oil basins). \\ \hline
\textbf{Gravimetry} & Density ($\rho$, kg/m$^3$) & High $\rho$ (Massive sulfides, Chromite, BIF) or Low $\rho$ (Salt domes, Sedimentary basins) vs Host & Massive sulfide lenses (VMS), Chromite pods, Basin architecture (Douala Basin salt tectonics), Regional crustal studies. \\ \hline
\textbf{Resistivity / IP} & Electrical Resistivity ($\rho_e$, $\Omega$m) \& Chargeability ($M$, mV/V) & Conductive (Graphite, Massive Sulfides, Clay) vs Resistive (Quartzite, Fresh Granite). IP: Disseminated sulfides (chargeable). & \textbf{IP is king for disseminated porphyry Cu-Au and VMS.} Resistivity for graphite, groundwater, structural mapping. \\ \hline
\textbf{Seismic Reflection} & Acoustic Impedance ($Z = \rho \times V_p$) & Contrast in $\rho \times V_p$ at layer boundaries. High frequency = high resolution, low penetration. & \textbf{Petroleum exploration (Douala, Rio del Rey, Logone Birni basins).} Coal seams. Deep crustal structure. \\ \hline
\textbf{Seismic Refraction} & P-wave Velocity ($V_p$) & Velocity contrast at interfaces. Low resolution, deep penetration. & Depth to bedrock (engineering), Basement topography, Groundwater aquifers. \\ \hline
\textbf{Radiometrics} & Gamma rays (K, U, Th) & High K (Potassic alteration), U/Th (Uranium, Rare Earths, Granites). & Uranium (Polí area), REE (Lom/Pangar), Potassic alteration halos (Porphyry Cu). \\ \hline
\textbf{EM (Electromagnetic)} & Electrical Conductivity ($\sigma$) & High conductivity (Massive sulfides, Graphite, Saline water). & Airborne VTEM/HELITEM for VMS (Volcanogenic Massive Sulfides) in greenstone belts. \\ \hline
\end{tabularx}
\end{center}
\end{propriete}

\begin{experience}[Interpreting a Resistivity / IP Profile over a Conductive Ore Body]
Consider a profile across a steeply dipping massive sulfide lens (pyrite/chalcopyrite) hosted in resistive quartzite, covered by 20m of conductive lateritic clay.
\begin{popfigure}
\begin{tikzpicture}[scale=0.9]
% Surface
\draw[thick, popdark] (-5,0) -- (5,0);
\draw[popmuted] (-5,-0.1) -- (5,-0.1);
\node at (-5.5,0) {Surface};

% Laterite layer (Conductive)
\fill[poporange!30] (-5,0) -- (5,0) -- (5,-2) -- (-5,-2) -- cycle;
\draw[dashed] (-5,-2) -- (5,-2);
\node[align=center] at (0,-1){\textbf{Lateritic Clay}\\\textit{Low Resistivity ($\sim$10-50 $\Omega$m)}};

% Host Rock (Resistive)
\fill[popblue!10] (-5,-2) -- (5,-2) -- (5,-6) -- (-5,-6) -- cycle;
\node[align=center] at (0,-4){\textbf{Quartzite Host}\\\textit{High Resistivity ($>$1000 $\Omega$m)}};

% Ore Body (Conductive & Chargeable)
\fill[popred!50] (-0.5,-2.5) -- (0.5,-2.5) -- (0.8,-5) -- (-0.8,-5) -- cycle;
\draw[thick, popred] (-0.5,-2.5) -- (0.5,-2.5) -- (0.8,-5) -- (-0.8,-5) -- cycle;
\node[popred, font=\bfseries] at (0,-3.7) {Massive Sulfide Lens};

% Electrodes (Wenner Array)
\foreach \x in {-4,-2,0,2,4} {
    \draw[thick, popdark] (\x,0.1) -- (\x,-0.1);
    \node[below, font=\tiny] at (\x,0.15) {E};
}

% Apparent Resistivity Curve (Sketch)
\begin{scope}[yshift=-7cm, xshift=-5cm]
\draw[->] (0,0) -- (10,0) node[right] {Position (m)};
\draw[->] (0,0) -- (0,3) node[above] {$\rho_a$ ($\Omega$m)};
\draw[thick, popblue, domain=-4:4, samples=50, variable=\x] plot ({\x+5}, {2.5 - 2*exp(-(\x)^2/0.5)});
\node[popblue, anchor=south] at (5,2.5) {Apparent Resistivity Low over Ore};
\draw[dashed, popmuted] (0,2.5) -- (10,2.5) node[right] {Background};
\end{scope}

% Chargeability Curve (Sketch)
\begin{scope}[yshift=-7cm, xshift=5cm]
\draw[->] (0,0) -- (10,0) node[right] {Position (m)};
\draw[->] (0,0) -- (0,3) node[above] {$M$ (mV/V)};
\draw[thick, poporange, domain=-4:4, samples=50, variable=\x] plot ({\x+5}, {2*exp(-(\x)^2/0.8)});
\node[poporange, anchor=south] at (5,2) {Chargeability High over Ore};
\draw[dashed, popmuted] (0,0) -- (10,0);
\end{scope}
\end{tikzpicture}
\legende{Resistivity/IP profile over a massive sulfide body. The conductive laterite masks the resistive host. The ore body appears as a resistivity \textbf{low} (conductive) and a chargeability \textbf{high} (polarizable). IP distinguishes metallic sulfides (chargeable) from clay/graphite (non-chargeable).}
\end{popfigure}
\end{experience}

\courssub{Seismic Reflection: The Backbone of Petroleum Exploration}

While mining exploration relies heavily on magnetics, gravity, and IP, \textbf{petroleum exploration is dominated by seismic reflection}. Hydrocarbon traps (anticlines, fault blocks, stratigraphic pinch-outs) are large-scale structural/stratigraphic features requiring imaging of sedimentary layering at depths of 1--5 km.

\begin{propriete}[Seismic Reflection Principle]
An acoustic wave (source: Vibroseis on land, Airgun array offshore) travels downwards. At each interface where \textbf{Acoustic Impedance} $Z = \rho \times V_p$ changes, a fraction of energy is \textbf{reflected} back to the surface (geophones/hydrophones) and the rest is \textbf{transmitted}.
\begin{itemize}
\item \textbf{Two-Way Time (TWT):} $t = 2 \int_0^z \frac{dz}{V_p(z)}$. The recorded section is in \textbf{time (ms)}, not depth.
\item \textbf{Vertical Resolution:} $\approx \lambda/4 = V_p / (4f)$. High frequency ($f$) $\to$ high resolution but low penetration. Typical: 10--60 Hz (land), 10--120 Hz (marine).
\item \textbf{Processing Flow:} Demultiplex $\to$ Deconvolution $\to$ Velocity Analysis (CMP stacking) $\to$ Migration (collapses diffractions, puts dips in true position) $\to$ Time-to-Depth Conversion (using velocity model).
\end{itemize}
\end{propriete}

\begin{popfigure}
\begin{tikzpicture}[scale=0.85]
% Layers
\fill[popblue!5] (-6,0) -- (6,0) -- (6,-1) -- (-6,-1) -- cycle; % Surface layer
\fill[popgreen!10] (-6,-1) -- (6,-1) -- (6,-3) -- (-6,-3) -- cycle; % Reservoir interval
\fill[poporange!10] (-6,-3) -- (6,-3) -- (6,-5) -- (-6,-5) -- cycle; % Deeper

% Anticline Structure
\draw[thick, popdark] plot[smooth, tension=0.7] coordinates {(-6,-1) (-3,-0.5) (0,-0.2) (3,-0.5) (6,-1)};
\draw[thick, popdark] plot[smooth, tension=0.7] coordinates {(-6,-3) (-3,-2.2) (0,-1.8) (3,-2.2) (6,-3)};
\draw[thick, popdark] plot[smooth, tension=0.7] coordinates {(-6,-5) (-3,-4.5) (0,-4.2) (3,-4.5) (6,-5)};

% Fault cutting through
\draw[thick, popred, decorate, decoration={zigzag, segment length=4, amplitude=1}] (-1,-0.5) -- (-1.5,-5);
\node[popred, rotate=-15, font=\bfseries\small] at (-2.5,-2.5) {Normal Fault};

% Source and Receivers
\draw[popdark, thick] (0,0.3) circle (0.15) node[above] {Source (Vibroseis)};
\foreach \x in {-3,-2,-1,1,2,3} {
    \draw[popdark, thick] (\x,0.3) -- (\x,0.1) -- (\x+0.1,0.1) -- (\x+0.1,0.3);
    \node[font=\tiny] at (\x+0.05,0.4) {Geo};
}

% Ray paths (Source at 0, Reflector at anticline crest ~0,-0.2)
\draw[popblue, ->, thick] (0,0.3) -- (0,-0.2) node[midway, right] {Normal Incidence};
\draw[popblue, ->] (0,0.3) -- (-2,-0.6);
\draw[popblue, ->] (0,0.3) -- (2,-0.6);
\draw[popblue, ->] (-2,0.3) -- (-2,-0.6);
\draw[popblue, ->] (2,0.3) -- (2,-0.6);

% Synthetic Seismogram (Right side)
\begin{scope}[xshift=8cm, yshift=-1cm]
\draw[->] (0,0) -- (0,-4) node[below] {TWT (s)};
\draw[->] (0,0) -- (2,0) node[right] {Amplitude};
\foreach \y/\lab in {0/0.0, -1/0.5, -2/1.0, -3/1.5, -4/2.0} {
    \draw (\y,0.05) -- (\y,-0.05) node[left] {\tiny \lab};
}
% Wiggle trace sketch
\draw[thick, popdark] (0,0) 
    .. controls (0.5,-0.2) and (1.5,-0.5) .. (0,-0.8) % First arrival
    .. controls (-1.5,-1.0) and (1.5,-1.2) .. (0,-1.5) % Anticline crest reflection (Strong +ve)
    .. controls (-1.0,-1.8) and (1.0,-2.2) .. (0,-2.5) % Deeper reflector
    .. controls (-0.5,-2.8) and (0.5,-3.2) .. (0,-3.8);
\node[font=\small, align=center] at (1,-1.5) {Strong Reflection\\ (High Impedance\\ Contrast)};
\node[font=\small, align=center] at (1,-3) {Fault Plane\\ (Diffraction)};
\end{scope}

\node[align=center, font=\bfseries] at (0,0.8) {Seismic Section (Time Domain)};
\node[align=center, font=\small] at (0,-5.5) {Depth (km)};
\end{tikzpicture}
\legende{Seismic reflection profile across an anticlinal trap cut by a normal fault. The source generates a wavefield; reflections from impedance contrasts (e.g., shale/sand, top reservoir) are recorded. Processing (migration) collapses the diffraction hyperbola from the fault tip and positions the dipping reflectors correctly. The anticline crest shows a strong, continuous reflection (hydrocarbon trap potential).}
\end{popfigure}

\begin{savaistu}[Cameroon's Petroleum Basins \& Seismic]
Cameroon's hydrocarbon production comes from the \textbf{Rio del Rey Basin} (offshore/onshore, Cretaceous) and the \textbf{Douala/Kribi-Campo Basin} (offshore, Cretaceous to Tertiary). The \textbf{Logone Birni Basin} (inland, Cretaceous) is a frontier basin. Modern 3D seismic surveys (e.g., by SNH -- Société Nationale des Hydrocarbures, and partners like Perenco, Bowleven, New Age) have revolutionized the imaging of \textbf{rollover anticlines} and \textbf{stratigraphic traps} in the "Gamba" and "Mungo" formations. Seismic attribute analysis (RMS amplitude, coherence, spectral decomposition) is now standard for Direct Hydrocarbon Indicators (DHIs / "flat spots", "bright spots").
\end{savaistu}

\courssub{Remote Sensing and Aerial Photography in Reconnaissance}

Before a geologist sets foot in the forest, remote sensing defines the playing field.

\begin{itemize}
\item \textbf{Aerial Photography / High-Res Satellite (Pleiades, WorldView, Drone):} 0.3--0.5m resolution. Maps \textbf{geomorphology} (drainage patterns -- dendritic, trellis, radial indicating structure), \textbf{lineaments} (faults, fractures, dykes), and \textbf{alteration zones} (iron oxide staining visible in RGB/False Colour Composites). Essential for planning access in difficult terrain (e.g., Mount Cameroon flanks, Mamfe Basin).
\item \textbf{Multispectral (Landsat 8/9 OLI, Sentinel-2 MSI, ASTER):} 10--30m resolution. \textbf{Band Ratios} map hydrothermal alteration minerals:
    \begin{align*}
    \text{Iron Oxides (Hematite/Goethite)} &: \text{Band 4 / Band 2 (Red/Blue)} \text{ or } \text{Band 5 / Band 4 (SWIR/NIR)} \\
    \text{Clay Minerals (Kaolinite, Illite, Smectite)} &: \text{Band 6 / Band 7 (SWIR1/SWIR2)} \text{ (ASTER bands 4/5, 5/6, 6/7)} \\
    \text{Silica / Quartz} &: \text{Band 7 / Band 6}
    \end{align*}
    \textit{Example: Mapping the argillic/propylitic halo around a porphyry Cu-Au target in the North-West Region.}
\item \textbf{Hyperspectral (HyMap, PRISMA, EnMAP):} Hundreds of narrow bands. Identifies \textbf{specific mineral species} (e.g., distinguishing Al-OH in white mica -- phengite vs muscovite -- indicating temperature of alteration). Still largely research/airborne in Cameroon.
\item \textbf{Radar (SAR -- Sentinel-1, ALOS PALSAR):} Penetrates cloud cover (critical for Cameroon's rainforest!) and vegetation canopy (L-band). Detects surface roughness and structural lineaments invisible on optical data. Interferometry (InSAR) measures ground deformation (subsidence over mines, volcano monitoring on Mt. Cameroon).
\end{itemize}

\courssub{Exploration Drilling: Core, Logging, and Reserve Estimation}

Drilling provides the \textbf{ground truth}. It is the most expensive stage (cost per metre: \$100--\$300+ for diamond core), so hole placement must be justified by all previous data.

\begin{methode}[Diamond Core Drilling Workflow]
\begin{enumerate}
\item \textbf{Collar Survey:} Precise GPS (DGPS/RTK) for X, Y, Z. Azimuth and Dip set by rig alignment.
\item \textbf{Core Recovery \& RQD:} \textbf{Recovery \%} = (Length of core recovered / Length drilled) $\times$ 100. \textbf{RQD (Rock Quality Designation)} = \% of core pieces $> 10$ cm. Critical for geotechnical stability and sample representativity. Target: $> 90\%$ recovery, $> 75\%$ RQD in ore zone.
\item \textbf{Core Orientation:} Using a core orientation tool (e.g., Reflex ACT III). Allows measurement of \textbf{true dip/dip direction} of structures (veins, foliation, contacts) in the core reference frame. Essential for structural modelling.
\item \textbf{Geotechnical Logging:} Fracture frequency, joint condition, weathering grade (ISRM scale).
\item \textbf{Geological Logging:} Lithology, alteration, mineralization (%), structure (alpha/beta angles). Standardized codes (e.g., Geobank, AcQuire).
\item \textbf{Sampling:} \textbf{Half-core} sawn lengthwise (quarter-core for duplicates). Sample intervals honour geological boundaries (min 0.3m, max 1.5m typically). \textbf{QA/QC:} Insert Certified Reference Materials (CRMs/Standards), Blanks, Field Duplicates, Pulp Duplicates $\approx$ 5--10\% insertion rate.
\item \textbf{Downhole Surveys:} Single-shot or multi-shot (Gyro/Reflex EZ-Shot) every 30--50m to track hole deviation (azimuth/dip change).
\item \textbf{Downhole Geophysics (Logging):} Run in open hole or cased hole.
    \begin{itemize}
    \item \textbf{Gamma / Density / Neutron:} Lithology, porosity (oil).
    \item \textbf{Resistivity / Induction:} Fluid saturation (oil), massive sulfides.
    \item \textbf{Deviation Survey:} Continuous (Gyro).
    \item \textbf{Borehole Camera / Televiewer (Acoustic/Optical):} Structural orientation, fracture aperture (in situ stress).
    \end{itemize}
\end{enumerate}
\end{methode}

\begin{definition}[Mineral Resources vs. Mineral Reserves (JORC / NI 43-101 / SAMREC Codes)]
Exploration results must be reported using standardized codes (Canada: NI 43-101; Australia: JORC; South Africa: SAMREC; Europe: PERC). Cameroon follows international best practice (often JORC/NI 43-101 for listed companies).
\begin{center}
\begin{tabularx}{\linewidth}{|l|X|X|}\hline
\rowcolor{popblueL} \textbf{Category} & \textbf{Definition} & \textbf{Geological Confidence / Data Spacing} \\ \hline
\textbf{Inferred Resource} & Quantity/Grade estimated from limited sampling. \textbf{Cannot be used for feasibility.} & Wide spacing (e.g., 100m x 100m). Geological continuity assumed. \\ \hline
\textbf{Indicated Resource} & Quantity/Grade estimated from adequate sampling. \textbf{Can support PFS (Pre-Feasibility).} & Moderate spacing (e.g., 50m x 50m). Geological continuity confirmed. \\ \hline
\textbf{Measured Resource} & Quantity/Grade estimated from detailed sampling. \textbf{Can support FS (Feasibility).} & Close spacing (e.g., 25m x 25m or less). High confidence in geology/grade. \\ \hline
\textbf{Probable Reserve} & Indicated Resource + Modifying Factors (mining, metallurgy, economic, legal, ESG) = \textbf{Mineable}. & Includes dilution, recovery, cut-off grade. Positive cash flow at PFS level. \\ \hline
\textbf{Proven Reserve} & Measured Resource + Modifying Factors = \textbf{Mineable}. & Highest confidence. Detailed mine design possible. \\ \hline
\end{tabularx}
\end{center}
\textbf{Key Rule:} You \textbf{cannot} convert Inferred Resources directly to Reserves. They must first be upgraded to Indicated/Measured by infill drilling.
\end{definition}

\begin{exoresolu}[Resource Estimation: Inverse Distance Weighting (IDW) -- Conceptual]
\textbf{Statement:} A vein system is drilled on a 50m x 50m grid. Four composite samples surround a block centroid at coordinates (0,0). Grades and distances:
\begin{itemize}
\item Sample A (50m East): 4.2 g/t Au, Distance = 50m
\item Sample B (50m West): 3.8 g/t Au, Distance = 50m
\item Sample C (50m North): 5.0 g/t Au, Distance = 50m
\item Sample D (50m South): 2.0 g/t Au, Distance = 50m
\end{itemize}
Estimate the block grade using Inverse Distance Squared ($ID^2$) weighting.

\tcblower
\textbf{Solution:}
Weight $w_i = \frac{1}{d_i^2}$. Since all distances are equal ($d=50$), all weights are equal: $w = 1/2500$.
Estimated Grade $= \frac{\sum (w_i \times g_i)}{\sum w_i} = \frac{w(4.2 + 3.8 + 5.0 + 2.0)}{4w} = \frac{15.0}{4} = \mathbf{3.75\ g/t\ Au}$.

\textbf{Comment:} With equal spacing, ID$^2$ reduces to a simple arithmetic mean. In reality, distances vary. \textbf{Kriging} (geostatistics) is the industry standard because it minimizes estimation variance and provides a kriging variance (error estimate), but ID$^2$ is a valid deterministic method for GCE-level understanding. Always honour geological boundaries (domaining) before estimating!
\end{exoresolu}

\courssub{Choosing the Exploration Strategy: Commodity-Specific Workflows}

A competent exploration geologist does not apply every method everywhere. The strategy is dictated by the \textbf{deposit model} (geological concept of the target).

\begin{center}
\begin{tabularx}{\linewidth}{|l|X|X|X|}\hline
\rowcolor{popblueL} \textbf{Commodity / Deposit Type} & \textbf{Key Physical/Chemical Signature} & \textbf{Priority Methods (Ranked)} & \textbf{Cameroon Example} \\ \hline
\textbf{Petroleum (Oil \& Gas)} & Acoustic Impedance Contrast (Traps), Source Rock Maturity, Seal & \textbf{1. Seismic Reflection (2D/3D)} $\to$ 2. Gravity/Magnetics (Basement) $\to$ 3. Geochemistry (Oil seeps, source rock) & Rio del Rey Basin (Offshore), Logone Birni (Onshore) \\ \hline
\textbf{Volcanogenic Massive Sulfides (VMS) Cu-Zn-Pb-Ag-Au} & \textbf{High Density, High Magnetite/Pyrrhotite, High Conductivity, Strong IP Chargeability} & \textbf{1. Airborne EM (VTEM) + Magnetics} $\to$ 2. Ground IP/Resistivity $\to$ 3. Soil Geochem (Cu, Zn, Pb, As) $\to$ 4. Drilling & Potential in Pan-African belt (North Cameroon), though less explored than W. Africa. \\ \hline
\textbf{Orogenic Gold (Shear-hosted, Quartz Veins)} & \textbf{Geochemical (Au, As, Sb, W)}, Structural Control, Low Phys. Contrast (Quartz veins = Resistive, Non-mag) & \textbf{1. Structural Mapping + Soil/Stream Geochem (Au pathfinders)} $\to$ 2. Ground Magnetics (Structures) $\to$ 3. IP/Resistivity (Sulfides in wallrock) $\to$ 4. Drilling & \textbf{Batouri, Bétaré-Oya, Kamina} (East/North Regions). Laterite cover requires deep soil/auger drilling. \\ \hline
\textbf{Porphyry Cu-Au / Epithermal} & \textbf{Stockwork veining (Chargeable), Magnetite destruction (Mag Low), Alteration Zoning (Spectral)} & \textbf{1. Magnetics (Mag low core)} $\to$ 2. IP/Resistivity (Chargeability shell) $\to$ 3. Hyperspectral/Alteration Mapping $\to$ 4. Soil Geochem (Cu, Mo, Au) $\to$ 5. Drilling & Potential in Pan-African granitoids (North Cameroon). Rare in Cenozoic volcanic line. \\ \hline
\textbf{Bauxite / Lateritic Ni-Co-Mn} & \textbf{Surface/Sub-surface Chemistry, Density contrast (Low density ore vs Saprock)} & \textbf{1. Geomorphology (Plateaus)} $\to$ 2. Auger/RC Drilling Grid $\to$ 3. Geochem (Al, Fe, Ni, Co, Mn) $\to$ 4. Density logging & \textbf{Minim-Martap, Ngaoundal (Bauxite)}; \textbf{Nkamouna, Mada (Ni-Co-Mn)} \\ \hline
\textbf{Heavy Mineral Sands (Ilmenite, Rutile, Zr, Monazite)} & \textbf{High Density, Magnetic (Ilmenite), Radiometric (Monazite/Th)} & \textbf{1. Airborne Radiometrics + Magnetics} $\to$ 2. Auger Drilling (Beach/Strandlines) $\to$ 3. Heavy Mineral Separation (Lab) & \textbf{Kribi-Campo Coast} (Potential). \\ \hline
\end{tabularx}
\end{center}

\begin{attention}[Common Examination Pitfalls]
\begin{enumerate}
\item \textbf{"Geophysics finds minerals."} \textbf{False.} Geophysics finds \textbf{physical property anomalies}. Only drilling + assay confirms minerals. A magnetic high could be a barren mafic dyke; an IP anomaly could be graphite or clay.
\item \textbf{"Seismic reflection works for gold."} \textbf{Generally False.} Seismic wavelengths (10--50m) are too long to resolve narrow gold veins (cm--m scale). Seismic is for \textbf{large-scale structures} (basins, faults at depth). Use IP/Resistivity/Magnetics for gold.
\item \textbf{Confusing "Anomaly" with "Ore".} An anomaly is a \emph{deviation from background}. It requires geological validation. Many anomalies are "false positives" (e.g., cultural noise, black shales, laterite).
\item \textbf{Ignoring "Modifying Factors" in Reserve Definition.} A high-grade Measured Resource is \textbf{not} a Proven Reserve if metallurgical recovery is unknown, or if there is no mining licence, or if the ESIA is rejected. Reserves are \emph{economic}; Resources are \emph{geological}.
\end{enumerate}
\end{attention}

\begin{exercice}[Integrated Exploration Design]
A junior company acquires a Permis de Recherche over a 500 km$^2$ area in the \textbf{Pan-African Belt of North Cameroon} (greenstone belt setting). Regional data suggests potential for \textbf{Orogenic Gold} and \textbf{VMS (Cu-Zn)}.
\begin{enumerate}
\item Design a \textbf{3-year exploration program} (Year 1: Reconnaissance; Year 2: Prospection; Year 3: Detailed Exploration) with budget allocation rationale (\% of total).
\item For the \textbf{VMS target}, specify the optimal airborne geophysical survey parameters (sensor, line spacing, direction).
\item For the \textbf{Gold target}, explain why stream sediment sampling is more effective than soil sampling in Year 1, and why soil sampling becomes critical in Year 2.
\item Define the drilling grid spacing required to declare an \textbf{Indicated Resource} for a tabular vein system dipping 60$^\circ$.
\end{enumerate}
\end{exercice}

\begin{corrige}[Exercise: Integrated Exploration Design]
\begin{enumerate}
\item \textbf{Year 1 (Reconnaissance -- 15\% Budget):} Compile existing MINMIDT data, geological maps, aeromagnetics. Acquire/Process Sentinel-2/Landsat/ASTER for alteration/lineaments. \textbf{Regional Stream Sediment Survey} (1 sample/km$^2$ at confluences) for Au, As, Cu, Zn. \textbf{Airborne Magnetics + Radiometrics} (if not existing) at 200m line spacing. \textit{Output: Target Areas (5--10 km$^2$ each).}
\item \textbf{Year 2 (Prospection -- 35\% Budget):} \textbf{Geological Mapping 1:25,000} on targets. \textbf{Infill Stream Sediment} (1 sample/0.5 km$^2$). \textbf{Soil Grids} (200m x 50m) over geochem anomalies + structural corridors (Au pathfinders: As, Sb, W). \textbf{Ground IP/Resistivity + Magnetics} profiles across VMS conductors. \textit{Output: Drill Targets (Ranked A, B, C).}
\item \textbf{Year 3 (Detailed Exploration -- 50\% Budget):} \textbf{Diamond Drilling} (HQ/NQ) on Rank A targets. Grid spacing 50m x 50m (Indicated). Core orientation, density, QA/QC. \textbf{Metallurgical testwork} (gravity, cyanidation, flotation). \textbf{Resource Estimation} (JORC/NI 43-101). \textit{Output: Maiden Resource Statement.}
\item \textbf{VMS Airborne Survey:} \textbf{Time-Domain EM (VTEM or ZTEM)} + \textbf{High-Res Magnetics}. Line spacing \textbf{100--200m}. Flight height \textbf{30--50m} (drape). Lines \textbf{perpendicular to regional strike} (typically NW-SE for NE-SW striking Pan-African belts) to cross-cut conductors. High sampling rate (0.1s).
\item \textbf{Gold Sampling Logic:} Year 1: \textbf{Stream sediment} integrates large catchment areas (km$^2$) cheaply; detects "blind" mineralization under laterite cover via hydromorphic dispersion. Year 2: \textbf{Soil grids} provide high-resolution (metre-scale) anomaly shape, orientation, and tenor over a specific target defined in Year 1; essential for drill hole placement.
\item \textbf{Drilling Grid for Indicated Resource (Vein):} Geological continuity must be proven along strike and down-dip. For a tabular vein: \textbf{50m along strike $\times$ 50m down-dip} (or 40m x 40m). Drill holes typically fanned from surface or underground platforms to intersect vein at $\sim$ perpendicular angles. Two holes intersecting the vein at the same level (separated by 50m) + continuity of grade/width $\to$ Indicated.
\end{enumerate}
\end{corrige}

\begin{aretenir}[Key Takeaways for GCE Revision]
\begin{itemize}
\item Exploration is a \textbf{funnel}: Reconnaissance $\to$ Prospection $\to$ Detailed Exploration $\to$ Feasibility. Cost $\uparrow$, Area $\downarrow$, Confidence $\uparrow$.
\item \textbf{Geology first:} Mapping and structural control guide everything. Geophysics/Geochem test geological hypotheses.
\item \textbf{Method Selection:} Match the method to the \textbf{Physical Property Contrast} of the \textbf{Deposit Model}.
    \begin{itemize}
    \item Oil/Gas $\to$ \textbf{Seismic Reflection} (Acoustic Impedance).
    \item Massive Sulfides (VMS) $\to$ \textbf{Airborne EM + Magnetics + Ground IP}.
    \item Orogenic Gold $\to$ \textbf{Structural Mapping + Soil Geochem (Pathfinders) + Ground Mag/IP}.
    \item Lateritic Ni/Co/Bauxite $\to$ \textbf{Auger Drilling Grid + Geochem}.
    \end{itemize}
\item \textbf{Drilling:} Diamond core (orientation, recovery, RQD, QA/QC) is the only way to get \textbf{Reserves}.
\item \textbf{Reporting Codes:} Know the hierarchy: Inferred $\to$ Indicated $\to$ Measured (Resources) $\to$ Probable $\to$ Proven (Reserves). Modifying factors are mandatory for Reserves.
\item \textbf{Cameroon Context:} Tropical laterite cover demands deep soil/auger sampling; rainforest demands Radar/Remote Sensing; Pan-African belt = Gold/VMS potential; Sedimentary Basins = Oil/Gas potential.
\end{itemize}
\end{aretenir}

\coursec{Mining Methods and the Environment}

In the previous section we learned how geologists \emph{find} an ore body: geological mapping, geochemical sampling, geophysical surveys and finally drilling. But discovering a deposit is only the beginning of the story. Imagine a team that has just confirmed, in the East Region of Cameroon, a gold-bearing zone a few tens of metres below the forest. A new question immediately arises: \textbf{how do we actually get the ore out of the ground, and at what cost to people and nature?} This final section answers that question. We first describe the two great families of mining methods --- \cle{surface mining} and \cle{underground mining} --- then the criteria used to choose between them, the life cycle of a mine, and finally the environmental and social impacts of mining together with the mitigation measures that the GCE examiner expects you to propose.

\courssub{Surface Mining Methods}

\begin{definition}[Surface mining]
\cle{Surface mining} (or open-cast mining) is the extraction of ore by removing the barren or low-grade material that covers it --- the \cle{overburden} --- and working the deposit directly from the surface, without underground workings. It is used when the ore body is \textbf{shallow}, \textbf{laterally extensive} and often of \textbf{low grade}, so that large volumes must be moved cheaply.
\end{definition}

There are four main types of surface working, and you must be able to distinguish them precisely.

\textbf{1. Open-pit mining.} The deposit is excavated as a large funnel-shaped pit, with \emph{benches} (terraced steps) cut into the walls for stability and access. A spiral haul road allows trucks to reach the bottom. Open pits suit massive or disseminated deposits such as porphyry copper or the iron ore bodies of the Mbalam--Nabeba region on the Cameroon--Congo border.

\textbf{2. Strip mining.} Where the ore forms a \emph{horizontal, tabular layer} close to the surface (a coal seam, a bauxite horizon), the overburden is stripped off in long parallel strips. The waste from one strip is dumped into the previous, already mined strip, which makes reclamation easier. This is the classic method for coal and for the bauxite plateaux of the Adamaoua (Minim--Martap, Ngaoundal).

\textbf{3. Quarrying.} A quarry is a surface working that extracts \emph{rock itself} as the product --- dimension stone (granite, marble), limestone for cement, sand and gravel for construction. Blasting is controlled so as not to shatter the blocks. Quarries around Yaound\'e and Douala supply aggregate for roads and buildings.

\textbf{4. Placer (alluvial) mining.} Dense, resistant minerals (gold, cassiterite, diamonds, rutile) are concentrated by rivers in \cle{placer deposits}. The miner simply washes the gravel: water separates the heavy minerals from the light sand by gravity, using pans, sluice boxes or dredges. This is the method used by artisanal gold miners along the rivers of East Cameroon (Lom and Dj\'erem, Kadey divisions) --- no blasting, no deep excavation, but heavy use of water and, too often, mercury.

\courssub{Underground Mining Methods}

\begin{definition}[Underground mining]
\cle{Underground mining} extracts ore through excavations made beneath the surface, leaving most of the overlying rock in place. It is used when the deposit is \textbf{too deep} for surface working, or when the overburden is too thick to remove economically. Key terms: a \cle{shaft} is a vertical access opening; an \cle{adit} is a horizontal or gently inclined access tunnel driven from a hillside; \cle{levels} are horizontal galleries at regular depth intervals; \cle{stopes} are the excavated rooms from which ore is actually removed.
\end{definition}

The three classical underground methods you must know are:

\begin{itemize}
\item \textbf{Room-and-pillar mining:} ore is removed from a network of ``rooms'', leaving columns of ore (the \emph{pillars}) in place to support the roof. Used for flat-lying, strong, tabular bodies (coal seams, salt, limestone). Simple and cheap, but up to half the ore may be left behind in the pillars.
\item \textbf{Cut-and-fill mining:} ore is removed in horizontal slices; each mined slice is backfilled with waste rock or cemented fill, which supports the walls while the next slice is taken. Used for steep, irregular, high-grade veins where selectivity matters (gold veins). More expensive, but almost all the ore is recovered and the surface does not subside.
\item \textbf{Longwall mining:} a mechanised shearer cuts along a long face of a coal seam while hydraulic supports hold up the roof; the supports advance and the roof behind is allowed to collapse in a controlled way. Very high productivity for thick, continuous coal seams, but it causes deliberate surface \emph{subsidence}.
\end{itemize}

\begin{popfigure}
\begin{center}
\begin{tikzpicture}[scale=0.92, every node/.style={font=\small}]
% ===== LEFT: open pit =====
\begin{scope}
\draw[fill=popgreenL] (-0.2,0) -- (5.4,0) -- (5.4,-0.4) -- (-0.2,-0.4) -- cycle;
\node[anchor=west] at (-0.2,0.25) {surface};
% pit with benches
\draw[fill=popgold!35,draw=popdark,thick]
 (0.6,0) -- (1.4,-1.0) -- (1.8,-1.0) -- (2.2,-1.9) -- (2.6,-1.9) -- (3.0,-2.7)
 -- (3.4,-2.7) -- (3.8,-1.9) -- (4.2,-1.9) -- (4.6,-1.0) -- (5.0,-1.0) -- (5.4,0) -- cycle;
\node[align=center] at (2.9,-1.35) {ore body};
\node[align=center,text=popdark] at (2.9,-3.15) {open pit (benches)};
% haul road arrow
\draw[->,thick,popblue] (0.9,-0.15) .. controls (1.6,-0.9) and (2.2,-1.7) .. (2.9,-2.5);
\node[text=popblue,anchor=west,align=center] at (0.4,-2.35) {haul\\road};
\node[align=center] at (2.7,0.8) {\textbf{Surface (open-pit) mine}};
\end{scope}
% ===== RIGHT: underground =====
\begin{scope}[xshift=7.2cm]
\draw[fill=popgreenL] (-0.2,0) -- (5.4,0) -- (5.4,-0.4) -- (-0.2,-0.4) -- cycle;
\node[anchor=west] at (-0.2,0.25) {surface};
% shaft
\draw[fill=gray!40,draw=black,thick] (1.1,0) rectangle (1.6,-3.4);
\node[align=center,anchor=east] at (1.05,-1.6) {shaft};
% headframe
\draw[thick] (1.0,0) -- (1.35,0.6) -- (1.7,0);
% levels
\draw[fill=gray!25,draw=black] (1.6,-1.3) rectangle (4.6,-1.05);
\draw[fill=gray!25,draw=black] (1.6,-2.4) rectangle (4.6,-2.15);
\node[anchor=west] at (4.65,-1.17) {level 1};
\node[anchor=west] at (4.65,-2.27) {level 2};
% ore vein / stope
\draw[fill=popgold!45,draw=popdark,thick] (3.0,-1.05) -- (3.5,-1.05) -- (3.9,-3.4) -- (3.4,-3.4) -- cycle;
\node[align=center,text=popdark] at (2.35,-3.1) {ore vein\\(stope)};
\node[align=center] at (2.7,0.8) {\textbf{Underground mine}};
\end{scope}
\end{tikzpicture}
\legende{Comparison cross-section: an open-pit mine removes all material above the ore (benched walls, haul road); an underground mine reaches the ore through a shaft and horizontal levels, leaving the overburden in place.}
\end{center}
\end{popfigure}

\courssub{Choosing the Method: Depth, Geometry, Grade and Economics}

No method is ``better'' in absolute terms; the choice is an \emph{engineering and economic decision}. The key criteria are:

\begin{itemize}
\item \textbf{Depth and geometry:} shallow, horizontal, extensive bodies $\Rightarrow$ surface mining; deep, steep, narrow veins $\Rightarrow$ underground mining.
\item \textbf{Grade:} low-grade ores need the low cost per tonne of surface methods; high-grade ores can pay for expensive underground work.
\item \textbf{Strength of ore and wall rocks:} weak, fractured ground makes underground stoping dangerous and favours open pits; very hard rock may favour underground selective mining.
\item \textbf{Economics:} the decisive quantity for surface mining is the \cle{stripping ratio}.
\end{itemize}

\begin{definition}[Stripping ratio]
The \cle{stripping ratio} (SR) of a surface mine is the volume (or mass) of waste overburden that must be removed to extract one unit of ore:
\[ \mathrm{SR} = \frac{\text{volume (or tonnes) of waste removed}}{\text{volume (or tonnes) of ore recovered}} \]
A deposit is worked at the surface only while the cost of stripping the waste remains lower than the extra cost of underground mining. The \emph{break-even stripping ratio} is the SR at which the two costs are equal; beyond it, the mine must go underground or close.
\end{definition}

\begin{methode}[Deciding between surface and underground mining]
\begin{enumerate}
\item \textbf{Locate the ore body:} from exploration data, determine its depth, thickness, dip and lateral extent.
\item \textbf{Compute the stripping ratio} for a hypothetical open pit: how many tonnes of waste per tonne of ore?
\item \textbf{Compare costs:} estimate the cost per tonne of ore by open pit (including stripping) and by underground mining (shaft sinking, support, ventilation, pumping).
\item \textbf{Check ground conditions:} rock strength, groundwater inflow, stability of pit walls or stopes.
\item \textbf{Decide:} if SR is below the break-even value and the body is shallow, choose surface mining; otherwise choose underground mining (or a combination: open pit first, underground at depth).
\end{enumerate}
\end{methode}

\begin{exoresolu}[Applying the stripping ratio]
An iron ore body near Mbalam lies under an average of $40\ \mathrm{m}$ of overburden. The ore horizon itself is $20\ \mathrm{m}$ thick and roughly horizontal. (a) Estimate the stripping ratio by volume. (b) The break-even stripping ratio for this project is $3:1$. Should the deposit be mined at the surface or underground? Justify.
\tcblower
\textbf{(a)} Over a given area $A$ of the deposit, the waste volume is $40A$ and the ore volume is $20A$. Hence
\[ \mathrm{SR} = \frac{40A}{20A} = 2 \quad \text{that is } 2:1. \]
\textbf{(b)} Since $\mathrm{SR} = 2:1 < 3:1$ (the break-even value), stripping the waste costs less than underground mining would. The deposit should therefore be mined by \textbf{surface (open-pit or strip) methods}. Note the reasoning pattern: compute the SR, compare it with the break-even SR, then conclude --- exactly what the examiner rewards.
\end{exoresolu}

\courssub{The Life Cycle of a Mine}

A mine is not a hole that appears overnight; it is a project with a planned life cycle, and the GCE syllabus expects you to know its stages:

\begin{enumerate}
\item \textbf{Development:} after exploration proves the deposit, the mine is built --- access roads, shaft sinking or box-cut for the pit, processing plant, waste dumps, water and power supply. This stage requires the largest investment and, in Cameroon, an \emph{environmental impact assessment} (EIA) before any permit is granted.
\item \textbf{Exploitation (production):} ore is extracted, transported and processed (crushing, concentration, smelting) over years or decades. Waste rock and tailings accumulate continuously.
\item \textbf{Closure:} when the ore is exhausted or uneconomic, extraction stops. Machinery is removed, shafts are sealed, and pit walls are made safe.
\item \textbf{Rehabilitation:} the site is restored --- regrading of dumps, replacement of topsoil, reforestation with local species, treatment of contaminated water --- so that the land can be used again (agriculture, forestry, sometimes a lake in the old pit).
\end{enumerate}

\begin{attention}[Closure is part of the plan]
A common weakness in exam answers is to treat rehabilitation as an optional afterthought. Modern mining law --- including the Cameroonian mining code --- requires closure and rehabilitation to be \emph{planned and financed from the start} of the project. A mine that is simply abandoned is a failure of management, not an inevitability.
\end{attention}

\courssub{Environmental Impacts of Mining}

Every mining method disturbs the environment, but in different ways. Learn the impacts \emph{by category}, and always link each impact to its cause.

\textbf{1. Deforestation and land degradation.} Surface mines, access roads and waste dumps destroy vegetation and fertile topsoil over large areas. In the forest zone of East Cameroon, both industrial sites and the scattering of artisanal gold diggings fragment the forest and leave abandoned pits that fill with stagnant water.

\textbf{2. Water pollution --- acid mine drainage.} This is the single most important chemical pollution problem of mining, and it is examinable as a property.

\begin{propriete}[Acid mine drainage]
\cle{Acid mine drainage} (AMD) forms when sulphide minerals --- especially pyrite, \ce{FeS2} --- are exposed to \textbf{oxygen and water} by excavation. The pyrite oxidises and produces sulphuric acid and dissolved iron:
\[ \ce{4FeS2 + 15O2 + 14H2O -> 4Fe(OH)3 + 8H2SO4} \]
The acidic water (pH sometimes below 3) dissolves heavy metals (Cu, Pb, Zn, Cd, As) from the waste rock and carries them into streams, killing aquatic life and making water unfit for use. AMD continues for decades \emph{after} a mine closes, which is why waste dumps must be covered and water treated.
\end{propriete}

\begin{popfigure}
\begin{center}
\begin{tikzpicture}[scale=1.0, every node/.style={font=\small}]
\node[draw,fill=popgold!40,rounded corners,align=center,text width=3.2cm] (pyr) at (0,0) {pyrite \ce{FeS2} in ore \& waste rock};
\node[draw,fill=popblueL,rounded corners,align=center,text width=2.6cm] (exp) at (4.3,0) {excavation exposes it to \ce{O2} + \ce{H2O}};
\node[draw,fill=poporange!35,rounded corners,align=center,text width=2.8cm] (ox) at (8.6,0) {oxidation $\to$ \ce{H2SO4} + \ce{Fe(OH)3}};
\node[draw,fill=poppink!35,rounded corners,align=center,text width=2.9cm] (leach) at (8.6,-2.2) {acid dissolves heavy metals (Cu, Pb, Zn, As)};
\node[draw,fill=popgreenL,rounded corners,align=center,text width=3.0cm] (river) at (4.3,-2.2) {acid, metal-rich water reaches streams \& groundwater};
\node[draw,fill=gray!25,rounded corners,align=center,text width=2.9cm] (harm) at (0,-2.2) {fish killed; water unfit; soils contaminated};
\draw[->,thick] (pyr) -- (exp);
\draw[->,thick] (exp) -- (ox);
\draw[->,thick] (ox) -- (leach);
\draw[->,thick] (leach) -- (river);
\draw[->,thick] (river) -- (harm);
\draw[->,thick,dashed,popdark] (harm.north) -- (pyr.south) node[midway,left,align=center,text width=2.4cm]{persists after closure};
\end{tikzpicture}
\legende{Formation of acid mine drainage from the oxidation of pyrite in exposed waste rock and abandoned workings.}
\end{center}
\end{popfigure}

\textbf{3. Mercury pollution from artisanal gold mining.} In the goldfields of East Cameroon (B\'etar\'e-Oya, Batouri, Ngoura), artisanal miners mix gold-bearing sediment with \cle{mercury}, which binds the gold into an amalgam; the amalgam is then heated, vaporising the mercury. Mercury is a potent neurotoxin: it contaminates rivers, accumulates in fish, and poisons the miners themselves and downstream communities. This is a flagship Cameroonian example --- quote it in your essays.

\textbf{4. Dust and noise.} Blasting, crushing and haulage release dust (including silica, dangerous to lungs) and continuous noise, affecting workers and nearby villages.

\textbf{5. Waste dumps and tailings.} Enormous volumes of waste rock and fine processing residues (\emph{tailings}) must be stored. Poorly built dumps can fail catastrophically, and tailings dams can release contaminated slurry into valleys.

\begin{center}
\begin{tabularx}{\linewidth}{|l|X|X|}
\hline
\rowcolor{popblueL} \textbf{Impact} & \textbf{Main cause} & \textbf{Typical mitigation} \\
\hline
Deforestation, land loss & clearing for pits, dumps, roads & regrading, topsoil replacement, reforestation \\
\hline
Acid mine drainage & pyrite oxidation in waste rock & cover dumps, divert water, lime treatment \\
\hline
Mercury pollution & artisanal gold amalgamation & retorts, mercury-free methods, regulation \\
\hline
Dust and noise & blasting, crushing, trucks & water spraying, buffers, working hours \\
\hline
Waste dumps / tailings & ore processing & engineered dams, monitoring, revegetation \\
\hline
\end{tabularx}
\end{center}

\courssub{Social Impacts, Conflicts and Mitigation}

Mining transforms societies as well as landscapes. \textbf{Industrial mining} is capital-intensive, employs relatively few but trained workers, and is (in principle) regulated and taxed. \textbf{Artisanal and small-scale mining} employs many more people --- tens of thousands in eastern Cameroon --- with simple tools, little capital and often no legal title. It provides vital rural income, but it is associated with child labour, dangerous pits, mercury use, smuggling and conflict over land.

Social impacts include: \emph{displacement of villages} and loss of farmland; \emph{in-migration} of job seekers, straining water, health and security; \emph{conflicts} between communities, artisanal miners and companies over land and benefits; and the paradox of resource-rich regions that remain poor when revenues are badly managed.

\begin{aretenir}[Mitigation measures to cite in the GCE]
\begin{itemize}
\item \textbf{Environmental Impact Assessment (EIA):} a compulsory study \emph{before} mining begins, identifying impacts and fixing mitigation and monitoring plans.
\item \textbf{Progressive rehabilitation:} regrade and revegetate worked-out areas \emph{during} exploitation, not only at closure; reforest with local species.
\item \textbf{Water treatment:} divert clean water away from waste, collect contaminated water, neutralise acidity with lime (\ce{CaCO3} / \ce{Ca(OH)2}) and settle out metals before discharge.
\item \textbf{Dust and noise control:} water spraying of roads, vegetation buffer zones, restricted blasting hours.
\item \textbf{Legal framework:} the \textbf{Cameroonian mining code} (Law No.\ 2016/017 of 14 December 2016) requires EIAs, rehabilitation plans and financial guarantees, and recognises artisanal mining through formal licensing --- a \emph{GCE-useful extra} that earns credit when quoted correctly.
\item \textbf{Community measures:} compensation for displaced people, local employment, formalisation and training of artisanal miners, mercury-free gold recovery (retorts, gravimetric concentrators).
\end{itemize}
\end{aretenir}

\begin{attention}[The classic exam trap: impacts without mitigation]
The most frequent mistake in GCE essays on this topic is to write a long, well-organised list of environmental impacts --- and stop there. A question such as ``Discuss the environmental consequences of mining in Cameroon'' implicitly demands \textbf{both} the impacts \textbf{and} the mitigation measures (EIA, rehabilitation, water treatment, legal control). Structure your answer in pairs: \emph{impact $\rightarrow$ cause $\rightarrow$ remedy}. An answer without mitigation rarely earns more than half marks.
\end{attention}

\begin{savaistu}[Mining in Cameroon: promise and paradox]
Cameroon is often described as a ``geological scandal'': it hosts world-class bauxite at Minim--Martap and Ngaoundal, iron ore at Mbalam, cobalt--nickel at Nkamouna, rutile at Akonolinga and widespread gold in the East --- yet large industrial mines have been slow to materialise, and most actual production remains artisanal gold and quarry materials. Understanding the methods, impacts and mitigation you have just studied is therefore not abstract: it is the key to how Cameroon can turn its geology into development without destroying its forests and rivers.
\end{savaistu}

\begin{exercice}[Choosing and justifying a mining method]
A cassiterite (tin) deposit occurs as (i) a deep, steeply dipping vein at $300\ \mathrm{m}$ depth, and (ii) a shallow alluvial gravel along a nearby river. For each occurrence, state the most appropriate mining method and justify your choice in two or three sentences. Then give two environmental risks specific to the alluvial operation and one mitigation measure for each.
\end{exercice}

\begin{corrige}[Exercise 1]
\textbf{(i) Deep vein:} underground mining, by shaft access with levels and a selective method such as cut-and-fill. Justification: at $300\ \mathrm{m}$ depth the stripping ratio for an open pit would be enormous (far above break-even), and the steep, narrow geometry suits selective stoping; backfill supports the walls and prevents subsidence.
\textbf{(ii) Alluvial gravel:} placer (alluvial) mining by washing and gravity separation (sluices, jigs). Justification: the tin mineral is already liberated and concentrated by the river at the surface, so no blasting or deep excavation is needed; gravity methods exploit cassiterite's high density.
\textbf{Risks and mitigation:} (a) \emph{River siltation and muddy water} from washing $\rightarrow$ build settling ponds and recycle process water; (b) \emph{destruction of riverbanks and abandoned pits} $\rightarrow$ regrade worked-out sections, replace topsoil and revegetate progressively. (Mercury is not normally used for cassiterite, but if gold is associated, mercury control would also be required.)
\end{corrige}

\begin{exercice}[Essay practice: impacts and mitigation]
``Artisanal gold mining is both a blessing and a curse for the East Region of Cameroon.'' Discuss this statement under the headings: (a) socio-economic benefits; (b) environmental impacts; (c) realistic mitigation measures. Your answer should contain at least three points under each heading.
\end{exercice}

\begin{corrige}[Exercise 2]
\textbf{(a) Benefits:} rural employment and income where few alternatives exist; cash circulating in local markets; a pathway for formal small-scale business if licensed; discovery of new prospects that may later attract industrial investment.
\textbf{(b) Impacts:} deforestation and abandoned water-filled pits; mercury contamination of rivers, fish and people; river siltation from washing; dust, noise and accidents in unstable diggings; social tensions, smuggling and child labour.
\textbf{(c) Mitigation:} formal licensing and organisation of miners into cooperatives; training in mercury-free recovery (retorts, gravimetric concentrators); progressive regrading, topsoil replacement and reforestation of worked-out sites; settling ponds for wash water; enforcement of the mining code and EIA requirements; fair marketing channels to reduce smuggling.
A strong GCE essay balances the two sides of the statement and ends with mitigation --- never with impacts alone.
\end{corrige}

\newpage

\coursec{Integration activity}

\courssub{Aims of the activity}

This integration activity brings together everything you have studied in this chapter: the economic uses of rocks, the engineering properties of geological materials, the siting of dams and reservoirs, slope stability, and geophysical exploration methods. It is the kind of multi-document problem you may meet in the GCE Advanced Level practical and structured papers. By the end of the activity you should be able to:

\begin{itemize}
\item read and interpret a simplified geological map, cross-section and borehole logs;
\item select and justify a dam axis using the strength, permeability and structure of the foundation rocks;
\item locate a reservoir basin with minimal risk of leakage;
\item choose a quarry site for aggregates and dimension stone, giving geological reasons;
\item recognise slopes at risk of failure and propose realistic stabilisation measures;
\item select an appropriate geophysical method to determine bedrock depth before construction;
\item defend a technical decision orally, as a site engineer would before a development committee.
\end{itemize}

\begin{aretenir}[Skills mobilised]
Map and cross-section reading; knowledge of rock properties (granite, basalt, shale); dam and reservoir geology; slope stability analysis; choice of exploration method; scientific communication.
\end{aretenir}

\courssub{Situation}

The town of \textbf{Nkwen--Mankon Heights}, a fast-growing settlement on the outskirts of Bamenda (North-West Region), suffers from severe water shortage every dry season and lacks a local source of construction material. The council has hired your team of junior engineering geologists to study the \textbf{Mezam Valley}, \SI{8}{km} north of the town, where a small concrete-gravity dam (height \SI{25}{m}), a reservoir and a quarry are planned, together with a \SI{6}{km} access road.

Your team receives the following documents.

\textbf{1. Simplified geological map and cross-section of the Mezam Valley.} The valley trends roughly east--west. From north to south, the map shows four units:

\begin{center}
\begin{tabularx}{\linewidth}{|l|X|X|}\hline
\rowcolor{popblueL} \textbf{Unit} & \textbf{Description} & \textbf{Position} \\ \hline
A -- Granite & Massive, coarse-grained, very hard, widely jointed but joints tight; forms the northern valley wall and a narrow gorge at site G1 & Northern flank, crosses the valley at G1 \\ \hline
B -- Fractured basalt & Hard but intensely fractured and weathered to \SI{6}{m} depth; highly permeable through open joints & Central valley floor and southern flank around site G2 \\ \hline
C -- Shale & Soft, fissile, impermeable, weathers to slippery clay; dips $30^{\circ}$ towards the valley (south) & Southern valley wall, above the basalt \\ \hline
F -- Fault zone & Crushed rock over \SI{15}{m} width, trending N--S, cuts all units near site G3 & Eastern end of the valley \\ \hline
\end{tabularx}
\end{center}

Three possible dam axes are marked: \textbf{G1} (granite gorge, valley width \SI{80}{m}), \textbf{G2} (fractured basalt, valley width \SI{150}{m}), and \textbf{G3} (on the fault zone, valley width \SI{120}{m}).

\begin{popfigure}
\begin{tikzpicture}[scale=1.0, font=\small]
% granite north
\fill[popblueL] (0,2.2) rectangle (10,4.2);
\draw[popblue] (0,2.2) -- (10,2.2);
\node[popdark] at (2.2,3.6) {A -- Granite (northern wall)};
% basalt centre
\fill[popgreenL] (0,0.6) rectangle (10,2.2);
\node[popdark] at (2.2,1.4) {B -- Fractured basalt (valley floor)};
% shale south
\fill[popgold!25] (0,-1.6) rectangle (10,0.6);
\draw[popgold] (0,0.6) -- (10,0.6);
\node[popdark] at (2.2,-0.5) {C -- Shale, dip $30^{\circ}$ S};
% fault zone
\fill[poporange!35] (7.6,-1.6) rectangle (8.2,4.2);
\draw[poporange, thick] (7.6,-1.6) -- (7.6,4.2);
\draw[poporange, thick] (8.2,-1.6) -- (8.2,4.2);
\node[popdark, rotate=90] at (7.9,3.4) {F -- Fault zone};
% axes
\draw[popink, very thick, dashed] (2.0,2.2) -- (2.0,0.6);
\node[popink, above] at (2.0,2.2) {G1};
\draw[popink, very thick, dashed] (5.0,2.2) -- (5.0,0.6);
\node[popink, above] at (5.0,2.2) {G2};
\draw[popink, very thick, dashed] (7.9,2.2) -- (7.9,0.6);
\node[popink, above] at (7.9,2.2) {G3};
% river
\draw[popblue, thick, ->] (0.3,1.35) -- (9.7,1.35);
\node[popblue, below] at (9.3,1.35) {river};
% reservoir arrow
\draw[popdark, thick, ->] (2.0,1.9) -- (0.6,1.9);
\node[popdark, align=center, text width=2.6cm] at (1.2,2.6) {reservoir extends west};
% slopes
\node[popdark] at (5.6,-1.2) {S1 (shale slope)};
\draw[popdark, ->] (5.6,-1.0) -- (5.6,-0.2);
\node[popdark] at (4.4,3.0) {S2 (road cut)};
\draw[popdark, ->] (4.4,2.85) -- (4.4,2.3);
% road
\draw[black, thick, loosely dotted] (0,2.0) -- (10,2.0);
\node[black, above right] at (8.4,2.0) {access road};
% north arrow
\draw[->, thick] (9.6,3.4) -- (9.6,4.0);
\node[above] at (9.6,4.0) {N};
\end{tikzpicture}
\legende{Simplified sketch map of the Mezam Valley (not to scale): four geological units, three candidate dam axes G1--G3, the fault zone F, and the two slopes at risk S1 and S2.}
\end{popfigure}

\textbf{2. Borehole logs} (simplified) drilled along the three axes:

\begin{center}
\begin{tabularx}{\linewidth}{|l|X|X|X|}\hline
\rowcolor{popblueL} \textbf{Borehole} & \textbf{0--3 m} & \textbf{3--8 m} & \textbf{Below 8 m} \\ \hline
BH1 (axis G1) & Topsoil and laterite & Slightly weathered granite & Fresh, massive granite \\ \hline
BH2 (axis G2) & Topsoil & Highly fractured, weathered basalt (open joints, clay filling) & Fractured basalt, water loss during drilling \\ \hline
BH3 (axis G3) & Topsoil & Crushed rock and fault gouge (clay) & Broken rock, artesian water rise \\ \hline
\end{tabularx}
\end{center}

\begin{attention}[Reading borehole logs]
\emph{Water loss during drilling} (BH2) means the drilling fluid escapes into open joints: direct evidence of high \cle{permeability}. An \emph{artesian water rise} (BH3) means groundwater under pressure rises up the borehole: evidence of a confined, fractured aquifer that would complicate excavation and grouting.
\end{attention}

\textbf{3. Additional data.} The water table lies at \SI{2}{m} depth on the valley floor. The shale slope above the proposed access road (slope S1, southern wall) already shows small landslide scars after the rainy season. A second slope, S2, is a steep cut planned through weathered basalt on the northern side of the road. The reservoir basin would extend upstream (west) of the chosen axis, mostly over granite and shale; the fault zone crosses its eastern edge. Aggregates (crushed stone) and dimension stone are needed in large quantity; haulage distance to the town should be short.

\courssub{Tasks}

\begin{enumerate}
\item \textbf{Choosing the dam axis.} For each of the three candidate axes (G1, G2, G3), state the foundation rock, its strength and permeability, and one geological advantage and one geological disadvantage. Conclude by choosing the best axis and justify your choice in three or four sentences.
\item \textbf{Locating the reservoir.} Explain why the reservoir should extend westwards over the granite and shale rather than eastwards across the fault zone. Name the property of shale that reduces leakage, and explain why the fault zone and the fractured basalt are dangerous for water tightness.
\item \textbf{Selecting the quarry.} Propose a site for (a) aggregates for concrete and (b) dimension stone. Which rock unit do you choose for each, and why? Give two reasons why the shale is unsuitable, and one practical reason for placing the quarry close to the access road.
\item \textbf{Slope stability along the access road.} Identify the two slopes at risk (S1 and S2). For each one, name the likely type of failure, the geological factor responsible (think of dip direction, fissility, weathering, water), and propose two stabilisation measures for each slope.
\item \textbf{Checking bedrock depth.} Before construction, the engineers want to map the depth of fresh bedrock along the dam axis and the road. Which geophysical method do you recommend, and why? Name one alternative method and state its principle in one sentence.
\item \textbf{Presentation.} Assemble your answers into a one-page justified site plan (sketch map plus short notes) and be ready to defend it in a class debate, playing the roles of geologist, engineer and a worried village representative.
\end{enumerate}

\courssub{Model answer}

\textbf{1. Choosing the dam axis.}

\begin{center}
\begin{tabularx}{\linewidth}{|l|X|X|X|}\hline
\rowcolor{popblueL} \textbf{Axis} & \textbf{Foundation} & \textbf{Advantage} & \textbf{Disadvantage} \\ \hline
G1 & Fresh massive granite below \SI{3}{m} (BH1) & Very strong, impermeable, narrow gorge (\SI{80}{m}) so a short, cheap dam & Joints must be cleaned and grouted \\ \hline
G2 & Fractured, weathered basalt & Hard rock in bulk & Open joints, water loss in BH2: high permeability, weak weathered zone, wide valley (\SI{150}{m}) \\ \hline
G3 & Fault zone: crushed rock and gouge & Narrower than G2 & Crushed, weak, permeable ground; artesian water; risk of leakage and sliding along the fault \\ \hline
\end{tabularx}
\end{center}

\textbf{Chosen axis: G1.} The granite at G1 is a \cle{fresh}, \cle{massive} rock of very high compressive strength, able to carry the weight of a concrete dam without settlement. Its tight joints make it almost impermeable, so the foundation will not leak. The narrow gorge reduces the volume of concrete and therefore the cost. G2 is rejected because the fractured basalt would leak and would need deep excavation and grouting; G3 is rejected because a fault zone is the worst possible foundation: weak, crushed, permeable and potentially unstable.

\textbf{2. Locating the reservoir.} The reservoir must extend \emph{westwards} over the granite and shale. Granite is impermeable when fresh, and shale, although soft, has the key property of being \cle{impermeable} (its clay-sized particles leave no connected pores), so it forms a natural seal on the valley sides and floor. Eastwards, the fault zone is a band of crushed, broken rock that acts as a \cle{drainage path}: water would escape along it out of the basin. Likewise the fractured basalt of the centre of the valley has open joints through which water would leak. Keeping the reservoir away from both features minimises leakage risk.

\textbf{3. Selecting the quarry.} (a) \emph{Aggregates:} the \textbf{granite} of the northern wall near G1 is chosen, because it is hard, tough and resistant to crushing, giving excellent crushed stone for concrete. The unweathered basalt could also serve, but its fractures make blasting less controlled. (b) \emph{Dimension stone:} the same massive granite, where joints are widely spaced, can be cut into regular blocks for building stone, as is done in many quarries of the North-West. The shale is unsuitable because (i) it is soft and weak, so it crushes to poor-quality flaky fragments, and (ii) it is fissile and weathers rapidly to clay, so blocks and aggregates would disintegrate. The quarry is placed beside the access road to shorten haulage and reduce transport cost.

\textbf{4. Slope stability along the access road.}

\emph{Slope S1 (shale, southern wall).} Likely failure: a \cle{planar slide} or mudslide. The shale beds dip $30^{\circ}$ towards the valley, i.e.\ in the same direction as the slope face, so bedding planes are daylighted and act as sliding surfaces; the shale is fissile and weathers to slippery clay, and rain water lubricates the planes. Stabilisation: (i) cut the slope back to a gentler angle than the dip (or bench it), and (ii) install surface drains and a lined ditch at the crest to keep water out; a retaining wall at the toe can be added.

\emph{Slope S2 (weathered basalt, northern side).} Likely failure: \cle{rock falls} and small wedge failures of loose blocks. The basalt is intensely fractured and weathered to \SI{6}{m}, so blocks are detached along joints, especially when the road cut is steepened. Stabilisation: (i) scale off loose blocks and cover the face with wire mesh (or bolt unstable blocks), and (ii) reduce the cut slope angle and provide a drainage ditch plus a catch ditch at the toe to trap falling debris.

\textbf{5. Checking bedrock depth.} The recommended method is \cle{seismic refraction}: the contrast in seismic wave velocity between the slow weathered layer (topsoil, weathered rock) and the fast fresh bedrock (granite or basalt) allows the depth to bedrock to be calculated along a profile, exactly what is needed along the dam axis and the road. It is cheap, rapid and well suited to layered ground with velocity increasing with depth. Alternative: \cle{electrical resistivity surveying}, in which an electric current is injected into the ground and the measured resistivity distinguishes conductive weathered, water-bearing material from resistive fresh bedrock, giving the thickness of the weathered zone.

\textbf{6. Presentation.} The final site plan shows: the dam axis at G1 across the granite gorge; the reservoir extending west over granite and shale, stopping short of the fault zone; the quarry in the massive granite near the road; slopes S1 and S2 marked with their stabilisation works; and the seismic refraction lines along the axis and road. In the debate, the geologist defends the geological choices, the engineer discusses cost and safety, and the village representative questions water security and road safety.

\begin{attention}[Common mistakes to avoid]
Do not choose an axis only because the valley is narrow: rock quality comes first. Do not confuse ``hard'' with ``impermeable'': fractured basalt is hard but leaks. A bed dipping \emph{towards} the valley is dangerous; a bed dipping \emph{into} the slope is stable. Finally, never site a dam or reservoir across an active-looking fault zone.
\end{attention}

\begin{exercice}[Going further]
A second valley, \SI{12}{km} away, offers a single possible dam site on \emph{slightly fractured sandstone} overlying shale, with beds dipping $15^{\circ}$ \emph{upstream} (into the northern slope). (a) State one advantage and one risk of this site compared with G1. (b) Explain why the upstream dip of the beds is favourable for slope stability but must still be checked for leakage along bedding planes. (c) Suggest one geophysical method to verify the thickness of the sandstone before design.
\end{exercice}

\begin{corrige}[Exercise 1]
(a) Advantage: sandstone is generally strong and, where cemented and tight, a fair foundation; the shale below adds an impermeable floor. Risk: sandstone is porous, and fractures plus bedding planes may allow leakage under the dam, unlike the massive granite of G1. (b) Beds dipping upstream (into the slope) are not daylighted on the valley-side faces, so no sliding surface is exposed: the slope is stable. However, water under reservoir pressure can still travel along bedding planes and fractures beneath the dam, so grouting or a cut-off trench may be needed. (c) Seismic refraction (or electrical resistivity) profiling across the axis gives the depth to fresh, unweathered sandstone and to the shale contact.
\end{corrige}

\newpage

\coursec{Methods and worked examples}

\coursec{Economic Geology}

\courssub{Skill 1: Determining Coal Rank and Interpreting a Coalification Sequence}

\begin{definition}[Coal and Coalification]
\cle{Coal} is a combustible sedimentary rock formed from the accumulation and preservation of plant debris in swampy environments, subsequently buried and altered by heat and pressure. \cle{Coalification} (or coal maturation) is the diagenetic to low-grade metamorphic process that transforms peat into successively higher ranks of coal through devolatilisation (loss of H, O as \ce{H2O}, \ce{CO2}, \ce{CH4}) and carbon enrichment.
\end{definition}

\begin{propriete}[Coal Rank Series and Key Properties]
The standard rank series (increasing coalification) is:
\cle{Peat} $\rightarrow$ \cle{Lignite} (Brown Coal) $\rightarrow$ \cle{Sub-bituminous Coal} $\rightarrow$ \cle{Bituminous Coal} $\rightarrow$ \cle{Anthracite} $\rightarrow$ \cle{Meta-anthracite} (graphite transition).
Key trends with increasing rank:
\begin{itemize}
\item \textbf{Fixed Carbon} increases ($\sim$50\% $\rightarrow$ $>$90\%).
\item \textbf{Moisture} (inherent) decreases ($\sim$40\% $\rightarrow$ $<$5\%).
\item \textbf{Volatile Matter} decreases ($\sim$45\% $\rightarrow$ $<$10\%).
\item \textbf{Calorific Value} increases ($\sim$15 $\rightarrow$ $>$35 MJ/kg).
\item \textbf{Reflectance ($R_o$)} of vitrinite increases (standard rank parameter).
\item \textbf{Hardness, lustre, brittleness} increase.
\end{itemize}
\end{propriete}

\begin{methode}[Analysing a coal rank problem]
\begin{enumerate}
\item \textbf{Recall the rank series} in order: Peat $\rightarrow$ Lignite $\rightarrow$ Sub-bituminous $\rightarrow$ Bituminous $\rightarrow$ Anthracite.
\item \textbf{Identify controlling factors}: increasing depth of burial, temperature (geothermal gradient $\sim$30°C/km), pressure, and time. Tectonic compression or igneous intrusions (contact metamorphism) can locally accelerate rank increase.
\item \textbf{Track property trends}: Fixed C $\uparrow$, Moisture $\downarrow$, Volatiles $\downarrow$, Calorific Value $\uparrow$, Lustre/Hardness $\uparrow$.
\item \textbf{Assign rank} by matching given data (proximate analysis: moisture, volatile matter, fixed carbon, ash; calorific value) to the trends. Justify explicitly (e.g. ``Sample A has 45\% fixed C and 35\% moisture $\rightarrow$ Lignite'').
\item \textbf{Link to geology}: Higher rank implies deeper burial, tectonic shortening, or proximity to a heat source. If coals of same age differ in rank, invoke tectonic or magmatic cause.
\end{enumerate}
\textbf{Reflex}: Coalification is primarily a \emph{temperature-time} process (Arrhenius kinetics). Vitrinite reflectance ($R_o$) is the most reliable rank indicator in the lab.
\end{methode}

\begin{exemplebox}[Proximate Analysis and Rank Boundaries (Typical Values)]
\begin{center}
\begin{tabularx}{\linewidth}{|l|X|X|X|X|X|}
\hline
\rowcolor{popblueL} \textbf{Rank} & \textbf{Fixed C (\%)} & \textbf{Moisture (\%)} & \textbf{Volatile Matter (\%)} & \textbf{Calorific Value (MJ/kg)} & \textbf{Vitrinite Reflectance $R_o$ (\%)} \\
\hline
Peat & $<$50 & $>$50 & $>$50 & $<$15 & $<$0.3 \\
\hline
Lignite & 50--65 & 30--45 & 40--50 & 15--20 & 0.3--0.5 \\
\hline
Sub-bituminous & 65--75 & 15--30 & 35--45 & 20--26 & 0.5--0.65 \\
\hline
Bituminous & 75--90 & 5--15 & 15--35 & 26--35 & 0.65--2.0 \\
\hline
Anthracite & $>$90 & $<$5 & $<$10 & $>$35 & $>$2.0 \\
\hline
\end{tabularx}
\end{center}
\end{exemplebox}

\begin{exoresolu}[Comparing three coal samples]
Three coal samples X, Y and Z were collected from different coalfields. Their proximate analyses give:

\begin{center}
\begin{tabularx}{\linewidth}{|l|X|X|X|}
\hline
\rowcolor{popblueL} \textbf{Property} & \textbf{X} & \textbf{Y} & \textbf{Z} \\
\hline
Fixed carbon (\%) & 45 & 78 & 92 \\
\hline
Moisture (\%) & 35 & 8 & 3 \\
\hline
Volatile matter (\%) & 40 & 18 & 6 \\
\hline
Calorific value (MJ/kg) & 17 & 29 & 34 \\
\hline
\end{tabularx}
\end{center}

\begin{enumerate}
\item Assign a rank (lignite, sub-bituminous, bituminous, anthracite) to each sample, justifying your answer.
\item State and explain the geological conditions that transform sample X into sample Z.
\item Explain why anthracite is economically more valuable than lignite.
\end{enumerate}
\tcblower
\textbf{1. Rank assignment.}

Coalification expels water and volatiles, concentrating carbon. We rank by fixed carbon and calorific value:
\begin{itemize}
\item \textbf{Sample X}: 45\% fixed C, very high moisture (35\%) and volatile matter (40\%), low CV (17 MJ/kg). Matches \cle{lignite} (brown coal).
\item \textbf{Sample Y}: 78\% fixed C, 8\% moisture, 18\% volatiles, CV 29 MJ/kg. Falls squarely in \cle{bituminous coal} range (high-volatile bituminous).
\item \textbf{Sample Z}: 92\% fixed C, 3\% moisture, 6\% volatiles, CV 34 MJ/kg. Highest rank: \cle{anthracite} (typically bright, conchoidal fracture, metallic lustre).
\end{itemize}
(No sample fits sub-bituminous; the gap between X and Y spans that rank.)

\textbf{2. Conditions transforming X into Z.}

The transformation is \emph{coalification}, driven by:
\begin{itemize}
\item \textbf{Burial}: accumulation of overlying sediments increases lithostatic pressure and temperature (geothermal gradient $\sim$30°C/km). Entry into the ``oil window'' ($\sim$60--120°C) corresponds to high-volatile bituminous; anthracite requires $>$200--250°C (anchizone/low-grade metamorphism).
\item \textbf{Time}: reactions are kinetically slow; millions of years at moderate T, or shorter time at higher T.
\item \textbf{Tectonic compression} (orogenic belts) or \textbf{igneous heat} (intrusions) can accelerate the process locally (contact metamorphism of coal $\rightarrow$ natural coke).
\end{itemize}
Progressive devolatilisation: \ce{H2O} (peat $\rightarrow$ lignite), \ce{CO2} (lignite $\rightarrow$ sub-bituminous), \ce{CH4} + heavier hydrocarbons (bituminous $\rightarrow$ anthracite).

\textbf{3. Economic value.}

Anthracite yields more heat per kilogram (34 vs 17 MJ/kg), burns with little smoke (low volatiles $\rightarrow$ domestic/industrial preference), and contains fewer impurities (less ash, sulphur) to dispose of. Transport costs per unit of energy are lower and combustion efficiency higher, so anthracite commands a premium price. Lignite is used near mine-mouth power stations due to high moisture and low energy density.
\end{exoresolu}

\begin{savaistu}[Coal in Cameroon]
Cameroon has no commercial coal production. Small occurrences of lignitic coal exist in the \emph{Mamfe Basin} (Cretaceous) and near \emph{Ngaoundéré} (Tertiary), but they are thin, high in sulphur/ash, and not economically viable. The country relies on hydroelectricity, imported petroleum products, and increasingly on natural gas (e.g., Kribi gas-to-power). World-class coal deposits are found in the Karoo basins of Southern Africa (South Africa, Botswana, Mozambique) and in the Gondwana basins of India and Australia.
\end{savaistu}

\courssub{Skill 2: Analysing a Petroleum System and Recognising a Trap}

\begin{definition}[The Petroleum System]
A \cle{petroleum system} comprises five essential elements that must coexist in space and time:
\begin{enumerate}
\item \textbf{Source Rock}: organic-rich sediment (shale, marl, carbonate) with TOC $>$0.5\%, deposited in anoxic conditions. Generates hydrocarbons in the \emph{oil window} (60--120°C, $\sim$2--4 km depth) and \emph{gas window} ($>$120°C).
\item \textbf{Reservoir Rock}: porous \emph{and} permeable rock (sandstone, fractured limestone, dolomite) to store and transmit fluids. Porosity $\phi$ (storage), Permeability $k$ (flow).
\item \textbf{Cap Rock / Seal}: impermeable rock (shale, evaporite, tight carbonate) overlying the reservoir, preventing vertical escape. Capillary entry pressure must exceed buoyancy pressure of hydrocarbon column.
\item \textbf{Trap}: geometric configuration that concentrates hydrocarbons. \emph{Structural} (anticline, fault trap, salt dome) or \emph{Stratigraphic} (pinch-out, unconformity, reef/lens, diagenetic).
\item \textbf{Migration Pathway}: permeable carrier beds or faults connecting source to trap. \emph{Primary migration} (source $\rightarrow$ carrier), \emph{Secondary migration} (carrier $\rightarrow$ trap).
\end{enumerate}
\emph{Critical moment}: trap must exist \emph{before} or \emph{during} migration; seal must be intact.
\end{definition}

\begin{methode}[Solving a petroleum-system exercise]
\begin{enumerate}
\item \textbf{List the five elements} and the physical criteria for each (lithology, porosity/permeability, geometry).
\item \textbf{On a cross-section}, identify each element from lithologies: shale = source or seal; sandstone/limestone = reservoir; fold/fault = structural trap; pinch-out = stratigraphic trap.
\item \textbf{Classify the trap} (structural vs stratigraphic) and name its type.
\item \textbf{Explain fluid arrangement}: Gas (lightest) on top, Oil in middle, Water (heaviest) below $\rightarrow$ \cle{Gas-Oil Contact (GOC)} and \cle{Oil-Water Contact (OWC)} are horizontal surfaces in the reservoir.
\item \textbf{Assess prospectivity}: All five elements present? Seal effective before migration? Source rock mature? Reservoir quality adequate?
\end{enumerate}
\textbf{Reflex}: No cap rock $\Rightarrow$ no accumulation, however good the reservoir. A ``bright spot'' on seismic suggests gas but is not proof of commercial oil.
\end{methode}

\begin{exoresolu}[Interpreting a petroleum trap]
The cross-section below shows an anticline cored by shale, overlain by a porous sandstone bed, itself overlain by a thick shale. A well drilled on the flank of the anticline found only water.

\begin{popfigure}
\begin{tikzpicture}[scale=1.1]
% Source shale (bottom)
\fill[popgreenL] (0,-2.2) rectangle (7,-3.0);
% Reservoir sandstone (middle, folded)
\fill[popblueL] (0,-0.6) .. controls (1.5,0.2) and (3.5,0.2) .. (5,-0.6) -- (5,-1.4) .. controls (3.5,-0.6) and (1.5,-0.6) .. (0,-1.4) -- cycle;
% Cap shale (top)
\fill[gray!40] (0,0.8) .. controls (1.5,1.6) and (3.5,1.6) .. (5,0.8) -- (5,-0.6) .. controls (3.5,0.2) and (1.5,0.2) .. (0,-0.6) -- cycle;
% Core shale (anticline core)
\fill[popgreenL] (0,-1.4) .. controls (1.5,-0.6) and (3.5,-0.6) .. (5,-1.4) -- (5,-2.2) -- (0,-2.2) -- cycle;
% Contacts
\draw[dashed, thick, popdark] (0.5,-0.9) -- (4.5,-0.9) node[midway, above, popdark] {\small OWC};
\draw[dashed, thick, popdark] (1.0,-0.75) -- (4.0,-0.75) node[midway, above, popdark] {\small GOC};
% Fluids in reservoir
\fill[popgold] (1.0,-0.75) .. controls (2.5,-0.5) and (3.5,-0.5) .. (4.0,-0.75) -- (4.0,-0.9) .. controls (3.5,-0.65) and (2.5,-0.65) .. (1.0,-0.9) -- cycle;
\fill[poporange] (0.5,-0.9) .. controls (2.5,-0.65) and (3.5,-0.65) .. (4.5,-0.9) -- (4.5,-1.4) .. controls (3.5,-1.15) and (1.5,-1.15) .. (0.5,-1.4) -- cycle;
% Labels
\node at (2.5,1.1) {\small Cap Rock: Shale (seal)};
\node[text width=2.8cm,align=center] at (2.5,-0.3) {\small Reservoir: Sandstone};
\node at (2.5,-1.05) {\small Gas};
\node at (2.5,-1.2) {\small Oil};
\node at (2.5,-1.7) {\small Water};
\node at (2.5,-2.6) {\small Source Rock: Organic Shale};
% Migration arrows
\draw[->,thick,popdark] (1.0,-2.5) -- (1.5,-1.6);
\draw[->,thick,popdark] (4.0,-2.5) -- (3.5,-1.6);
\node[popdark] at (0.6,-2.7) {\small Migration};
% Dry well on flank
\draw[thick, popdark] (5.5,-0.2) -- (5.5,-1.8);
\node[popdark, rotate=90] at (5.7,-1.0) {\small Dry well (water only)};
\end{tikzpicture}
\legende{Anticlinal trap: gas (orange), oil (gold), water (blue) in sandstone reservoir sealed by shale. Source rock in core and below.}
\end{popfigure}

\begin{enumerate}
\item Name the type of trap and classify it (structural or stratigraphic).
\item Identify the source rock, reservoir rock and cap rock, justifying each choice.
\item Explain why the flank well found only water, and state where a second well should be drilled.
\end{enumerate}
\tcblower
\textbf{1. Type of trap.} The reservoir sandstone is folded into an upwarp (anticline) and closed by an overlying impermeable shale. This is an \cle{anticlinal trap}, a \emph{structural trap} because it results from tectonic deformation (folding) of the strata after deposition.

\textbf{2. Elements of the system.}
\begin{itemize}
\item \textbf{Source rock}: the organic-rich shale at the core and base of the anticline. Shales deposited in quiet, oxygen-poor (anoxic) marine or lacustrine environments preserve organic matter (kerogen). Burial to 2--4 km brings it into the oil window (60--120°C), generating oil and gas.
\item \textbf{Reservoir rock}: the sandstone. Its intergranular \emph{porosity} (15--25\%) stores fluids and its \emph{permeability} (interconnected pores) allows flow to a wellbore.
\item \textbf{Cap rock}: the thick shale overlying the sandstone. Its very low permeability (capillary entry pressure $>$ buoyancy pressure) prevents upward escape of buoyant hydrocarbons.
\end{itemize}

\textbf{3. The dry flank well.} Hydrocarbons are less dense than formation water ($\rho_{\text{gas}} < \rho_{\text{oil}} < \rho_{\text{water}}$). Within the permeable reservoir, they migrate \emph{up-dip} and accumulate at the \emph{crest} of the anticline, floating on water. The oil--water contact (OWC) and gas--oil contact (GOC) are horizontal. A well on the flank penetrates the reservoir \emph{below the OWC}, hence produces only water. The second well should be drilled at the \textbf{crest (structural culmination) of the anticline}, where the hydrocarbon column is thickest. This logic underpins exploration in Cameroon's \emph{Rio del Rey} and \emph{Douala/Kribi-Campo} basins, where Tertiary sandstone reservoirs are trapped in rollover anticlines and fault blocks sealed by marine shales.
\end{exoresolu}

\begin{attention}[Common Trap Misidentification]
\begin{itemize}
\item Do not confuse a \emph{syncline} (downfold) with an anticline trap $\rightarrow$ hydrocarbons escape up the limbs unless a stratigraphic pinch-out closes the syncline.
\item A \emph{fault trap} requires the fault to juxtapose reservoir against seal (shale gouge ratio). If reservoir is against reservoir across the fault, it leaks.
\item \emph{Stratigraphic traps} (e.g., reef buildup, channel sand pinch-out) have no structural closure; they rely on lateral facies change. Seismic recognition is harder.
\end{itemize}
\end{attention}

\courssub{Skill 3: Estimating Reserves and the Life of a Deposit}

\begin{definition}[Resources vs Reserves (McKelvey Classification)]
\begin{itemize}
\item \textbf{Mineral Resources}: total concentration of mineral in the ground with reasonable prospects for economic extraction. Sub-divided: \emph{Inferred} (low confidence), \emph{Indicated} (moderate), \emph{Measured} (high confidence).
\item \textbf{Mineral Reserves}: the economically mineable part of a \emph{Measured} or \emph{Indicated} Resource. Includes diluting materials and losses. Sub-divided: \emph{Probable Reserve}, \emph{Proven Reserve}.
\item \textbf{Only Reserves} enter mine-life calculations. Resources require further drilling/study to convert.
\end{itemize}
\end{definition}

\begin{methode}[Reserve and mine-life calculations]
\begin{enumerate}
\item \textbf{Volume of ore body}: $V = \text{Area} \times \text{Average Thickness}$ (convert all to metres: $1\ \mathrm{km^2} = 10^6\ \mathrm{m^2}$).
\item \textbf{Tonnage (in-place)}: $M = V \times \rho_b$, with $\rho_b$ = bulk density (t/m$^3$).
\item \textbf{Metal/Commodity content}: $M_{\text{commodity}} = M \times \text{Grade}$ (grade as decimal, e.g. 42\% = 0.42).
\item \textbf{Recoverable Reserve}: $M_{\text{rec}} = M \times \text{Recovery Factor}$ (mining recovery $\times$ processing recovery). Typical: open pit 85--95\%, underground 60--85\%, oil 20--60\%.
\item \textbf{Life of Mine (LoM)}: $t = \dfrac{M_{\text{rec}}}{\text{Annual Production Rate}}$.
\item \textbf{Check units} at every step; express final answer with appropriate significant figures (usually 2--3).
\end{enumerate}
\textbf{Reflex}: Stripping ratio (waste:ore) determines open-pit depth limit. If stripping ratio exceeds economic threshold, switch to underground or abandon.
\end{methode}

\begin{exoresolu}[Life of a bauxite deposit]
A lateritic bauxite plateau near Minim-Martap (Adamaoua, Cameroon) covers an area of $20\ \mathrm{km^2}$ with an average ore thickness of $8\ \mathrm{m}$. The bulk density of the bauxite is $2.2\ \mathrm{t/m^3}$ and the average grade is $42\%$ \ce{Al2O3}. Mining recovery is $90\%$ and the planned plant treats $5$ million tonnes of ore per year.
\begin{enumerate}
\item Calculate the total tonnage of ore in place.
\item Calculate the recoverable tonnage and the tonnage of contained \ce{Al2O3}.
\item Estimate the life of the mine.
\end{enumerate}
\tcblower
\textbf{1. Tonnage in place.}

Area: $20\ \mathrm{km^2} = 20 \times 10^{6}\ \mathrm{m^2}$.
Volume: $V = 20\times 10^{6}\ \mathrm{m^2} \times 8\ \mathrm{m} = 1.6\times 10^{8}\ \mathrm{m^3}$.
Tonnage: $M = V \times \rho_b = 1.6\times 10^{8} \times 2.2 = 3.52\times 10^{8}\ \mathrm{t} = 352\ \text{million tonnes (Mt)}$.

\textbf{2. Recoverable tonnage and alumina content.}

Recoverable ore (mining recovery 90\%):
$M_{\text{rec}} = 352 \times 0.90 = 316.8\ \text{Mt} \approx 317\ \text{Mt}$.
Contained \ce{Al2O3} (on in-place ore, grade 42\%):
$M_{\ce{Al2O3}} = 352 \times 0.42 = 147.8\ \text{Mt} \approx 148\ \text{Mt of } \ce{Al2O3}$.
(If processing recovery is given, multiply again; here only mining recovery stated.)

\textbf{3. Life of the mine.}
$t = \dfrac{316.8\ \text{Mt}}{5\ \text{Mt/year}} = 63.36\ \text{years} \approx 63\ \text{years}$.

\textbf{Conclusion}: the deposit could supply the plant for roughly six decades. The Adamaoua bauxite plateaux (Minim-Martap, Ngaoundal) are a strategic resource for Cameroon's aluminium ambitions. \emph{Exam reflex}: apply recovery factor \emph{before} dividing by annual production.
\end{exoresolu}

\courssub{Skill 4: Assessing a Dam or Tunnel Site from a Geological Cross-Section}

\begin{propriete}[Engineering Rock Mass Quality]
The behaviour of a rock mass is controlled by \emph{discontinuities} (joints, faults, bedding planes, foliation) more than intact rock strength. Key parameters:
\begin{itemize}
\item \textbf{RQD (Rock Quality Designation)}: \% of core pieces $>$10 cm in a drill run. $>$75\% = excellent, $<$25\% = very poor.
\item \textbf{Joint spacing, persistence, aperture, roughness, infilling} (clay, calcite, open).
\item \textbf{Groundwater}: pressure, flow, chemical aggressiveness.
\item \textbf{Weathering grade}: Fresh $\rightarrow$ Slightly $\rightarrow$ Moderately $\rightarrow$ Highly $\rightarrow$ Completely $\rightarrow$ Residual soil.
\end{itemize}
\end{propriete}

\begin{methode}[Evaluating an engineering site (dam or tunnel)]
\begin{enumerate}
\item \textbf{Read the section}: list rock types, dip/dip direction, all discontinuities (faults, joints, unconformities, karst).
\item \textbf{Rank foundation quality}:
\begin{itemize}
\item \emph{Excellent}: fresh massive igneous/metamorphic (granite, gneiss, basalt), thick competent limestone.
\item \emph{Good}: moderately jointed competent sedimentary (sandstone, well-cemented conglomerate).
\item \emph{Poor}: weathered rock, shale/clay (swelling, low shear strength), fractured zones, karstified limestone.
\item \emph{Very poor}: unconsolidated sediments (sand, gravel, silt), peat, landslide debris.
\end{itemize}
\item \textbf{Check watertightness (dams)}: reservoir must hold water. Faults, open joints, karst conduits, permeable beds crossing dam axis $\rightarrow$ leakage risk. Require grout curtain, cut-off wall, or upstream blanket.
\item \textbf{Check stability (tunnels/dam abutments)}: avoid tunnels parallel to steeply dipping weak beds (wedge failure). For dam abutments, avoid beds dipping \emph{towards} the valley (planar sliding). Prefer beds dipping \emph{into} the abutment.
\item \textbf{Recommend best site} by elimination, justify with geology, propose treatment for residual defects (grouting, rock bolts, drainage, dental concrete).
\end{enumerate}
\textbf{Reflex}: ``Strongest rock'' is wrong if cut by open fault or karst $\rightarrow$ \emph{discontinuities control behaviour}.
\end{methode}

\begin{exoresolu}[Choosing a dam site]
A hydroelectric dam is planned on a river in the Western Highlands of Cameroon. Three sites are available:
\begin{itemize}
\item \textbf{Site A}: fresh, massive basalt flows; joints are tight and widely spaced ($>$2 m).
\item \textbf{Site B}: thick limestone with caves and solution channels (karst); a fault crosses the valley perpendicular to the axis.
\item \textbf{Site C}: deeply weathered (lateritised) gneiss, soft to $15\ \mathrm{m}$ depth, overlying sound gneiss.
\end{itemize}
\begin{enumerate}
\item Evaluate each site for dam foundation and reservoir watertightness.
\item Choose the best site and justify your choice.
\item Suggest one treatment that could make Site C acceptable.
\end{enumerate}
\tcblower
\textbf{1. Evaluation.}
\begin{itemize}
\item \textbf{Site A}: Massive basalt has high compressive strength ($>$150 MPa), low deformability, excellent bearing capacity. Tight, widely spaced joints transmit negligible water $\rightarrow$ good natural watertightness. \emph{Excellent foundation and seal.}
\item \textbf{Site B}: Intact limestone can be strong, but \emph{karst} (caves, solution channels) creates high secondary permeability $\rightarrow$ reservoir would leak severely. The fault provides a seepage path and a zone of crushed, weak rock (fault gouge). Foundation support undermined by cavities. \emph{Unsuitable without extremely costly treatment (extensive grouting, concrete diaphragm wall).}
\item \textbf{Site C}: Upper $15\ \mathrm{m}$ of lateritised gneiss is soft, compressible, erodible $\rightarrow$ poor foundation if left in place. However, sound gneiss at depth is excellent. \emph{Potentially acceptable} if weathered mantle is removed or bypassed.
\end{itemize}

\textbf{2. Choice.} \textbf{Site A} is the best: it combines a strong, stable foundation with natural watertightness, requiring minimal treatment and therefore minimal cost and construction risk.

\textbf{3. Treatment of Site C.} Excavate (strip) the entire weathered mantle down to sound gneiss before placing dam concrete (open excavation or cut-off trench). Alternatively, found the dam on deep foundations (piles/caissons) socketed into fresh gneiss. A \emph{grouted cut-off curtain} beneath the dam axis would seal residual joint seepage in the fresh rock. (Standard practice in tropical terrains like Cameroon where lateritic weathering profiles are deep.)
\end{exoresolu}

\begin{experience}[Tunnel Site Investigation]
For a tunnel, the key question is \emph{stand-up time} of the unsupported span. In the \emph{Western Highlands} (volcanic rocks), a tunnel through fresh basalt with wide joint spacing needs little support. Through weathered tuff or fault gouge, it requires immediate support (shotcrete, rock bolts, steel ribs). The \emph{New Austrian Tunnelling Method (NATM)} uses the rock mass itself as a load-bearing component: monitor deformation, apply support as needed. Geological mapping of the tunnel face every advance is mandatory.
\end{experience}

\courssub{Skill 5: Assessing Slope Stability with the Factor of Safety}

\begin{definition}[Factor of Safety (FoS)]
For a potential sliding mass on a defined slip surface:
\[ F = \frac{\text{Resisting (stabilising) Forces}}{\text{Driving (destabilising) Forces}} = \frac{\sum \tau_{\text{resist}} \times L}{\sum \tau_{\text{drive}} \times L} \]
where $\tau_{\text{resist}} = c' + (\sigma_n - u) \tan \varphi'$ (Mohr--Coulomb criterion).
\begin{itemize}
\item $c'$ = effective cohesion (kPa), $\varphi'$ = effective friction angle ($^\circ$).
\item $\sigma_n$ = total normal stress on slip surface, $u$ = pore water pressure.
\item $(\sigma_n - u) = \sigma'_n$ = effective normal stress.
\end{itemize}
Interpretation: $F > 1$ stable; $F = 1$ limit equilibrium (failure imminent); $F < 1$ failure occurring.
\end{definition}

\begin{methode}[Factor-of-safety problems]
\begin{enumerate}
\item \textbf{Identify the slip surface} (planar, circular, wedge) and the sliding mass geometry.
\item \textbf{Calculate driving force}: component of weight $W$ parallel to slip surface ($W \sin \alpha$ for planar), plus water pressure forces, seismic forces, surcharge.
\item \textbf{Calculate resisting force}: shear strength integrated along slip surface. For planar: $R = c' L + (W \cos \alpha - U) \tan \varphi'$, where $U = u \times L$ is uplift force.
\item \textbf{Compute $F = R / D$}.
\item \textbf{Analyse water effect}: rising $u$ reduces $\sigma'_n$ $\rightarrow$ reduces friction; water adds weight ($W \uparrow$); saturation may reduce $c'$ (softening). Heavy rain is the classic trigger.
\item \textbf{Propose stabilisation}: lower $u$ (drainage), reduce $W$ or $\alpha$ (regrading, benching), increase $R$ (retaining wall, rock bolts, soil nails, buttress), protect surface (vegetation, shotcrete).
\end{enumerate}
\end{methode}

\begin{exoresolu}[A road-cut slope in the West Region]
A road cutting between Bafoussam and Mbouda is excavated in weathered gneiss overlying a clay-rich layer that dips $25^\circ$ towards the road. During the rainy season, blocks slide along the clay layer. The resisting force along the slip surface is estimated at $900\ \mathrm{kN}$ per metre of slope; the driving force is $600\ \mathrm{kN/m}$ in the dry season, rising to $1050\ \mathrm{kN/m}$ when the clay is saturated.
\begin{enumerate}
\item Calculate the factor of safety in the dry season and comment.
\item Calculate the factor of safety in the rainy season and comment.
\item Explain, using $\tau = c + (\sigma - u)\tan\varphi$, why saturation causes failure.
\item Propose two engineering measures to stabilise the slope.
\end{enumerate}
\tcblower
\textbf{1. Dry season.}
$F_{\text{dry}} = \dfrac{900}{600} = 1.5 > 1$.
The slope is \emph{stable} in dry conditions, with a comfortable margin (typical design $F \ge 1.3$--$1.5$ for permanent slopes).

\textbf{2. Rainy season.}
$F_{\text{wet}} = \dfrac{900}{1050} \approx 0.86 < 1$.
Driving forces exceed resisting forces: the slope \emph{fails}, as observed.

\textbf{3. Role of water.} Shear resistance follows Mohr--Coulomb: $\tau = c' + (\sigma_n - u)\tan\varphi'$.
\begin{itemize}
\item Infiltration raises pore pressure $u$ in the clay layer $\rightarrow$ effective normal stress $\sigma'_n = \sigma_n - u$ decreases $\rightarrow$ frictional resistance drops.
\item Water increases unit weight of the sliding mass $\rightarrow$ driving force $W \sin \alpha$ increases.
\item Clay softens when saturated $\rightarrow$ cohesion $c'$ may decrease.
\end{itemize}
The combination of reduced resistance and increased driving force pushes $F$ below 1.

\textbf{4. Stabilisation measures} (any two, but best are):
\begin{itemize}
\item \textbf{Drainage}: deep trench drains or horizontal boreholes to lower water table in the clay layer, reducing $u$ (most effective, treats cause).
\item \textbf{Regrading}: flatten slope angle or create benches to reduce driving moment.
\item \textbf{Toe support}: retaining wall, rock fill buttress, or soil nails/rock bolts to add resisting force.
\item \textbf{Surface protection}: vegetation, geomembrane, or shotcrete to limit infiltration.
\end{itemize}
\end{exoresolu}

\begin{attention}[Slope Stability Traps]
\begin{itemize}
\item Do not confuse \emph{total stress} analysis (undrained, $\varphi_u=0$) with \emph{effective stress} analysis (drained, $c', \varphi'$). Long-term stability of clay slopes requires effective stress parameters.
\item A factor of safety of 1.0 is \emph{not} safe; it is the point of collapse. Design requires $F \ge 1.3$ (temporary) to 1.5 (permanent).
\item \emph{Tension cracks} at the slope crest fill with water, adding a hydrostatic driving force and reducing the resisting length $\rightarrow$ always check for tension cracks in analyses.
\end{itemize}
\end{attention}

\courssub{Skill 6: Selecting the Appropriate Prospection Method}

\begin{propriete}[Geophysical Methods and Physical Contrasts]
\begin{center}
\begin{tabularx}{\linewidth}{|l|X|X|}
\hline
\rowcolor{popblueL} \textbf{Method} & \textbf{Physical Property Measured} & \textbf{Typical Targets} \\
\hline
Magnetic (ground/airborne) & Magnetic susceptibility, remanence & Magnetite/Fe ores, Ni-Cr laterites (serpentinite), basement structure, dykes, faults \\
\hline
Gravity & Density contrast ($\Delta\rho$) & Massive sulphides (dense), chromite, Fe-Mn; Salt domes (low $\rho$ $\rightarrow$ oil traps); Basin architecture \\
\hline
Electrical Resistivity / IP & Resistivity ($\rho$), Chargeability & Disseminated sulphides (IP), groundwater, weathering thickness, clay layers, geothermal \\
\hline
Seismic Reflection/Refraction & Acoustic impedance ($Z=\rho V_p$), Velocity & Sedimentary layering, petroleum traps (anticlines, faults, pinch-outs), bedrock depth \\
\hline
Radiometric (Gamma-ray) & K, U, Th concentration & Uranium, thorium, potash; mapping granites, alteration zones \\
\hline
Electromagnetic (EM) & Electrical conductivity & Massive sulphides, graphite, saline water, permafrost \\
\hline
\end{tabularx}
\end{center}
\end{propriete}

\begin{methode}[Matching geophysical methods to targets]
\begin{enumerate}
\item \textbf{Define the target's physical contrast} with host rocks: density, magnetism, resistivity/conductivity, seismic velocity, radioactivity.
\item \textbf{Select the method(s)} that best detect that contrast (see table above). Often a \emph{combination} is used (e.g., Magnetics + Gravity for Fe ore; Seismic + Gravity for salt domes).
\item \textbf{Sequence the campaign}:
\begin{enumerate}
\item Remote sensing (satellite, air photos) + Geological mapping (lithology, structure).
\item Regional geophysics (airborne mag, radiometric) + Geochemical surveys (soil, stream sediment, laterite).
\item Detailed ground geophysics (grid-based) over anomalies.
\item \textbf{Trenching / Drilling} (the \emph{only} direct proof of ore, thickness, grade).
\end{enumerate}
\item \textbf{Justify every choice} by the physical property contrast $\rightarrow$ this justification earns marks.
\end{enumerate}
\textbf{Reflex}: Geophysics gives \emph{anomalies} (indirect evidence). Drilling gives \emph{intersections} (direct evidence). Never claim a deposit is ``proven'' by geophysics alone.
\end{methode}

\begin{exoresolu}[Choosing methods for three targets]
A mining company holds three permits in Cameroon:
\begin{enumerate}
\item Search for nickel--chromium laterite over ultramafic rocks in the East Region.
\item Search for oil traps in the offshore Rio del Rey basin.
\item Search for a buried magnetite iron body near the Congo craton margin.
\end{enumerate}
For each permit, choose the most appropriate geophysical method and justify your choice; then state the final step that proves the deposit.
\tcblower
\textbf{(i) Ni--Cr laterite (East Region).}
The ore is a \emph{weathering profile} (laterite) whose thickness, grade, and clay content vary vertically and laterally.
\begin{itemize}
\item \textbf{Primary method}: \emph{Electrical Resistivity} (or Induced Polarisation). The soft, clay-rich, water-saturated laterite has very low resistivity ($<50\ \Omega\!\cdot\!m$) compared to fresh ultramafic bedrock ($>500\ \Omega\!\cdot\!m$). Resistivity sections map the base of laterite and internal zonation.
\item \textbf{Complementary}: \emph{Magnetic survey} (airborne then ground) to outline the serpentinised ultramafic body (high susceptibility) before detailed resistivity.
\item \textbf{Geochemical}: Auger/soil sampling on grid for Ni, Cr, Co grades.
\end{itemize}

\textbf{(ii) Offshore petroleum (Rio del Rey basin).}
\begin{itemize}
\item \textbf{Primary method}: \emph{Seismic Reflection} (2D then 3D). Contrasts in acoustic impedance between sedimentary beds reveal geometry: anticlines, growth faults, rollover structures, stratigraphic pinch-outs $\rightarrow$ potential \emph{traps} at 2--4 km depth.
\item \textbf{Complementary}: \emph{Gravity} (and magnetic) to map basement structure and locate low-density salt domes (major trap type in Gulf of Guinea).
\item \textbf{Note}: Seismic reveals \emph{structure}, not fluid content. ``Bright spots'' (amplitude anomalies) suggest gas but require drilling to confirm.
\end{itemize}

\textbf{(iii) Buried magnetite body (Congo craton margin).}
Magnetite ($\ce{Fe3O4}$) is strongly ferrimagnetic and very dense ($\rho \approx 5.2\ \mathrm{t/m^3}$).
\begin{itemize}
\item \textbf{Primary method}: \emph{Magnetic survey} (ground or high-resolution airborne). A strong, sharp positive anomaly directly over the body.
\item \textbf{Confirmation}: \emph{Gravity survey} confirms a coincident high-density mass. Combined modelling constrains depth, shape, and tonnage.
\end{itemize}

\textbf{Final proof (all three cases).}
Geophysics only indicates \emph{anomalies}. The deposit is proven and evaluated only by \textbf{drilling} (diamond core drilling for ore, rotary for oil), with \emph{core logging, sampling, and assay/analysis}. This provides direct evidence of mineralisation, thickness, grade, and continuity. The sequence \emph{indirect $\rightarrow$ direct} is the universal logic of mineral and petroleum exploration.
\end{exoresolu}

\begin{savaistu}[Exploration in Cameroon]
Cameroon's mining code (Law No. 2016/017) requires an \emph{Environmental and Social Impact Assessment (ESIA)} before any exploration drilling. The \emph{Cameroon Geological and Mining Research Institute (IRGM)} and \emph{National Hydrocarbons Corporation (SNH)} manage geological data. Major discoveries: \emph{Nkamouna} (Co-Ni-Mn laterite, East), \emph{Mbalam/Nabeba} (Fe, border with Congo), \emph{Logbaba} (gas, Douala basin), \emph{Sanharo} (Au, East). Exploration follows the sequence: stream sediment geochemistry $\rightarrow$ soil grid $\rightarrow$ trenching $\rightarrow$ drilling.
\end{savaistu}

\courssub{Skill 7: Choosing a Mining Method and Assessing Its Environmental Impact}

\begin{propriete}[Mining Method Selection Criteria]
\begin{center}
\begin{tabularx}{\linewidth}{|l|X|X|}
\hline
\rowcolor{popblueL} \textbf{Deposit Characteristic} & \textbf{Surface Mining} & \textbf{Underground Mining} \\
\hline
Depth & Shallow ($<$100--200 m typically) & Deep ($>$100--200 m) \\
\hline
Geometry & Tabular, horizontal/low dip, blanket, placer & Steeply dipping veins, pipes, massive lenses \\
\hline
Rock Strength & Weak to moderate (ore and waste) & Competent ore, strong hanging/footwall \\
\hline
Stripping Ratio (waste:ore) & Low to moderate (economic limit $\sim$5--10:1) & N/A (but high stripping $\rightarrow$ go underground) \\
\hline
Production Rate & High (10$^4$--10$^6$ t/day) & Lower (10$^2$--10$^4$ t/day) \\
\hline
Capital Cost & High initial (equipment, prestripping) & High initial (shaft/decline, infrastructure) \\
\hline
Operating Cost & Lower (\$/t) & Higher (\$/t) \\
\hline
Safety & Better (open environment) & Higher risk (ground control, ventilation) \\
\hline
\end{tabularx}
\end{center}
\end{propriete}

\begin{methode}[Mining method + environmental assessment]
\begin{enumerate}
\item \textbf{Classify deposit geometry} (depth, dip, thickness, continuity, rock strength).
\item \textbf{Choose method}:
\begin{itemize}
\item \emph{Surface}: Open pit (hard rock), Strip mining (coal, flat seams), Quarry (construction materials), Placer/Dredging (alluvial), Mountaintop removal (coal).
\item \emph{Underground}: Room-and-pillar, Cut-and-fill, Sublevel stoping, Block caving, Longwall (coal).
\end{itemize}
\item \textbf{Apply stripping-ratio test} for surface mining: economic only while waste:ore ratio $<$ threshold.
\item \textbf{List environmental impacts}:
\begin{itemize}
\item \emph{Surface}: Land scarring, deforestation, habitat loss, dust, noise, waste dumps, \emph{Acid Mine Drainage (AMD)} from sulphide oxidation ($\ce{FeS2 + 3.5O2 + H2O -> Fe^2+ + 2SO4^2- + 2H+}$), siltation of rivers, visual intrusion.
\item \emph{Underground}: Subsidence, mine water pumping/treatment, spoil heaps, lower land take but potential for AMD, processing chemicals (cyanide, mercury).
\end{itemize}
\item \textbf{Propose mitigation} (hierarchy: Avoid $\rightarrow$ Minimise $\rightarrow$ Mitigate $\rightarrow$ Offset):
\begin{itemize}
\item Progressive rehabilitation (backfill, regrade, topsoil, revegetate with native species).
\item Water management: diversion drains, settling ponds, AMD treatment (lime, constructed wetlands).
\item Dust/noise control: water sprays, enclosures, scheduling.
\item Waste management: separate potentially acid-generating (PAG) waste, encapsulate.
\item \textbf{EIA mandatory before permitting} (Cameroon Law No. 96/12, Decree 2013/0171).
\end{itemize}
\item \textbf{Balance} economic benefits (jobs, GDP, exports, infrastructure) vs environmental/social costs in essay answers.
\end{enumerate}
\end{methode}

\begin{exoresolu}[Exploiting a coastal placer and a deep gold vein]
Two projects are proposed in Cameroon:
\begin{itemize}
\item \textbf{Project 1}: Rutile-rich sands in surface deposits along the coastal plain, thickness $3\ \mathrm{m}$, extending over several km$^2$.
\item \textbf{Project 2}: A gold-bearing quartz vein, $1.5\ \mathrm{m}$ thick, dipping $70^\circ$, traced to $400\ \mathrm{m}$ depth in the East Region.
\end{itemize}
\begin{enumerate}
\item For each project, choose and justify a mining method.
\item State two environmental impacts specific to each method.
\item Outline a rehabilitation plan for Project 1.
\end{enumerate}
\tcblower
\textbf{1. Choice of method.}

\emph{Project 1 (Coastal Rutile Placer):} The ore is unconsolidated sand, at surface, thin but laterally extensive over km$^2$. Ideal for \textbf{surface mining} by \emph{strip mining} (dry) or \emph{dredging} (wet, if below water table). Heavy minerals (rutile, ilmenite, zircon) separated by \emph{gravity concentration} (spirals, shaking tables, wet high-intensity magnetic separation - WHIMS). No blasting needed; low stripping ratio (overburden often minimal).

\emph{Project 2 (Deep Gold Vein):} Narrow vein ($1.5\ \mathrm{m}$), steep dip ($70^\circ$), depth $400\ \mathrm{m}$. Open pit would require removing enormous waste rock (stripping ratio $>$20:1) $\rightarrow$ uneconomic. Requires \textbf{underground mining}: \emph{Vertical shaft} (or decline/adit if topography permits) to access depth; \emph{levels} driven along strike; \emph{stoping} (e.g., cut-and-fill or longhole stoping) to extract the vein. Ore hauled to shaft, hoisted to surface for processing (crushing, grinding, gravity/flotation, cyanide leaching).

\textbf{2. Environmental impacts.}
\begin{itemize}
\item \emph{Project 1 (Surface/Dredge)}:
\begin{enumerate}
\item Large-scale vegetation clearance and topsoil loss over km$^2$; disruption of farming/fishing communities; habitat destruction in coastal forest/mangrove ecotone.
\item Disturbance of water table; discharge of process water (slimes, residual reagents) into coastal streams/lagoons $\rightarrow$ siltation, turbidity, potential chemical pollution affecting fisheries.
\end{enumerate}
\item \emph{Project 2 (Underground)}:
\begin{enumerate}
\item \textbf{Ground subsidence} above stoped areas $\rightarrow$ damage to surface infrastructure, forest, water courses.
\item \textbf{Acid Mine Drainage (AMD)} if vein/sulphides (pyrite, arsenopyrite) oxidise in mine workings and waste rock; discharge of acidic, metal-laden water. \textbf{Cyanide/mercury risk} from gold processing plant tailings.
\end{enumerate}
\end{itemize}

\textbf{3. Rehabilitation Plan for Project 1 (Coastal Placer).}
Adopt \emph{progressive rehabilitation} concurrent with mining:
\begin{enumerate}
\item \textbf{Topsoil stripping \& stockpiling}: Remove and store topsoil (0.3--0.5 m) separately in berms before mining each block.
\item \textbf{Backfilling}: As the dredge/strip face advances, immediately backfill the worked-out void with barren sand tailings (from the separation plant) to restore original topography.
\item \textbf{Landforming}: Regrade backfilled surface to original gentle slopes; ensure drainage integration with natural streams.
\item \textbf{Topsoil replacement}: Spread stockpiled topsoil evenly over reshaped surface.
\item \textbf{Revegetation}: Plant native pioneer species (legumes, grasses) followed by climax forest species; include agroforestry species if community requests.
\item \textbf{Water management}: Closed-circuit process water with settling ponds; monitor effluent quality (pH, TSS, heavy metals) before any discharge.
\item \textbf{Monitoring \& Handover}: Monitor vegetation survival, soil fertility, water quality for 3--5 years post-closure; obtain regulatory sign-off before final relinquishment.
\end{enumerate}
This restores land capability (agriculture, forest) while keeping the operation profitable $\rightarrow$ standard expected under Cameroonian environmental regulations (Decree 2013/0171 on ESIA).
\end{exoresolu}

\begin{aretenir}[Key Exam Points for Economic Geology]
\begin{itemize}
\item \textbf{Coal}: Rank series, proximate analysis trends, coalification drivers (T, P, time), economic value vs rank.
\item \textbf{Petroleum}: 5 elements (source, reservoir, seal, trap, migration), trap types (structural/stratigraphic), fluid contacts (GOC, OWC), migration timing.
\item \textbf{Reserves}: Resources $\neq$ Reserves; Volume $\times$ Density = Tonnage; Grade $\times$ Tonnage = Metal; Recovery factor; Life = Reserve / Production.
\item \textbf{Engineering}: Discontinuities control rock mass; Dam = foundation + watertightness; Tunnel = stand-up time; Slope $F = R/D$; Water reduces $F$ via pore pressure.
\item \textbf{Prospection}: Physical contrast $\rightarrow$ Method; Sequence: Remote $\rightarrow$ GeoChem/GeoPhys $\rightarrow$ Drilling (proof).
\item \textbf{Mining}: Geometry $\rightarrow$ Method (Surface vs Underground); Stripping ratio limit; Impacts (AMD, subsidence, land take); Mitigation hierarchy; EIA mandatory.
\end{itemize}
\end{aretenir}

\newpage

\coursec{Exercises}

\courssub{I check my knowledge}

\begin{exercice}[Multiple Choice Questions]
\textbf{1.} Which of the following sequences correctly represents the increase in rank of coal?
\begin{enumerate}[label=\Alph*.]
\item Peat $\rightarrow$ Lignite $\rightarrow$ Bituminous $\rightarrow$ Anthracite
\item Lignite $\rightarrow$ Peat $\rightarrow$ Anthracite $\rightarrow$ Bituminous
\item Bituminous $\rightarrow$ Anthracite $\rightarrow$ Lignite $\rightarrow$ Peat
\item Anthracite $\rightarrow$ Bituminous $\rightarrow$ Lignite $\rightarrow$ Peat
\end{enumerate}

\textbf{2.} A petroleum trap formed by the termination of a reservoir bed against an impermeable fault plane is classified as a:
\begin{enumerate}[label=\Alph*.]
\item Structural trap (anticline)
\item Stratigraphic trap (pinch-out)
\item Structural trap (fault trap)
\item Combination trap (salt dome)
\end{enumerate}

\textbf{3.} The most important property of a reservoir rock for petroleum accumulation is:
\begin{enumerate}[label=\Alph*.]
\item High compressive strength
\item High porosity and permeability
\item High density
\item Low capillary pressure
\end{enumerate}

\textbf{4.} In slope stability analysis, the factor of safety (FoS) for a planar failure along a discontinuity is given by:
\begin{enumerate}[label=\Alph*.]
\item $\dfrac{cA + W\cos\psi \tan\phi}{W\sin\psi}$
\item $\dfrac{cA + W\sin\psi \tan\phi}{W\cos\psi}$
\item $\dfrac{W\sin\psi}{cA + W\cos\psi \tan\phi}$
\item $\dfrac{W\cos\psi}{cA + W\sin\psi \tan\phi}$
\end{enumerate}
where $c$ = cohesion, $A$ = area of sliding surface, $W$ = weight of block, $\psi$ = dip of discontinuity, $\phi$ = angle of internal friction.
\end{exercice}

\begin{exercice}[True/False with Justification]
State whether each statement is True or False. Justify your answer in one or two sentences.
\begin{enumerate}
\item Coal is a sedimentary rock formed exclusively from the remains of marine organisms.
\item Natural gas (methane) can be generated at lower temperatures than crude oil during the maturation of kerogen.
\item A dam founded on a limestone bedrock with well-developed karst features is generally safe without grouting.
\item The gravity method of geophysical prospection is most suitable for detecting massive sulfide ore bodies.
\item In open-pit mining, the stripping ratio increases as the pit deepens, assuming constant slope angles.
\end{enumerate}
\end{exercice}

\begin{exercice}[Definitions and Key Concepts]
Define the following terms precisely, using geological terminology:
\begin{enumerate}
\item \cle{Vitrinite reflectance ($R_o$)} and its use in petroleum exploration.
\item \cle{Source rock}, \cle{Reservoir rock}, \cle{Cap rock} and \cle{Trap} in a petroleum system.
\item \cle{Stripping ratio} and \cle{Break-even stripping ratio} in surface mining.
\item \cle{Angle of repose} and its significance in waste dump design.
\item \cle{Grout curtain} in dam engineering.
\end{enumerate}
\end{exercice}

\begin{exercice}[Direct Calculations]
\begin{enumerate}
\item A coal sample has the following proximate analysis (air-dried basis): Moisture = 8\%, Volatile Matter = 32\%, Fixed Carbon = 50\%, Ash = 10\%. Calculate the \cle{Calorific Value} (CV) on a dry, mineral-matter-free basis using the \cle{Dulong formula} (approximate): $CV (\mathrm{kJ/kg}) = 337C + 1442(H - O/8) + 93S$, where C, H, O, S are weight percentages on a dry, mineral-matter-free basis. Assume the dry, mineral-matter-free composition is C=80\%, H=5\%, O=10\%, S=1\%, N=2\%, others=2\%.
\item A sandstone reservoir has a porosity of 22\% and a water saturation ($S_w$) of 30\%. If the formation volume factor of oil ($B_o$) is 1.3 rb/stb, calculate the recoverable oil volume per acre-foot of reservoir rock (1 acre-foot = 43,560 ft$^3$). Assume recovery factor = 35\%.
\item A planar slide occurs on a joint dipping at $35^\circ$ in a slope face dipping at $55^\circ$. The joint has cohesion $c = 15\ \mathrm{kPa}$ and friction angle $\phi = 30^\circ$. The rock unit weight is $26\ \mathrm{kN/m^3}$. For a block height of 10 m and unit width, calculate the Factor of Safety (FoS) using the formula from Exercise 1, Question 4.
\item A gravity survey over a horizontal cylindrical ore body (radius 50 m, depth to centre 150 m, density contrast 2.0 g/cm$^3$) gives a maximum anomaly. Estimate the maximum Bouguer anomaly (in mGal) using the formula $\Delta g_{max} = \dfrac{2G \Delta \rho \pi R^2}{z}$, where $G = 6.67 \times 10^{-11}\ \mathrm{N\,m^2/kg^2}$, $\Delta \rho$ in kg/m$^3$, $R$ and $z$ in m. (1 mGal = $10^{-5}\ \mathrm{m/s^2}$).
\end{enumerate}
\end{exercice}

\courssub{I practise}

\begin{exercice}[Coal Rank Determination and Basin History]
The table below shows proximate analysis results (dry, mineral-matter-free basis) for five coal samples collected at increasing depths in a sedimentary basin.

\begin{center}
\begin{tabular}{|c|ccccc|}\hline
Sample & A & B & C & D & E \\ \hline
Depth (m) & 200 & 500 & 800 & 1200 & 1800 \\
Volatile Matter (\%) & 42 & 35 & 28 & 18 & 8 \\
Fixed Carbon (\%) & 58 & 65 & 72 & 82 & 92 \\ \hline
\end{tabular}
\end{center}

\begin{enumerate}
\item Classify each sample into a coal rank (Peat, Lignite, Sub-bituminous, Bituminous, Anthracite) using standard rank parameters (e.g., ASTM D388).
\item Plot a graph of Volatile Matter vs. Depth and Fixed Carbon vs. Depth. Describe the trend.
\item Deduce the thermal and burial history of the basin. What was the approximate maximum paleo-temperature reached by Sample E? (Use the approximation: $R_o (\%) \approx 0.018 \times \text{Temp}(^\circ\mathrm{C}) - 0.6$ for $R_o > 0.5\%$, and the relationship between $R_o$ and volatile matter for bituminous coals: $VM_{daf} \approx 60 - 50 R_o$).
\item Discuss the economic implications of this rank variation for mining and utilisation (e.g., coking vs. steam coal).
\end{enumerate}
\end{exercice}

\begin{exercice}[Petroleum System Analysis]
A seismic cross-section through the \cle{Rio del Rey Basin} (offshore Cameroon) reveals the following features from bottom to top:
\begin{itemize}
\item Pre-rift basement (Precambrian).
\item Syn-rift sequence: Continental clastics (sandstones, shales) with lacustrine shales (potential source rocks).
\item Post-rift sequence: Marine shales (excellent source rocks, TOC up to 5\%), overlain by deltaic sandstones (reservoirs) and sealed by marine shales (cap rocks).
\item A major growth fault cuts the syn-rift and lower post-rift sequences, with rollover anticlines on the downthrown side.
\item An unconformity truncates the lower post-rift sequence, overlain by younger shales.
\end{itemize}

\begin{enumerate}
\item Identify and name \textbf{four} distinct types of petroleum traps present in this setting. For each, sketch a small labelled cross-section (use TikZ) and explain the trapping mechanism.
\item The lacustrine shales are Type I/II kerogen, while the marine shales are Type II. Explain the implications for oil vs. gas generation and the likely timing of expulsion relative to trap formation.
\item A well drilled on the crest of a rollover anticline encounters the reservoir sandstone at 3200 m TVD (True Vertical Depth). The geothermal gradient is $3.2^\circ\mathrm{C}/100\ \mathrm{m}$ and surface temperature is $25^\circ\mathrm{C}$. Calculate the reservoir temperature. If the oil window for Type II kerogen is $60-120^\circ\mathrm{C}$, has the source rock reached peak oil generation at this depth? (Assume source rock is at 3500 m).
\item List the key data (well logs, core, seismic, geochemical) required to de-risk a prospect in this basin before drilling.
\end{enumerate}
\end{exercice}

\begin{exercice}[Slope Stability in Jointed Basalt]
A road cut on the flanks of \cle{Mount Cameroon} exposes a 25 m high slope in jointed basalt. The slope face dips at $60^\circ$ towards $120^\circ$ (azimuth). Two dominant joint sets are mapped:
\begin{itemize}
\item Set J1: Dip $45^\circ$, Dip Direction $115^\circ$ (rough, undulating, $JRC = 12$, $JCS = 150\ \mathrm{MPa}$).
\item Set J2: Dip $70^\circ$, Dip Direction $290^\circ$ (planar, smooth, $JRC = 4$, $JCS = 100\ \mathrm{MPa}$).
\end{itemize}
The basalt has unit weight $\gamma = 28\ \mathrm{kN/m^3}$, intact UCS = 200 MPa. Groundwater table is at the slope crest.

\begin{enumerate}
\item Using \cle{Markland's test} (kinematic analysis), determine which joint set(s) can produce planar sliding, wedge sliding, or toppling. Assume friction angle $\phi = 30^\circ$ for both joints.
\item For the critical failure mode, calculate the Factor of Safety (FoS) using the \cle{Hoek-Brown} criterion for the joint (or Barton-Bandis if preferred). Simplify: assume a planar block of height 25 m, unit width, sliding on J1. Water pressure acts on the joint.
\item Propose \textbf{three} practical stabilisation measures suitable for this volcanic rock slope, justifying each with respect to the failure mechanism.
\item Discuss how the presence of \cle{paleosols} (weathered horizons) between lava flows would modify the stability analysis.
\end{enumerate}
\end{exercice}

\begin{exercice}[Prospection Method Selection]
Match each of the following six commodities to the \textbf{most appropriate primary prospection method} (Geological mapping, Geochemical soil sampling, Magnetic survey, Gravity survey, IP/Resistivity, Seismic reflection, Radiometric survey, EM survey). Justify each choice in one sentence, referring to the physical/chemical property exploited.

\begin{enumerate}
\item \cle{Crude oil} in a Tertiary deltaic basin (e.g., Rio del Rey).
\item \cle{Gold} in shear zones within the Pan-African granitoids of East Cameroon.
\item \cle{Limestone} for cement production in the Mesozoic sedimentary cover of North Cameroon.
\item \cle{Iron ore (magnetite)} in Banded Iron Formations (BIF) of the Congo Craton (South Cameroon).
\item \cle{Rock salt (halite)} in a deep evaporite basin.
\item \cle{Groundwater} in the fractured basement aquifers of the Adamawa Plateau.
\end{enumerate}
\end{exercice}

\courssub{I prepare for the GCE}

\begin{exercice}[Structured Question: Petroleum Traps and Drilling Strategy] (25 marks)
The figure below (to be sketched by candidate or described) shows a geological cross-section through a sedimentary basin. The section includes:
\begin{itemize}
\item A basal unconformity separating folded Precambrian basement from flat-lying Cambrian sandstones.
\item An anticline in the Cambrian sandstones, cored by a shale unit.
\item A normal fault cutting the anticline, downthrowing the northern limb.
\item An overlying Cretaceous shale (source rock) and sandstone (reservoir) sequence, deposited unconformably on the Cambrian.
\item A younger Tertiary shale seal.
\end{itemize}

\begin{enumerate}
\item \textbf{(5 marks)} On the cross-section, identify and label \textbf{five} potential petroleum traps (structural, stratigraphic, combination). For each, state the trap type and the seal mechanism.
\item \textbf{(6 marks)} The Cambrian sandstone is a proven reservoir in a nearby field. The Cretaceous sandstone is a new target. Compare the \cle{risk factors} (Source, Reservoir, Seal, Trap, Timing) for the two reservoirs. Which has the higher chance of success? Justify.
\item \textbf{(4 marks)} A well is proposed to test the Cretaceous sandstone on the downthrown side of the fault. Calculate the \cle{drilling depth} (TVD) to the top of the Cretaceous sandstone given: Seismic two-way time (TWT) to top Cretaceous = 2.4 s; Average velocity above = 3.2 km/s; Kelly bushing elevation = 25 m above sea level; Sea level datum.
\item \textbf{(5 marks)} During drilling, a \cle{kick} is detected at 3200 m. The shut-in drill pipe pressure (SIDPP) is 8 MPa and shut-in casing pressure (SICP) is 12 MPa. Mud weight in use is 1.35 SG. Calculate the \cle{formation pressure} (MPa) and the \cle{kill mud weight} (SG) required to control the well. (Assume hydrostatic gradient of water = 0.1 MPa/m).
\item \textbf{(5 marks)} Discuss the environmental and social considerations specific to offshore drilling in the \cle{Rio del Rey Basin}, including waste management, oil spill contingency, and community relations (e.g., fishing communities).
\end{enumerate}
\end{exercice}

\begin{exercice}[Essay: Geological Factors in Dam Construction] (25 marks)
\textbf{``Discuss the geological factors that must be considered before constructing a large dam and reservoir, illustrating your answer with specific examples from Cameroon.''}

Structure your essay under the following headings (allocate marks approximately as shown):
\begin{enumerate}
\item \textbf{Site Geology and Foundation Conditions} (6 marks): Rock type, structure (folds, faults, joints), weathering, bearing capacity, settlement. Example: \cle{Lom Pangar Dam} (granite/gneiss basement) vs. \cle{Memve'ele Dam} (metamorphic basement with shear zones).
\item \textbf{Reservoir Integrity} (6 marks): Permeability of reservoir floor and flanks, karst, fault leakage, sedimentation rate. Example: Potential karst in limestone areas (e.g., North Cameroon), fault-controlled leakage.
\item \textbf{Seismicity and Seismic Hazard} (4 marks): Regional tectonics (Cameroon Volcanic Line, Benue Trough), induced seismicity, design earthquake parameters.
\item \textbf{Construction Materials} (4 marks): Availability of aggregates, riprap, clay core material, pozzolana (volcanic ash from CVL). Quarry site geology.
\item \textbf{Environmental and Social Geology} (5 marks): Slope stability around reservoir (landslides), water quality (mercury methylation, eutrophication), resettlement on stable ground. Example: \cle{Lake Nyos} gas hazard (though natural, relevant to reservoir gas).
\end{enumerate}
\end{exercice}

\begin{exercice}[Problem Solving: Mining Method Selection and Environmental Management] (25 marks)
A \cle{bauxite deposit} in the \cle{Adamawa Plateau} (lateritic profile over syenite) is being evaluated. The deposit is a flat-lying, laterally extensive blanket, 5--8 m thick, under 2--4 m of overburden (topsoil and lateritic gravel). The bauxite grade averages 45\% Al$_2$O$_3$, 5\% SiO$_2$. The water table is 10 m below surface. The area is a savannah ecosystem with gallery forests along streams. Local communities practice agriculture and grazing.

\begin{enumerate}
\item \textbf{(6 marks)} Compare \cle{open-pit (strip mining)} and \cle{underground (room-and-pillar)} methods for this deposit. Evaluate each under: Capital/Operating cost, Recovery percentage, Dilution, Production rate, Safety, Suitability for thin tabular deposit. Conclude with a recommended method.
\item \textbf{(7 marks)} For the chosen method, identify the \textbf{major environmental impacts} during: (a) Land clearing and overburden removal, (b) Mining and haulage, (c) Waste dump and tailings management, (d) Mine closure. For each, propose a specific mitigation measure.
\item \textbf{(6 marks)} Design a \cle{rehabilitation plan} for the post-mining landscape. Address: Landform design (slopes, drainage), Topsoil management (storage, respreading), Revegetation (species selection: native grasses, legumes, trees for agroforestry), Water management (sediment ponds, stream diversion restoration), Long-term monitoring and community handover.
\item \textbf{(6 marks)} The mine requires an \cle{Environmental Impact Assessment (EIA)} under Cameroonian law (Law No. 96/12). Outline the key components of the EIA report for this project, highlighting the role of the \cle{geologist} in baseline data collection (geology, hydrogeology, geochemistry, geomorphology).
\end{enumerate}
\end{exercice}

\newpage

\coursec{Solutions to the exercises}

\coursec{Corrections to Exercises}

\begin{corrige}[Exercise 1 — Multiple Choice Questions]
\textbf{1. Answer: A.}\par
Coal rank increases with the degree of \cle{coalification} (rising temperature and pressure during burial): Peat $\rightarrow$ Lignite $\rightarrow$ Sub-bituminous $\rightarrow$ Bituminous $\rightarrow$ Anthracite. With increasing rank, moisture and volatile matter decrease while fixed carbon and calorific value increase. Options B, C and D either reverse the sequence or place a low-rank coal after a high-rank one, which is thermodynamically impossible under normal burial.

\textbf{2. Answer: C.}\par
When a reservoir bed is truncated laterally by an \cle{impermeable fault plane} (fault seal), the trap is produced by deformation of the rocks after deposition: it is a \cle{structural trap}, more precisely a \emph{fault trap}. An anticline trap (A) is also structural but is caused by folding, not faulting. A pinch-out (B) is stratigraphic (depositional termination of the bed). A salt dome trap (D) is a combination trap related to diapirism.

\textbf{3. Answer: B.}\par
A reservoir rock must \emph{store} hydrocarbons (\cle{porosity}) and \emph{transmit} them to the wellbore (\cle{permeability}). Compressive strength and density are irrelevant to storage capacity, and a low capillary pressure alone cannot create storage space. Hence high porosity \emph{and} permeability is the essential combination.

\textbf{4. Answer: A.}\par
Resolving the weight $W$ of the block normal and parallel to the discontinuity (dip $\psi$):
\begin{itemize}
\item Driving (shear) force along the plane: $T = W\sin\psi$.
\item Normal force on the plane: $N = W\cos\psi$.
\item Resisting force (Mohr--Coulomb): $R = cA + N\tan\phi = cA + W\cos\psi\tan\phi$.
\end{itemize}
Therefore
\[
FoS \;=\; \frac{\text{resisting forces}}{\text{driving forces}} \;=\; \frac{cA + W\cos\psi\,\tan\phi}{W\sin\psi}.
\]
Option C is the inverse ratio; options B and D mix up $\sin\psi$ and $\cos\psi$ in the resolution of $W$.
\end{corrige}

\begin{corrige}[Exercise 2 — True/False with Justification]
\textbf{1. FALSE.} Coal forms from the accumulation and coalification of \emph{terrestrial plant} material (trees, ferns, mosses) in continental or paralic swamps and mires, not from marine organisms. Marine organic matter (plankton, algae) is the precursor of the kerogen that generates petroleum, not coal.

\textbf{2. TRUE.} During catagenesis, oil is generated mainly in the ``oil window'' (about $60$--$120^\circ$C). However, biogenic methane forms at shallow depth and low temperature by bacterial action, and thermogenic gas is also cracked from kerogen and oil at temperatures \emph{above} the oil window; small amounts of gas are released throughout maturation. Gas generation therefore spans a wider range and begins at lower temperatures than the main oil generation.

\textbf{3. FALSE.} Karstified limestone contains solution channels, cavities and sinkholes which create severe problems of foundation bearing capacity, differential settlement and, above all, reservoir leakage. Systematic \cle{grouting} (grout curtain, consolidation grouting) and careful cavity treatment are indispensable before such a dam can be considered safe.

\textbf{4. FALSE.} Massive sulfide bodies are dense, but they are above all excellent electrical conductors and are highly chargeable; the methods of choice are \cle{Induced Polarisation (IP)}, resistivity and electromagnetic (EM) surveys. The gravity method is best suited to large density contrasts (salt domes, chromite, barite, basin architecture); it is generally not the \emph{most} suitable tool for sulfides.

\textbf{5. TRUE.} With constant pit slope angles, each deeper bench requires the removal of a larger ring of waste around the ore, so the volume of overburden/waste grows faster than the volume of ore. The \cle{stripping ratio} (waste : ore) therefore increases with depth, until the \cle{break-even stripping ratio} is reached, which defines the ultimate pit limit.
\end{corrige}

\begin{corrige}[Exercise 3 — Definitions and Key Concepts]
\textbf{1. Vitrinite reflectance ($R_o$).} The percentage of incident light reflected by the maceral \emph{vitrinite} in a polished rock sample, measured under a microscope in oil immersion. $R_o$ increases irreversibly with the maximum temperature experienced by the rock, so it is a \cle{paleo-thermometer}: in petroleum exploration it is used to determine the thermal maturity of a source rock and to locate it relative to the oil window ($R_o \approx 0.6$--$1.3\%$) and gas window ($R_o > 1.3\%$).

\textbf{2. Elements of the petroleum system.}
\begin{itemize}
\item \cle{Source rock}: a fine-grained sedimentary rock rich in organic matter (kerogen) which, on burial and heating, generates and expels hydrocarbons (e.g. marine shale with high TOC).
\item \cle{Reservoir rock}: a rock with sufficient porosity and permeability to store and yield hydrocarbons (e.g. sandstone, fractured limestone).
\item \cle{Cap rock (seal)}: an impermeable layer (shale, evaporite) that prevents upward migration and confines hydrocarbons in the reservoir.
\item \cle{Trap}: a geometric configuration (structural, stratigraphic or combination) of reservoir and seal that concentrates and retains hydrocarbons in a closed accumulation.
\end{itemize}

\textbf{3. Stripping ratio and break-even stripping ratio.} The \cle{stripping ratio} is the volume (or tonnage) of waste rock that must be removed per unit volume (or tonnage) of ore in surface mining. The \cle{break-even stripping ratio} is the maximum stripping ratio at which surface mining remains economically equivalent to (or more profitable than) underground mining; it fixes the economic depth limit of the open pit.

\textbf{4. Angle of repose.} The maximum angle, measured from the horizontal, at which a pile of loose, unconsolidated material remains stable without sliding (about $35$--$38^\circ$ for most mine waste). In waste dump design, dump slopes must be kept at or below the angle of repose (or benched) to prevent raveling and flow slides.

\textbf{5. Grout curtain.} A vertical or inclined barrier of low permeability created beneath a dam (and into its abutments) by injecting cement or chemical grout into boreholes, filling joints, fissures and cavities in the foundation rock. Its purpose is to reduce seepage under the dam, control uplift pressure and prevent internal erosion (piping).
\end{corrige}

\begin{corrige}[Exercise 4 — Direct Calculations]
\textbf{1. Calorific value (Dulong formula).}\par
Given C $= 80\%$, H $= 5\%$, O $= 10\%$, S $= 1\%$:
\begin{align*}
CV &= 337C + 1442\left(H - \frac{O}{8}\right) + 93S\\
&= 337(80) + 1442\left(5 - \frac{10}{8}\right) + 93(1)\\
&= 26\,960 + 1442(3.75) + 93\\
&= 26\,960 + 5\,407.5 + 93\\
&\approx 32\,460\ \mathrm{kJ/kg} \;(\approx 32.5\ \mathrm{MJ/kg}).
\end{align*}
This high value is consistent with a high-rank bituminous coal.

\textbf{2. Recoverable oil per acre-foot.}\par
Bulk volume $= 43\,560\ \mathrm{ft^3}$, i.e. $7\,758$ bbl per acre-foot.
\begin{align*}
\text{OOIP} &= \frac{7\,758 \times \phi \times (1 - S_w)}{B_o}
= \frac{7\,758 \times 0.22 \times 0.70}{1.3}\\
&= \frac{1\,194.7}{1.3} \approx 919\ \mathrm{stb/acre\text{-}ft}.\\
\text{Recoverable oil} &= 919 \times 0.35 \approx 322\ \mathrm{stb/acre\text{-}ft}.
\end{align*}

\textbf{3. Factor of Safety of the planar slide.}\par
Length of sliding surface: $A = \dfrac{H}{\sin\psi} = \dfrac{10}{\sin 35^\circ} = 17.43\ \mathrm{m^2}$ (unit width).\par
Weight of the triangular block: $W = \tfrac{1}{2}\gamma H^2 \cot\psi = 0.5 \times 26 \times 10^2 \times \cot 35^\circ = 1\,856\ \mathrm{kN}$.
\begin{align*}
FoS &= \frac{cA + W\cos\psi\tan\phi}{W\sin\psi}\\
&= \frac{15(17.43) + 1\,856 \times 0.819 \times 0.5774}{1\,856 \times 0.5736}\\
&= \frac{261.5 + 877.8}{1\,064.6} = \frac{1\,139.3}{1\,064.6} \approx 1.07.
\end{align*}
$FoS \approx 1.07$: the slope is marginally stable (FoS only slightly above 1); any water pressure or vibration could trigger failure. Remediation (drainage, bolting) is advisable.

\textbf{4. Maximum Bouguer anomaly over a horizontal cylinder.}\par
$\Delta\rho = 2.0\ \mathrm{g/cm^3} = 2\,000\ \mathrm{kg/m^3}$, $R = 50$ m, $z = 150$ m.
\begin{align*}
\Delta g_{max} &= \frac{2G\,\Delta\rho\,\pi R^2}{z}
= \frac{2 \times 6.67\times10^{-11} \times 2\,000 \times \pi \times 50^2}{150}\\
&= \frac{2.095\times10^{-3}}{150} = 1.40\times10^{-5}\ \mathrm{m/s^2}.
\end{align*}
Since $1\ \mathrm{mGal} = 10^{-5}\ \mathrm{m/s^2}$: $\Delta g_{max} \approx 1.4\ \mathrm{mGal}$. This is easily detectable with a modern gravimeter (precision $\sim 0.01$ mGal), but requires careful terrain and drift corrections.
\end{corrige}

\begin{corrige}[Exercise 5 — Coal Rank Determination and Basin History]
\textbf{1. Rank classification.}\par
Using volatile matter (dry, ash-free) as the rank parameter (ASTM D388 logic):
\begin{center}
\begin{tabular}{|c|c|c|c|l|}\hline
\rowcolor{popblueL} Sample & Depth (m) & VM (\%) & FC (\%) & Rank \\ \hline
A & 200 & 42 & 58 & Sub-bituminous (low rank, near lignite boundary) \\
B & 500 & 35 & 65 & Sub-bituminous \\
C & 800 & 28 & 72 & Bituminous (high volatile) \\
D & 1200 & 18 & 82 & Bituminous (medium/low volatile) \\
E & 1800 & 8 & 92 & Anthracite \\ \hline
\end{tabular}
\end{center}
(Peat is excluded: all samples are true coals with VM $< 50\%$ daf.)

\textbf{2. Graph and trend.}\par
\begin{popfigure}
\begin{center}
\begin{tikzpicture}[scale=0.85]
\draw[->,thick] (0,0) -- (0,6.2) node[above]{\small Depth (m)};
\draw[->,thick] (0,0) -- (9.5,0) node[right]{\small VM (\%) / FC (\%)};
\foreach \y/\l in {0/0,1.5/500,3/1000,4.5/1500,6/2000}{\draw (0,\y) node[left]{\small \l};}
\foreach \x/\l in {0/0,2/20,4/40,6/60,8/80}{\draw (\x,0.02) node[below]{\small \l};}
% VM curve: (VM/10, depth/333.33)
\draw[very thick,poporange] (4.2,0.6) -- (3.5,1.5) -- (2.8,2.4) -- (1.8,3.6) -- (0.8,5.4);
\fill[poporange] (4.2,0.6) circle(2pt) (3.5,1.5) circle(2pt) (2.8,2.4) circle(2pt) (1.8,3.6) circle(2pt) (0.8,5.4) circle(2pt);
% FC curve
\draw[very thick,popblue] (5.8,0.6) -- (6.5,1.5) -- (7.2,2.4) -- (8.2,3.6) -- (9.2,5.4);
\fill[popblue] (5.8,0.6) circle(2pt) (6.5,1.5) circle(2pt) (7.2,2.4) circle(2pt) (8.2,3.6) circle(2pt) (9.2,5.4) circle(2pt);
\node[poporange,align=center] at (3.2,5.9) {\small VM};
\node[popblue,align=center] at (8.6,5.9) {\small FC};
\end{tikzpicture}
\end{center}
\legende{Volatile matter decreases and fixed carbon increases regularly with depth: the signature of progressive coalification during burial.}
\end{popfigure}
VM falls from $42\%$ to $8\%$ while FC rises from $58\%$ to $92\%$ with increasing depth: a regular, monotonic coalification gradient.

\textbf{3. Thermal and burial history.}\par
The regular increase of rank with depth indicates progressive burial in a subsiding basin under a normal geothermal regime (no evidence of local igneous heating, which would produce rank anomalies). For Sample E:
\[
VM_{daf} \approx 60 - 50R_o \;\Rightarrow\; 8 = 60 - 50R_o \;\Rightarrow\; R_o = \frac{52}{50} = 1.04\%.
\]
\[
R_o \approx 0.018\,T - 0.6 \;\Rightarrow\; T = \frac{R_o + 0.6}{0.018} = \frac{1.64}{0.018} \approx 91^\circ\mathrm{C}.
\]
Sample E therefore reached a maximum paleo-temperature of about $90$--$95^\circ$C. With a typical gradient of $30^\circ$C/km and a surface temperature near $25^\circ$C, this implies a maximum burial of roughly $2.2$ km, i.e. some overburden may have been eroded since maximum burial (uplift and denudation phase).

\textbf{4. Economic implications.}\par
\begin{itemize}
\item Samples A--B (sub-bituminous): high VM, moderate CV; suitable as \emph{steam coal} for power generation; often mined in shallow, low-cost open pits.
\item Samples C--D (bituminous): the medium/low-volatile bituminous coals (D especially) are potential \emph{coking coals} for the steel industry, commanding higher prices; deeper seams require underground mining with higher costs and gas (firedamp) hazards.
\item Sample E (anthracite): high CV, smokeless fuel for domestic and metallurgical uses, but deep, thin, gassy and difficult to mine; economic viability depends on seam thickness and continuity.
\end{itemize}
Thus value per tonne rises with rank, but mining cost also rises with depth; the optimum economic target in this basin is probably the bituminous interval (C--D).
\end{corrige}

\begin{corrige}[Exercise 6 — Petroleum System Analysis]
\textbf{1. Four trap types.}\par
\begin{itemize}
\item \textbf{(a) Rollover anticline (structural, fault-related):} formed on the downthrown side of the growth fault by reverse drag of beds into the fault; closure is provided by the fold, seal by the overlying marine shale.
\item \textbf{(b) Fault trap (structural):} where the growth fault juxtaposes the deltaic reservoir sandstone against impermeable shale across the fault plane (fault seal).
\item \textbf{(c) Unconformity (subcrop) trap (stratigraphic):} truncated lower post-rift reservoir beds sealed beneath the unconformity by the younger shales.
\item \textbf{(d) Pinch-out / onlap trap (stratigraphic):} syn-rift and basal post-rift sandstones wedging out updip against the basement high or onto the unconformity, sealed by lacustrine/marine shales.
\end{itemize}
\begin{popfigure}
\begin{center}
\begin{tikzpicture}[scale=0.95]
% basement
\fill[popmuted!40] (0,0) -- (10,0) -- (10,1.2) -- (7,1.2) -- (6.6,2.2) -- (3,1.4) -- (0,1.0) -- cycle;
\node at (1.2,0.5) {\small Basement};
% syn-rift
\fill[popgold!40] (0,1.0) -- (3,1.4) -- (6.6,2.2) -- (7,1.2) -- (10,1.2) -- (10,2.0) -- (7,2.0) -- (6.6,3.0) -- (3,2.2) -- (0,1.8) -- cycle;
\node at (2.2,1.75) {\small Syn-rift};
% post-rift reservoir
\fill[poporange!50] (0,1.8) -- (3,2.2) -- (6.6,3.0) -- (7,2.0) -- (10,2.0) -- (10,2.9) -- (7.4,2.9) -- (6.6,3.9) -- (3,3.1) -- (0,2.7) -- cycle;
\node at (2.0,2.6) {\small Reservoir ss};
% shale seal
\fill[popblue!30] (0,2.7) -- (3,3.1) -- (6.6,3.9) -- (7.4,2.9) -- (10,2.9) -- (10,3.8) -- (0,3.6) -- cycle;
\node at (1.6,3.35) {\small Shale seal};
% growth fault
\draw[very thick,popdark] (6.6,4.3) -- (6.9,3.4) -- (6.6,2.2) -- (7,1.2);
\node[popdark] at (7.6,4.2) {\small Growth fault};
% unconformity
\draw[very thick,poppink,decorate,decoration={zigzag,segment length=4,amplitude=1.2}] (0,3.6) -- (10,3.8);
\node[poppink] at (9.0,4.15) {\small Unconformity};
% accumulation
\fill[popgreen!70] (5.6,3.15) -- (6.55,3.85) -- (6.6,3.0) -- cycle;
\node[popgreen!60!black] at (5.0,4.35) {\small Oil};
\draw[->,popgreen!60!black] (5.2,4.2) -- (5.9,3.6);
\end{tikzpicture}
\end{center}
\legende{Schematic Niger-Delta-type section: rollover anticline and fault trap against a growth fault, unconformity trap above, pinch-out against basement.}
\end{popfigure}

\textbf{2. Kerogen type, product and timing.}\par
Type I/II lacustrine kerogen is highly oil-prone (high H/C ratio) and generates mainly light oil, often early in the oil window. Type II marine kerogen generates oil plus associated gas, with gas becoming dominant at higher maturity. In a deltaic setting with rapid sedimentation, maturation is driven by fast burial; growth faults and rollover anticlines develop \emph{synsedimentary}, i.e. \emph{before or during} hydrocarbon generation and migration. Timing is therefore favourable: traps were already in place when expulsion occurred (late Miocene to present in such deltas), so the structures could receive and retain the charge. The main risk is that early-formed traps may be breached by later fault reactivation.

\textbf{3. Reservoir temperature and maturity.}\par
At the reservoir (3\,200 m):
\[
T = 25 + \frac{3.2}{100}\times 3\,200 = 25 + 102.4 = 127.4^\circ\mathrm{C}.
\]
At the source rock (3\,500 m):
\[
T = 25 + \frac{3.2}{100}\times 3\,500 = 25 + 112 = 137^\circ\mathrm{C}.
\]
The oil window for Type II kerogen is $60$--$120^\circ$C. At $137^\circ$C the source rock is \emph{past peak oil generation}: it is in the late-oil to wet-gas/condensate zone, where remaining oil begins to crack to gas. Peak oil generation occurred earlier, at shallower burial; the expelled oil has since migrated into the shallower reservoirs. Present-day charge at this depth is increasingly gas-prone.

\textbf{4. Data required to de-risk the prospect.}\par
\begin{itemize}
\item \emph{Seismic:} 3D seismic for structural mapping, fault seal analysis, direct hydrocarbon indicators (bright spots, flat spots), depth conversion.
\item \emph{Well logs:} gamma ray, resistivity, sonic, density, neutron for porosity, saturation, lithology and correlation; pressure data (RFT) for fluid contacts.
\item \emph{Core/cuttings:} porosity--permeability measurements, sedimentology of the reservoir, seal capacity (capillary entry pressure) of the shale.
\item \emph{Geochemistry:} TOC and Rock-Eval of source rocks, kerogen typing, vitrinite reflectance ($R_o$), biomarker oil--source correlation, basin modelling (1D/3D) for timing of generation, migration and trap formation.
\end{itemize}
\end{corrige}

\begin{corrige}[Exercise 7 — Slope Stability in Jointed Basalt]
\textbf{1. Kinematic (Markland) analysis.}\par
Slope face: dip $60^\circ$, dip direction $120^\circ$; friction angle $\phi = 30^\circ$.
\begin{itemize}
\item \textbf{Planar sliding on J1:} J1 dips $45^\circ$ towards $115^\circ$, i.e. within $\pm 20^\circ$ of the slope dip direction, and $\phi (30^\circ) < \psi_{J1} (45^\circ) < \psi_{slope} (60^\circ)$. The plane daylights in the face and is steeper than the friction angle: \textbf{planar sliding is kinematically feasible on J1}.
\item \textbf{Planar sliding on J2:} J2 dips $70^\circ$ towards $290^\circ$, i.e. \emph{into} the slope (opposite dip direction): planar sliding on J2 is impossible. However, J2 is steeper than the slope face and dips into the face: blocks resting on J2 can rotate forward --- \textbf{toppling failure is feasible on J2} (especially in columnar-jointed basalt).
\item \textbf{Wedge sliding:} the line of intersection of J1 and J2 plunges moderately; given that J2 dips steeply into the slope, the intersection does not daylight favourably in the face, so wedge sliding is unlikely. The critical mode is \textbf{planar sliding on J1}.
\end{itemize}

\textbf{2. Factor of Safety (planar block on J1, with water).}\par
Geometry: $H = 25$ m, $\psi = 45^\circ$, unit width.
\[
A = \frac{H}{\sin\psi} = \frac{25}{0.7071} = 35.36\ \mathrm{m^2};\qquad
W = \tfrac{1}{2}\gamma H^2\cot\psi = 0.5 \times 28 \times 625 \times 1 = 8\,750\ \mathrm{kN}.
\]
Water table at the crest $\Rightarrow$ tension crack full; uplift on the joint:
\[
U = \tfrac{1}{2}\gamma_w \frac{H^2}{\sin\psi} = 0.5 \times 10 \times \frac{625}{0.7071} \approx 4\,419\ \mathrm{kN}.
\]
Effective normal stress on the joint:
\[
\sigma_n = \frac{W\cos\psi - U}{A} = \frac{8\,750(0.7071) - 4\,419}{35.36} = \frac{6\,187 - 4\,419}{35.36} = \frac{1\,768}{35.36} \approx 50\ \mathrm{kPa}.
\]
Barton--Bandis shear strength with $JRC = 12$, $JCS = 150$ MPa, $\phi_r = 30^\circ$:
\[
\tau = \sigma_n \tan\!\left[JRC \log_{10}\!\left(\frac{JCS}{\sigma_n}\right) + \phi_r\right]
= 50 \tan\!\left[12\log_{10}\!\left(\frac{150\,000}{50}\right) + 30^\circ\right].
\]
\[
12\log_{10}(3\,000) = 12 \times 3.477 = 41.7^\circ;\qquad
\tau = 50\tan(71.7^\circ) = 50 \times 3.03 \approx 152\ \mathrm{kPa}.
\]
\[
FoS = \frac{\tau A}{W\sin\psi} = \frac{152 \times 35.36}{8\,750 \times 0.7071} = \frac{5\,375}{6\,187} \approx 0.87.
\]
$FoS < 1$: with a full water table the slope is \textbf{unstable}; drainage is the first priority. (Without water, $FoS \approx \tau' A/(W\sin\psi)$ with $\sigma_n = 175$ kPa giving $\tau \approx 380$ kPa and $FoS \approx 2.2$: the slope is safe when dry --- water is the controlling factor.)

\textbf{3. Stabilisation measures.}
\begin{enumerate}
\item \textbf{Drainage} (horizontal drain holes, surface diversion channels, weep holes): lowers water pressure $U$ on J1, directly increasing effective normal stress and FoS --- it attacks the controlling factor identified above.
\item \textbf{Rock bolts / dowels and cable anchors} installed perpendicular to J1: add an active normal and shear force across the sliding plane, increasing frictional resistance; suitable in strong basalt which provides good anchorage.
\item \textbf{Regrading / benching of the slope face} (reduce face angle below $45^\circ$ or cut benches): reduces the driving force $W\sin\psi$ and removes daylighting of J1; combined with shotcrete and mesh to control raveling of the columnar jointing and to prevent toppling of J2 blocks.
\end{enumerate}

\textbf{4. Effect of paleosols.}\par
Weathered paleosols between lava flows are horizons of much lower strength (low $JCS$, low friction, high clay content) and higher permeability contrast: they act as \emph{preferential sliding surfaces} (weaker than J1, so the critical surface may relocate to a paleosol), as \emph{aquicludes} that perch groundwater and build up pore pressure at the base of the overlying flow, and as zones of deep weathering that reduce the effective cohesion of the rock mass. The analysis must therefore be repeated with paleosol shear parameters (residual $\phi$ possibly $15$--$20^\circ$, $c \approx 0$), a stepped failure path, and transient pore-pressure conditions after rain; the FoS would be significantly lower than for fresh jointed basalt alone.
\end{corrige}

\begin{corrige}[Exercise 8 — Prospection Method Selection]
\begin{center}
\begin{tabularx}{\linewidth}{|l|l|X|}\hline
\rowcolor{popblueL} Commodity & Primary method & Justification (property exploited) \\ \hline
1. Crude oil (Tertiary delta, Rio del Rey) & \textbf{Seismic reflection} & Images subsurface geometry (anticlines, growth faults, pinch-outs) from acoustic-impedance contrasts; the only method that directly delineates traps at kilometre depths. \\ \hline
2. Gold in shear zones (East Cameroon granitoids) & \textbf{Geochemical soil sampling} & Gold is detected by its chemical signature (Au and pathfinder elements As, Sb) in residual soils over mineralised structures; cheap and effective for follow-up trenching. \\ \hline
3. Limestone for cement (North Cameroon) & \textbf{Geological mapping} & Limestone is a surface/near-surface bedded unit; mapping of outcrops with sampling for CaO/MgO/SiO$_2$ directly defines extent, thickness and quality --- geophysics adds little. \\ \hline
4. Magnetite in BIF (Congo Craton) & \textbf{Magnetic survey} & Magnetite has a very high magnetic susceptibility, producing strong, sharp magnetic anomalies directly over the BIF horizons; airborne magnetics maps them even under cover. \\ \hline
5. Rock salt (deep evaporite basin) & \textbf{Gravity survey} & Halite ($\approx 2.2\ \mathrm{g/cm^3}$) is much less dense than surrounding sediments, so salt pillows and domes produce clear negative Bouguer anomalies (confirmed later by seismic). \\ \hline
6. Groundwater in fractured basement (Adamawa) & \textbf{IP/Resistivity (electrical methods)} & Water-filled fractures and weathered zones are far more conductive than fresh basement; vertical electrical soundings and resistivity profiling locate the low-resistivity fractured aquifers for borehole siting. \\ \hline
\end{tabularx}
\end{center}
\emph{Note:} radiometric and EM surveys are useful complements (EM for conductive sulfides, radiometrics for U/Th/K alteration mapping) but are not the \emph{primary} tool for any of the six cases above.
\end{corrige}

\begin{corrige}[Exercise 9 — Structured Question: Petroleum Traps and Drilling Strategy (25 marks)]
\textbf{Indicative mark scheme and model answers.}

\textbf{(a) Five traps (5 marks --- 1 mark per correctly identified trap + seal).}
\begin{enumerate}
\item \emph{Anticlinal trap} in the Cambrian sandstone, cored and sealed by the shale unit within the fold (structural; shale top seal).
\item \emph{Fault trap} on the upthrown side of the normal fault, where Cambrian reservoir is juxtaposed against shale across the fault (structural; fault seal).
\item \emph{Fault-dependent closure} on the downthrown northern limb (rollover/fault-bend fold) sealed by the Cretaceous shale (structural).
\item \emph{Unconformity/subcrop trap}: Cambrian sandstones truncated at the basal unconformity and sealed by overlying Cretaceous shale (stratigraphic).
\item \emph{Pinch-out/onlap trap}: Cretaceous reservoir sandstone onlapping and wedging out against the buried Cambrian high, sealed by the Tertiary shale (stratigraphic/combination).
\end{enumerate}
\emph{Examiner expectation:} clear labelling on the section; trap type named; seal mechanism stated (top seal, fault seal, unconformity seal). A trap without a stated seal earns at most half the mark.

\textbf{(b) Risk comparison (6 marks).}
\begin{center}
\begin{tabularx}{\linewidth}{|l|X|X|}\hline
\rowcolor{popblueL} Factor & Cambrian sandstone & Cretaceous sandstone \\ \hline
Source & Charged from overlying Cretaceous shale: requires downward/lateral migration --- moderate risk & Directly overlies/interfingers with its own source shale --- low risk \\
Reservoir & Proven quality in nearby field --- low risk & New target; porosity/permeability unproven --- moderate risk \\
Seal & Shale within fold proven, but fault may breach --- moderate & Tertiary shale thick and regionally extensive --- low risk \\
Trap & Anticline + fault closure well imaged --- low & Relies on fault seal/onlap --- moderate \\
Timing & Trap (folding) is pre-Cretaceous, predates generation --- favourable & Trap formed during/after deposition; must predate migration --- to verify \\ \hline
\end{tabularx}
\end{center}
Conclusion: the \textbf{Cambrian play has the higher chance of success} because reservoir, seal and trap are already proven by the nearby field; the Cretaceous play carries higher reservoir and fault-seal risk despite an excellent source--seal pair. (Full marks require a justified conclusion, not just a list.)

\textbf{(c) Drilling depth (4 marks).}
\[
\text{Depth below KB} = \frac{TWT}{2}\times v = \frac{2.4}{2}\ \mathrm{s}\times 3.2\ \mathrm{km/s} = 1.2 \times 3.2 = 3.84\ \mathrm{km} = 3\,840\ \mathrm{m}.
\]
TVD below sea level (datum): $3\,840 - 25 = 3\,815\ \mathrm{m\ TVDSS}$.
\emph{Examiner expectation:} halving of the two-way time (1 mark), correct use of velocity (1), subtraction of KB elevation for the sea-level datum (1), correct units (1).

\textbf{(d) Kick and kill mud weight (5 marks).}\par
\emph{Note: a hydrostatic gradient of $0.1\ \mathrm{MPa/m}$ is a misprint for the standard $\approx 0.01\ \mathrm{MPa/m}$ ($9.81\ \mathrm{kPa/m}$); the calculation below uses $0.01\ \mathrm{MPa/m}$, and candidates using the printed value consistently should not be penalised.}
\[
P_{mud} = 1.35 \times 0.01 \times 3\,200 = 43.2\ \mathrm{MPa}.
\]
\[
P_{formation} = P_{mud} + SIDPP = 43.2 + 8 = 51.2\ \mathrm{MPa}.
\]
\[
\text{Kill mud weight} = \frac{P_{formation}}{0.01 \times 3\,200} = \frac{51.2}{32} = 1.60\ \mathrm{SG}.
\]
(SICP $>$ SIDPP confirms an influx of low-density fluid, i.e. gas, in the annulus --- 1 mark for interpretation.)

\textbf{(e) Environmental and social considerations, offshore Rio del Rey (5 marks --- any five well-developed points).}
\begin{itemize}
\item \emph{Waste management:} treatment or reinjection of drilling muds and cuttings (especially oil-based muds), produced water treated before discharge; no discharge of plastics or sanitary waste.
\item \emph{Oil spill contingency:} blow-out preventer testing, oil-spill response plan with booms, skimmers and dispersants stockpiled; the basin borders mangroves of the Douala--Ed\'ea reserve and the Nigerian border, so a spill could have transboundary effects.
\item \emph{Fisheries and communities:} consultation and compensation of fishing communities (Limbe--Idenau--Kribi axis), exclusion zones negotiated, employment and local-content provisions.
\item \emph{Sensitive ecosystems:} protection of mangroves, sea turtles and manatee habitats; seismic surveys timed to avoid disturbing marine mammals (soft-start procedures).
\item \emph{Decommissioning:} platform removal and site restoration plan with financial provision.
\end{itemize}
\end{corrige}

\begin{corrige}[Exercise 10 — Essay: Geological Factors in Dam Construction (25 marks)]
\textbf{Model essay plan with indicative marking (examiner expectations in italics).}

\textbf{Introduction (within the 25 marks).} A dam is a structure whose safety depends entirely on the geology of its foundation and reservoir basin; geological investigation must precede, accompany and follow construction.

\textbf{1. Site geology and foundation conditions (6 marks).}
\begin{itemize}
\item Rock type and strength: massive, unweathered igneous/metamorphic rocks (granite, gneiss) give high bearing capacity --- e.g. the \cle{Lom Pangar Dam} on the Lom river is founded on the Pan-African granito-gneissic basement of East Cameroon, suitable for a high embankment dam after stripping of the weathered mantle.
\item Structure: faults and shear zones are weaknesses --- at \cle{Memve'ele} (Ntem river, South Cameroon) the metamorphic basement is cut by shear zones requiring dental concrete treatment and consolidation grouting; a dam axis should cross faults at high angle, never parallel.
\item Joints and weathering: joint orientation controls sliding of abutment blocks; deep tropical weathering (lateritisation) must be removed to reach sound rock, affecting excavation volumes and cost.
\item Settlement: compressible or heterogeneous foundations cause differential settlement and cracking of rigid (concrete) dams; flexible embankment dams tolerate it better.
\end{itemize}
\emph{Expectation: 3--4 developed factors with the two Cameroonian examples = full marks.}

\textbf{2. Reservoir integrity (6 marks).}
\begin{itemize}
\item Permeability of floor and flanks: a reservoir is useless if water escapes through permeable rocks; watertightness must be verified by permeability (Lugeon) tests.
\item Karst: in the limestone terrains of North Cameroon (e.g. the B\'enou\'e basin margins), solution cavities and sinkholes can drain a reservoir; sites on karst require grout curtains or must be avoided.
\item Fault leakage: open faults crossing the reservoir rim can leak water to adjacent valleys.
\item Sedimentation: high sediment yield from deforested catchments (e.g. the Sudano-Sahelian north) silts up reservoirs, shortening their life; catchment management and dead storage are geological responses.
\end{itemize}

\textbf{3. Seismicity and seismic hazard (4 marks).}
Cameroon is crossed by the \cle{Cameroon Volcanic Line} and bordered by the \cle{Benue Trough}; moderate seismicity occurs along these zones. Dams must be designed for a design earthquake (peak ground acceleration from a seismic hazard map), and the risk of \emph{reservoir-induced seismicity} (water load and pore-pressure diffusion reactivating faults) must be assessed for large reservoirs.

\textbf{4. Construction materials (4 marks).}
A dam needs millions of tonnes of materials close to site: aggregates and riprap from sound rock quarries (basalt, granite), impervious clay for the core of embankment dams, and \emph{pozzolana} (volcanic ash/scoria from the CVL, e.g. around Mount Cameroon) as a cement additive. Geological survey must locate, quantify and test these borrow areas before design is finalised --- transport distance often decides the dam type.

\textbf{5. Environmental and social geology (5 marks).}
\begin{itemize}
\item Slope stability: reservoir filling and drawdown destabilise banks (landslides into the reservoir can generate displacement waves); bank slopes must be mapped and monitored.
\item Water quality: flooding of vegetation and soils can cause eutrophication and mercury methylation; the \cle{Lake Nyos} disaster (1986, limnic CO$_2$ eruption) illustrates the danger of gas-charged deep waters and the need for degassing/monitoring of deep reservoir waters.
\item Resettlement: displaced populations must be relocated on geologically stable, non-floodable ground with adequate water supply.
\end{itemize}

\textbf{Conclusion.} Dam geology integrates rock mechanics, hydrogeology, seismic hazard and environmental geology; the success of Lom Pangar and Memve'ele rests on decades of geological investigation.\par
\emph{General marking: award full marks only for a structured essay with labelled examples, correct terminology and at least one sketch (cross-section of a dam foundation showing grout curtain, faults, weathered zone).}
\end{corrige}

\begin{corrige}[Exercise 11 — Problem Solving: Mining Method Selection and Environmental Management (25 marks)]
\textbf{(a) Method comparison (6 marks).}
\begin{center}
\begin{tabularx}{\linewidth}{|l|X|X|}\hline
\rowcolor{popblueL} Criterion & Open-pit (strip mining) & Underground (room-and-pillar) \\ \hline
Capital/operating cost & Low capital, very low unit cost (large equipment, no support) & High capital (shafts/declines, ventilation, support), high unit cost \\
Recovery & 90--100\% of the blanket & 40--60\% (ore left in pillars) \\
Dilution & Very low (selective stripping of 2--4 m overburden) & Low but contamination by floor/roof \\
Production rate & Very high, easily scaled & Limited by underground logistics \\
Safety & Excellent (no underground work); risks limited to benches and haulage & Roof falls, gas, heat; poor in weak lateritic ground \\
Suitability for thin, flat, shallow blanket & Ideal: 5--8 m ore under 2--4 m overburden gives stripping ratio $<1$ & Unsuitable: thin seam, weak roof, water table at 10 m would flood workings \\ \hline
\end{tabularx}
\end{center}
\textbf{Recommendation: open-pit strip mining} with progressive backfilling. The shallow depth, flat attitude, low stripping ratio and the water table at 10 m (above the base of any underground workings) make underground mining technically and economically absurd. \emph{(1 mark per row reasonably developed, 1 mark for the justified conclusion.)}

\textbf{(b) Environmental impacts and mitigation (7 marks).}
\begin{itemize}
\item \emph{Land clearing and overburden removal:} loss of savannah vegetation and gallery forest, destruction of fauna habitat, dust. \textbf{Mitigation:} clear only the active strip; protect gallery forests as buffer zones (legal riparian reserves); salvage and stockpile topsoil separately.
\item \emph{Mining and haulage:} dust (red bauxitic dust), noise, diesel emissions, road degradation. \textbf{Mitigation:} water spraying of haul roads, speed limits, maintenance of equipment, routing roads away from villages.
\item \emph{Waste dumps and tailings:} bauxite produces little tailings (direct shipping ore), but overburden dumps can erode and silt streams; alkaline runoff. \textbf{Mitigation:} dumps designed at angle of repose with toe drains and sediment traps; progressive backfilling of mined-out strips with overburden.
\item \emph{Mine closure:} abandoned pits, unstable faces, derelict land. \textbf{Mitigation:} closure plan with final landform regrading, revegetation and post-closure monitoring financed by a rehabilitation bond.
\end{itemize}

\textbf{(c) Rehabilitation plan (6 marks).}
\begin{enumerate}
\item \emph{Landform design:} backfill strips with overburden, regrade to gentle slopes ($< 10$--$15^\circ$) compatible with agriculture; re-establish natural drainage lines to avoid ponding and gully erosion.
\item \emph{Topsoil management:} strip topsoil ($20$--$30$ cm) separately, store in low heaps ($< 2$ m) for the shortest possible time to preserve seed bank and microbes; respread directly on regraded areas (direct return is best practice).
\item \emph{Revegetation:} rapid cover with native grasses and legumes (e.g. \emph{Stylosanthes}, local \emph{Brachiaria}) to fix nitrogen and control erosion; then plant native and agroforestry trees (e.g. \emph{Acacia}, fruit trees) to return land to agriculture/grazing use by local communities.
\item \emph{Water management:} retain sediment ponds until runoff is clean, then breach and restore stream courses; re-establish gallery forest along restored streams.
\item \emph{Monitoring and handover:} 5--10 years of monitoring (vegetation cover, soil quality, water quality), with formal handover criteria agreed with the communities and the administration.
\end{enumerate}

\textbf{(d) EIA under Law No. 96/12 (6 marks).}\par
Key components of the EIA report:
\begin{enumerate}
\item Project description (deposit, method, schedule, infrastructure).
\item \textbf{Baseline study} --- the geologist's core contribution: geological mapping of the deposit and surroundings; \emph{hydrogeology} (water table at 10 m, borehole inventory, risk of drawdown or contamination of village wells); \emph{geochemistry} (background levels of Al, Fe, trace metals in soils and streams to distinguish natural from mining pollution later); \emph{geomorphology and soils} (slope, erosion risk, soil types for rehabilitation).
\item Impact identification and evaluation (on air, water, soils, biodiversity, communities) for each project phase.
\item Mitigation and Environmental Management Plan (EMP), including the rehabilitation plan of part (c).
\item Public consultation of affected communities (agriculturalists and grazers of the Adamawa Plateau).
\item Monitoring programme and closure/decommissioning plan with financial provisions.
\end{enumerate}
\emph{Examiner expectation: the candidate must explicitly link the geologist's role to baseline data (geology, hydrogeology, geochemistry, geomorphology) and cite the legal framework; a generic EIA list without the geological content earns at most half marks.}
\end{corrige}

\newpage

\coursec{Summary sheet}

\begin{popfigure}
\begin{tikzpicture}[mindmap, grow cyclic, every node/.style={concept, text=white, font=\small},
root concept/.append style={concept color=popdark, font=\bfseries\small, minimum size=2.6cm},
level 1 concept/.append style={level distance=3.6cm, sibling angle=51.4, font=\scriptsize},
level 2 concept/.append style={level distance=2.4cm, font=\tiny}]
\node[root concept, align=center]{Economic\\Geology}
  child[concept color=poporange]{ node{Coal}
    child{ node{peat $\to$ lignite $\to$ bituminous $\to$ anthracite} }
    child{ node{swamps, burial, coalification} }
  }
  child[concept color=popblue]{ node{Petroleum \& gas}
    child{ node{source, reservoir, seal, trap} }
    child{ node{anticline, fault, stratigraphic traps} }
  }
  child[concept color=popgreen]{ node{Uses of rocks}
    child{ node{building, cement, aggregates} }
    child{ node{ores \& industrial minerals} }
  }
  child[concept color=poppurple]{ node{Engineering geology}
    child{ node{dams, reservoirs, tunnels} }
    child{ node{slope stability, factor of safety} }
  }
  child[concept color=poppink]{ node{Prospection}
    child{ node{geological, geochemical} }
    child{ node{geophysics: seismic, gravity, magnetic, electrical} }
  }
  child[concept color=popgold]{ node{Mining \& environment}
    child{ node{open pit vs underground} }
    child{ node{impacts \& rehabilitation} }
  };
\end{tikzpicture}
\legende{Mind map of the chapter Economic Geology: fossil fuels, uses of rocks, engineering geology, prospection and mining.}
\end{popfigure}

\begin{aretenir}[Key Definitions — Economic Geology]
\begin{itemize}
\item \cle{Coal}: a combustible sedimentary rock formed by the accumulation and transformation of plant debris in swamps, through \emph{peatification} (biochemical stage, near the surface) then \emph{coalification} (geochemical stage under burial, rising temperature and pressure).
\item \cle{Rank of coal}: the degree of coalification reached: peat $\to$ lignite $\to$ bituminous coal $\to$ anthracite. As rank increases, carbon content and calorific value increase, while moisture and volatile matter decrease.
\item \cle{Petroleum system}: the complete set of elements and processes needed for an oil or gas accumulation: a \emph{source rock} (organic-rich, containing kerogen) $\to$ \emph{maturation} in the oil window (about $60$--$120\ ^\circ$C) $\to$ primary migration (out of the source rock) then secondary migration (through permeable beds) $\to$ accumulation in a \emph{reservoir rock} (porous and permeable) capped by a \emph{seal} (impermeable cap rock) within a \emph{trap}.
\item \cle{Trap}: a geometric arrangement that prevents hydrocarbons from escaping: \emph{structural traps} (anticline, fault) or \emph{stratigraphic traps} (pinch-out, unconformity, fossil reef).
\item \cle{Porosity}: the percentage of void space in a rock; \cle{permeability}: the ability of a rock to transmit fluids through interconnected pores. Both are essential reservoir properties.
\item \cle{Ore}: a rock from which a metal or useful mineral can be extracted \emph{profitably}; the worthless associated material is the \emph{gangue}.
\item \cle{Factor of safety} of a slope: $F = \dfrac{\text{resisting forces}}{\text{driving forces}}$; the slope is stable if $F>1$ and unstable if $F<1$.
\end{itemize}
\end{aretenir}

\begin{aretenir}[Facts and Formulas to Know]
\begin{itemize}
\item Coalification series with increasing carbon content: peat (about $60\ \%$ C or less) $\to$ lignite $\to$ bituminous coal $\to$ anthracite (more than $90\ \%$ C).
\item The \emph{gas window} lies at higher temperatures than the oil window; excessive temperature destroys liquid oil, leaving only dry gas (methane) or, at extreme grades, graphite.
\item Porosity: $\phi = \dfrac{V_{\text{voids}}}{V_{\text{total}}}\times 100\ \%$.
\item Slope failure is favoured by water (raises pore-water pressure, adds weight, lubricates slip planes), slope steepening, weak layers (clay, shale) and discontinuities (joints, bedding) dipping out of the slope.
\item Dam sites require: a narrow valley, an impermeable and strong foundation rock, low fracture density and an adequate spillway; reservoirs need an impermeable basin (avoid karstic limestone and major open faults).
\item Tunnels: the alignment is chosen with respect to strike and dip for maximum stability; avoid fault zones, water-bearing horizons and squeezing ground (soft, swelling rocks).
\end{itemize}
\end{aretenir}

\begin{methode}[Prospection, Exploration and Mining]
\begin{enumerate}
\item \cle{Prospection}: geological mapping and stratigraphic study of the area; geochemical sampling (stream sediments, soils, rock chips) to detect anomalies.
\item \cle{Geophysical methods}: seismic reflection and refraction (subsurface structures, traps), gravimetry (density contrasts, salt domes), magnetometry (magnetic minerals, basement depth), electrical resistivity (water table, conductive ore bodies), and well logging in boreholes.
\item \cle{Mining}: open-pit methods (quarrying, strip mining) for shallow deposits; underground mining (shafts, galleries) for deep deposits. The choice depends on depth, geometry, grade of the deposit and economic factors.
\item \cle{Environmental management}: environmental impact assessment before mining, control of waste rock and tailings, treatment of mine water, then site rehabilitation and revegetation after closure.
\end{enumerate}
\end{methode}

\begin{attention}[Common Mistakes]
\begin{itemize}
\item Confusing \emph{porosity} with \emph{permeability}: a clay can be highly porous yet practically impermeable because its pores are not interconnected.
\item Saying oil forms \emph{in} the reservoir rock: it forms in the \emph{source rock} and migrates into the reservoir.
\item Reversing the coal rank order, or claiming that anthracite is the youngest (least transformed) stage.
\item Forgetting that a petroleum accumulation needs \emph{all} the elements: source rock, reservoir, seal, trap geometry and correct timing of trap formation relative to migration.
\item Believing faults are always good traps: an open (unsealed) fault can act as a migration pathway and drain the accumulation.
\item Ignoring water as the main trigger of landslides, or forgetting to mention it in slope-stability questions.
\end{itemize}
\end{attention}

\begin{methode}[GCE Examination Tips]
\begin{enumerate}
\item Draw and label a neat anticline-trap cross-section (gas--oil--water order from top to bottom, seal above, reservoir in the middle, source rock below): a classic, highly rewarded question.
\item Quote Cameroonian examples naturally: petroleum in the Rio del Rey and Douala/Kribi-Campo basins, limestone for cement at Figuil, the Lake Nyos gas hazard, volcanic materials of Mount Cameroon.
\item In essays on mining and the environment, always balance economic benefits (jobs, revenue, raw materials) with impacts (deforestation, water pollution, subsidence) and end with rehabilitation measures.
\item Define every technical term before using it; examiners reward precise vocabulary and clearly labelled diagrams.
\end{enumerate}
\end{methode}

\findecours

\end{document}
